\documentclass[
  justified,
  openright,
]{tufte-book}

% --------------------------------------------------
% Engine sanity (XeLaTeX or LuaLaTeX required)
% --------------------------------------------------
\usepackage{iftex}
\ifPDFTeX
  \PackageError{template}{This template requires XeLaTeX or LuaLaTeX}{}
\fi

% --------------------------------------------------
% Core math
% --------------------------------------------------
\usepackage{amsmath,amssymb,amsfonts}

% --------------------------------------------------
% Tufte: subsubsection is forbidden --- remap safely
% --------------------------------------------------
\makeatletter
\let\subsubsection\paragraph
\let\paragraph\paragraph
\makeatother

% --------------------------------------------------
% Fix Pandoc tabularray counter=none bug
% --------------------------------------------------
\makeatletter
\ExplSyntaxOn
\cs_if_exist:NF \c@none { \newcounter{none} }
\ExplSyntaxOff
\makeatother

\usepackage{fontspec}

\IfFontExistsTF{TeX Gyre Pagella}
  {\setmainfont{TeX Gyre Pagella}}
  {\setmainfont{Latin Modern Roman}}

\IfFontExistsTF{TeX Gyre Heros}
  {\setsansfont[Scale=0.92]{TeX Gyre Heros}}
  {\setsansfont{Latin Modern Sans}}

\IfFontExistsTF{TeX Gyre Cursor}
  {\setmonofont[Scale=0.82]{TeX Gyre Cursor}}
  {\setmonofont{Latin Modern Mono}}
% --------------------------------------------------
% Graphics and tables
% --------------------------------------------------
\usepackage{graphicx}
\usepackage{booktabs}
\usepackage{longtable}
\usepackage{array}
\usepackage{tabularray}
\UseTblrLibrary{booktabs}

\setkeys{Gin}{width=\maxwidth,height=\maxheight,keepaspectratio}

% --------------------------------------------------
% Colors
% --------------------------------------------------
\usepackage[table,xcdraw,svgnames]{xcolor}

\definecolor{DarkGray}{RGB}{55,60,70}
\definecolor{CalloutBlue}{HTML}{DCEEFF}
\definecolor{CalloutBorder}{HTML}{5B8BD9}
\definecolor{BrandBlue}{RGB}{18,65,124}
\definecolor{BrandAccent}{RGB}{0,153,204}
\definecolor{LightGray}{RGB}{245,246,248}
\definecolor{MidGray}{RGB}{175,180,190}
\definecolor{WarningAmber}{RGB}{230,100,0}
\definecolor{WarningBg}{RGB}{255,248,235}
\definecolor{CautionYellow}{RGB}{180,140,0}
\definecolor{CautionBg}{RGB}{255,252,230}
\definecolor{TipGreen}{RGB}{36,110,44}
\definecolor{TipBg}{RGB}{237,248,239}
\definecolor{NoteBg}{RGB}{235,243,255}
\definecolor{NoteBlue}{RGB}{18,90,160}
\definecolor{DangerRed}{RGB}{185,28,28}
\definecolor{DangerBg}{RGB}{255,237,237}
\definecolor{TableHeader}{RGB}{18,65,124}
\definecolor{TableRowAlt}{RGB}{235,241,252}
\definecolor{MarginRule}{RGB}{210,215,225}
\definecolor{SpineAccent}{RGB}{0,153,204}
% --------------------------------------------------
% Pandoc syntax highlighting macros
% --------------------------------------------------
\newcommand{\AlertTok}[1]{\textcolor{Red}{#1}}
\newcommand{\AnnotationTok}[1]{\textcolor{Gray}{#1}}
\newcommand{\AttributeTok}[1]{\textcolor{DarkSlateGray}{#1}}
\newcommand{\BaseNTok}[1]{\textcolor{Purple}{#1}}
\newcommand{\BuiltInTok}[1]{#1}
\newcommand{\CharTok}[1]{\textcolor{Teal}{#1}}
\newcommand{\CommentTok}[1]{\textcolor{MidGray}{#1}}
\newcommand{\CommentVarTok}[1]{\textcolor{MidGray}{#1}}
\newcommand{\ConstantTok}[1]{\textcolor{Purple}{#1}}
\newcommand{\ControlFlowTok}[1]{\textcolor{BrandBlue}{#1}}
\newcommand{\DataTypeTok}[1]{\textcolor{Purple}{#1}}
\newcommand{\DecValTok}[1]{\textcolor{Purple}{#1}}
\newcommand{\DocumentationTok}[1]{\textcolor{MidGray}{#1}}
\newcommand{\ErrorTok}[1]{\textcolor{Red}{#1}}
\newcommand{\ExtensionTok}[1]{#1}
\newcommand{\FloatTok}[1]{\textcolor{Purple}{#1}}
\newcommand{\FunctionTok}[1]{\textcolor{BrandBlue}{#1}}
\newcommand{\ImportTok}[1]{\textcolor{BrandAccent}{#1}}
\newcommand{\InformationTok}[1]{\textcolor{BrandAccent}{#1}}
\newcommand{\KeywordTok}[1]{\textcolor{BrandBlue}{#1}}
\newcommand{\NormalTok}[1]{#1}
\newcommand{\OperatorTok}[1]{\textcolor{DarkGray}{#1}}
\newcommand{\OtherTok}[1]{#1}
\newcommand{\PreprocessorTok}[1]{\textcolor{BrandAccent}{#1}}
\newcommand{\RegionMarkerTok}[1]{#1}
\newcommand{\SpecialCharTok}[1]{\textcolor{Teal}{#1}}
\newcommand{\SpecialStringTok}[1]{\textcolor{Teal}{#1}}
\newcommand{\StringTok}[1]{\textcolor{Teal}{#1}}
\newcommand{\VariableTok}[1]{#1}
\newcommand{\VerbatimStringTok}[1]{\textcolor{Teal}{#1}}
\newcommand{\WarningTok}[1]{\textcolor{WarningAmber}{#1}}
% --------------------------------------------------
% Ensure ALL Pandoc *Tok macros are defined (robust)
% --------------------------------------------------
\makeatletter
\@for\tok:=Alert,Annotation,Attribute,BaseN,BuiltIn,Char,Comment,CommentVar,%
Constant,ControlFlow,DataType,DecVal,Documentation,Error,Extension,Float,%
Function,Import,Information,Keyword,Normal,Operator,Other,Preprocessor,%
RegionMarker,SpecialChar,SpecialString,String,Variable,VerbatimString,Warning%
\do{%
  \expandafter\providecommand\csname\tok Tok\endcsname[1]{#1}%
}
\makeatother

% --------------------------------------------------
% Pandoc code block support (Shaded environment)
% --------------------------------------------------
\usepackage{fancyvrb}
\usepackage{framed}

\definecolor{shadecolor}{RGB}{248,248,248}

\newenvironment{Shaded}{%
  \begin{snugshade}
}{%
  \end{snugshade}
}

% --------------------------------------------------
% Pandoc Highlighting environment
% --------------------------------------------------
\DefineVerbatimEnvironment{Highlighting}{Verbatim}{
  commandchars=\\\{\},
  fontsize=\small
}

% --------------------------------------------------
% Tufte: subsubsection is forbidden --- remap safely
% --------------------------------------------------
\makeatletter
\let\subsubsection\paragraph
\let\paragraph\paragraph
\makeatother

% --------------------------------------------------
% Utilities
% --------------------------------------------------
\usepackage{etoolbox}
\usepackage{xparse}
\usepackage{ifthen}
\usepackage{textcase}
\usepackage{titlesec}
% --------------------------------------------------
% Fix Pandoc tabularray counter=none bug
% --------------------------------------------------
\makeatletter
\ExplSyntaxOn
\cs_if_exist:NF \c@none { \newcounter{none} }
\ExplSyntaxOff
\makeatother
% --------------------------------------------------
% Paragraph & section spacing
% --------------------------------------------------
\setlength{\parindent}{1em}
\setlength{\parskip}{0pt}

\titlespacing*{\section}{0pt}{2.0ex plus .2ex}{1.0ex}
\titlespacing*{\subsection}{0pt}{1.5ex plus .2ex}{0.8ex}

% --------------------------------------------------
% Margin notes (tufte)
% --------------------------------------------------
\setlength{\marginparwidth}{2.2in}

\newcommand{\mn}[1]{%
  \marginnote{\raggedright\footnotesize\color{DarkGray}#1}
}

\NewDocumentCommand{\margincallout}{m}{%
  \marginnote{%
    \setlength{\fboxsep}{6pt}%
    \colorbox{CalloutBlue}{%
      \begin{minipage}{\marginparwidth}
      \footnotesize #1
      \end{minipage}
    }
  }
}

% --------------------------------------------------
% Soul / tabularray / marking fixes
% --------------------------------------------------
\usepackage{soul}
\AfterEndEnvironment{longtblr}{\soulsuspend}

\makeatletter
\providecommand{\soulsuspend}{}
\providecommand{\soulsuspend*}{}
\providecommand{\soulsuspend@}{}
\makeatother

\DeclareRobustCommand{\checkbox}{\(\square\)}

\soulregister\checkbox{0}
\soulregister\textbf{1}
\soulregister\emph{1}
\soulregister\href{2}
\soulregister\label{0}
\soulregister\ref{0}
\soulregister\pageref{0}
\soulregister\footnote{1}
\soulregister\cite{0}

% --------------------------------------------------
% Headers / marks (no case mangling)
% --------------------------------------------------
\makeatletter
\renewcommand{\chaptermark}[1]{\markboth{#1}{}}
\renewcommand{\sectionmark}[1]{\markright{#1}}

\pretocmd\markboth{\let\allcapsspacing\relax\let\smallcapsspacing\relax}{}{}
\pretocmd\markright{\let\allcapsspacing\relax\let\smallcapsspacing\relax}{}{}
\makeatother

% --------------------------------------------------
% Blank page helper
% --------------------------------------------------
\newcommand{\blankpage}{%
  \clearpage
  \thispagestyle{empty}
  \null\vfill
  \begin{center}
    \color{MidGray}\sffamily\small\itshape
    (This page intentionally left blank.)
  \end{center}
  \vfill
  \clearpage
}

% --------------------------------------------------
% Fancyhdr
% --------------------------------------------------
\usepackage{fancyhdr}
\pagestyle{fancy}
\fancyhf{}
\fancyhead[LE]{\small\scshape\leftmark}
\fancyhead[RO]{\small\scshape\rightmark}
\fancyfoot[C]{\thepage}
\AtBeginDocument{\pagestyle{fancy}}

% --------------------------------------------------
% Hyperlinks
% --------------------------------------------------
\usepackage{hyperref}
\hypersetup{
  colorlinks=true,
  linkcolor=blue,
  urlcolor=blue,
  citecolor=blue,
  hypertexnames=false
}

% --------------------------------------------------
% Pandoc helpers
% --------------------------------------------------
\providecommand{\tightlist}{%
  \setlength{\itemsep}{0pt}\setlength{\parskip}{0pt}}

% Show only parts and chapters in TOC; number only parts
\setcounter{tocdepth}{1}
\setcounter{secnumdepth}{-1}

% --------------------------------------------------
% Metadata
% --------------------------------------------------



% --------------------------------------------------
% ============================================================
%  DOCUMENT
% ============================================================
\begin{document}
% -- Front matter -------------------------------------------
\frontmatter
\pagestyle{fancy}

% 1) Blank half title page
\thispagestyle{empty}
\null
\clearpage

% 2) Frontispiece (placeholder)
\thispagestyle{empty}
\begin{center}
  \vspace*{0.35\textheight}
  {\sffamily\large Frontispiece}
\end{center}
\clearpage

% 3) Full title page
\begin{titlepage}
    \vspace*{1.6in}
    {\sffamily\fontsize{9}{11}\selectfont\MakeUppercase{Training Manual}}
    \vspace{0.3in}
    {\color{BrandAccent}\rule{4.5in}{2pt}}
    \vspace{0.25in}
    % ADD YOUR TITLE TEXT HERE:
    {\sffamily\bfseries\fontsize{32}{38}\selectfont Coding With //ScreenReaders}
    \vspace{0.2in}
    {\color{BrandAccent}\rule{4.5in}{0.6pt}}
    \vspace{0.5in}
    \vfill
    {\sffamily\fontsize{7.5}{9}\selectfont\color{MidGray}%
      CONFIDENTIAL --- For Internal Use Only%
    }
  \clearpage
  \nopagecolor
\end{titlepage}

% 4) Copyright page (placeholder)
\thispagestyle{empty}
\vspace*{0.25\textheight}
{\small
\noindent Copyright \textcopyright\ \the\year\ Michael Ryan Hunsaker, M.Ed., Ph.D. Coding With ScreenReaders is licensed under a Creative Commons Attribution-NoNCommerical-NoDerivatives 4.0 International: https://creativecommons.org/licenses/by-nc-nd/4.0/
\noindent All rights reserved.
\noindent NO AI TRAINING: Without in any way limiting the author's rights under copyright, any use of this publication to “train” generative artificial intelligence (AI) technologies to generate text is expressly prohibited.
}
\clearpage

% 5) Contents
{
\hypersetup{linkcolor=BrandBlue}
\tableofcontents
\clearpage
}

% 6) Dedication (placeholder)
\chapter*{Dedication}
\thispagestyle{empty}
\begin{center}
  \vspace*{0.3\textheight}
  {\itshape To the countless individuals who have been told they lack the capability to code, yet remain determined to learn and create through accessible coding practices, this book is dedicated.}\\[1em]
\end{center}
\clearpage

% 7) Epigraph (placeholder)
\chapter*{Epigraph}
\thispagestyle{empty}
\begin{flushright}
  \vspace*{0.25\textheight}
  {\itshape The power of the Web is in its universality. Access by everyone regardless of disability is an essential aspect}\\[1em]
  --- Tim Berners-Lee, Inventor of the World Wide Web
\end{flushright}
\clearpage

% 8) Introduction (placeholder)
\newpage

% -- Main content -------------------------------------------
\mainmatter
\pagenumbering{arabic} % Switch to numbers (1, 2, 3...) for the main content


\chapter*{Introduction to Computer Science with Java}

\subsubsection{For Visually Impaired
Students}\label{for-visually-impaired-students}

\subsection{Course Overview}\label{course-overview}

This curriculum introduces computer science concepts using Java, adapted
for students who rely on screen readers. All exercises are text-based
and terminal-focused. No graphical output is required.

\textbf{Target Audience:} Students with no prior programming experience
using screen readers (NVDA, JAWS, VoiceOver, or Windows Narrator).

\begin{center}\rule{0.5\linewidth}{0.5pt}\end{center}

\subsection{Tools Required}\label{tools-required}

\begin{itemize}
\tightlist
\item
  \textbf{Java Development Kit (JDK) 17+} ---
  \url{https://adoptium.net/}
\item
  \textbf{Visual Studio Code} with the \textbf{Extension Pack for Java}
  by Microsoft
\item
  A screen reader (NVDA, JAWS, or VoiceOver)
\end{itemize}

\subsection{Accessibility Notes}\label{accessibility-notes}

\begin{itemize}
\tightlist
\item
  All program output is text-based (console/terminal).
\item
  VS Code's Screen Reader Optimized Mode works well with Java.
\item
  Configure your screen reader to announce indentation changes --- Java
  uses braces \texttt{\{\}} to define blocks, but consistent indentation
  is still essential for readability.
\item
  Use \texttt{Ctrl\ +\ Shift\ +\ P} -\textgreater{} ``Toggle Screen
  Reader Optimized'' in VS Code.
\item
  Use \texttt{F8} / \texttt{Shift+F8} to navigate between errors.
\end{itemize}

\begin{center}\rule{0.5\linewidth}{0.5pt}\end{center}

\subsection{Units}\label{units}

{\def\LTcaptype{none} % do not increment counter
\begin{longtable}[]{@{}ll@{}}
\toprule\noalign{}
Unit & Title \\
\midrule\noalign{}
\endhead
\bottomrule\noalign{}
\endlastfoot
Unit 0 & Setting Up Your Java Development Environment \\
Unit 1 & Foundations of Computing and Getting Started with Java \\
Unit 2 & Writing Simple Programs and Numeric Computations \\
Unit 3 & Control Flow --- Decisions and Loops \\
Unit 4 & Data Structures --- Arrays and ArrayLists \\
Unit 5 & User Input and String Processing \\
Unit 6 & Methods (Functions) \\
Unit 7 & Classes and Object-Oriented Programming \\
Unit 8 & Files and Exception Handling \\
Unit 9 & Collections and Basic Data Analysis \\
Unit 10 & Project-Based Learning \\
Unit 11 & Capstone Project \\
\end{longtable}
}

\begin{center}\rule{0.5\linewidth}{0.5pt}\end{center}

\subsection{References}\label{references}

\begin{itemize}
\tightlist
\item
  Eck, D.J. (2022). \emph{Introduction to Programming Using Java} (9th
  ed.). Free PDF available at \url{https://math.hws.edu/javanotes/}
\item
  Horstmann, C.S. (2018). \emph{Big Java: Early Objects} (6th ed.).
  Wiley.
\end{itemize}

\newpage

\section{Teacher's Guide: Introduction to Computer Science with
Java}\label{teachers-guide-introduction-to-computer-science-with-java}

\subsection{For Students with Visual Impairments --- Complete Instructor
Edition}\label{for-students-with-visual-impairments-complete-instructor-edition}

\begin{center}\rule{0.5\linewidth}{0.5pt}\end{center}

\subsection{Table of Contents}\label{table-of-contents}

\begin{enumerate}
\def\labelenumi{\arabic{enumi}.}
\tightlist
\item
  \hyperref[philosophy-and-why-java]{Philosophy and Why Java}
\item
  \hyperref[java-specific-accessibility-challenges]{Java-Specific
  Accessibility Challenges}
\item
  \hyperref[classroom-setup-and-configuration]{Classroom Setup and
  Configuration}
\item
  \hyperref[spoken-java-conventions]{Spoken Java Conventions}
\item
  \hyperref[unit-01-setup-and-foundations]{Unit 0--1: Setup and
  Foundations}
\item
  \hyperref[unit-2-io-and-the-math-class]{Unit 2: I/O and the Math
  Class}
\item
  \hyperref[unit-3-control-flow]{Unit 3: Control Flow}
\item
  \hyperref[unit-4-arrays-and-arraylists]{Unit 4: Arrays and ArrayLists}
\item
  \hyperref[unit-5-strings-and-string-processing]{Unit 5: Strings and
  String Processing}
\item
  \hyperref[unit-6-methods]{Unit 6: Methods}
\item
  \hyperref[unit-7-classes-and-oop]{Unit 7: Classes and OOP}
\item
  \hyperref[unit-8-files-and-exception-handling]{Unit 8: Files and
  Exception Handling}
\item
  \hyperref[unit-9-collections-analysis-and-testing]{Unit 9:
  Collections, Analysis, and Testing}
\item
  \hyperref[units-1011-projects-and-capstone]{Units 10--11: Projects and
  Capstone}
\item
  \hyperref[bridging-from-python-to-java]{Bridging from Python to Java}
\item
  \hyperref[assessment-strategies]{Assessment Strategies}
\item
  \hyperref[supporting-diverse-needs]{Supporting Diverse Needs}
\item
  \hyperref[alternative-and-supplementary-resources]{Alternative and
  Supplementary Resources}
\item
  \hyperref[quick-reference-sheets]{Quick Reference Sheets}
\end{enumerate}

\begin{center}\rule{0.5\linewidth}{0.5pt}\end{center}

\subsection{1. Philosophy and Why Java}\label{philosophy-and-why-java}

\subsubsection{The Case for Teaching Java to Blind
Students}\label{the-case-for-teaching-java-to-blind-students}

Java is not the easiest language for beginners, and its verbosity
creates overhead that requires specific pedagogical attention. Despite
this, Java has genuine advantages for visually impaired learners that
are often overlooked:

\textbf{Compile-time error detection.} Java's compiler catches type
errors, name errors, and structural mistakes before any code runs. A
blind student does not need to search through runtime output for subtle
type mismatches --- the compiler announces them precisely and
immediately with file name, line number, and error description.

\textbf{Strong, static typing is auditorily informative.} When
IntelliSense in VS Code suggests methods, it includes the return type:
``\texttt{double\ getBalance()} --- returns a double.'' This type
information helps a screen reader user understand what a method does
without having to navigate to its definition.

\textbf{Enforced code structure.} Every Java program requires a class;
every executable program requires a \texttt{main} method with a specific
signature. This structure is more overhead to type but less ambiguity
about how programs are organized. For a student who cannot visually scan
a file for structure, predictable structural conventions reduce
cognitive load.

\textbf{Exception specification.} Java's checked exceptions force
programmers to declare which exceptions a method can throw. This means a
screen reader user reading a method signature gets information about
failure modes that Python provides only through documentation.

\textbf{Enterprise relevance.} Java remains a dominant language in
enterprise software, Android development, and many backend systems. A
student who learns Java learns patterns that transfer directly to
professional employment.

\begin{center}\rule{0.5\linewidth}{0.5pt}\end{center}

\subsection{2. Java-Specific Accessibility
Challenges}\label{java-specific-accessibility-challenges}

\subsubsection{The Verbosity Problem}\label{the-verbosity-problem}

Java requires significantly more text than Python to accomplish the same
task:

\begin{Shaded}
\begin{Highlighting}[]
\CommentTok{// Java: Hello, World}
\KeywordTok{public} \KeywordTok{class}\NormalTok{ HelloWorld }\OperatorTok{\{}
    \KeywordTok{public} \DataTypeTok{static} \DataTypeTok{void} \FunctionTok{main}\OperatorTok{(}\BuiltInTok{String}\OperatorTok{[]}\NormalTok{ args}\OperatorTok{)} \OperatorTok{\{}
        \BuiltInTok{System}\OperatorTok{.}\FunctionTok{out}\OperatorTok{.}\FunctionTok{println}\OperatorTok{(}\StringTok{"Hello, World!"}\OperatorTok{);}
    \OperatorTok{\}}
\OperatorTok{\}}
\end{Highlighting}
\end{Shaded}

\begin{Shaded}
\begin{Highlighting}[]
\CommentTok{\# Python: Hello, World}
\BuiltInTok{print}\NormalTok{(}\StringTok{"Hello, World!"}\NormalTok{)}
\end{Highlighting}
\end{Shaded}

For a sighted student, extra text is a visual annoyance. For a blind
student, it is additional audio before they reach the meaningful code.
\textbf{Mitigation strategies:}

\begin{enumerate}
\def\labelenumi{\arabic{enumi}.}
\item
  \textbf{Code templates.} Provide a starter file the student copies for
  every new program. They modify only the method body for early lessons.
  This reduces the typing burden and prevents structural errors while
  the student is learning concepts.
\item
  \textbf{JShell (Java REPL).} Java includes an interactive shell since
  Java 9: run \texttt{jshell} in the terminal. JShell strips away class
  and method boilerplate. For exploration and experimentation, JShell is
  far more accessible than creating files:

\begin{verbatim}
jshell> int x = 5;
x ==> 5
jshell> System.out.println(x * 2);
10
\end{verbatim}
\item
  \textbf{Progressive disclosure.} In early lessons, treat
  \texttt{public\ static\ void\ main(String{[}{]}\ args)\ \{} as a
  formula, not something to understand. Students learn to recognize it
  by ear. Explain each word only when the concept it relates to (access
  modifiers, static methods, arrays) is being taught.
\end{enumerate}

\subsubsection{Integer Division
Surprises}\label{integer-division-surprises}

Java's \texttt{/} operator performs integer division when both operands
are integers:

\begin{itemize}
\tightlist
\item
  \texttt{15\ /\ 4} = \texttt{3} (not \texttt{3.75})
\item
  \texttt{15.0\ /\ 4} = \texttt{3.75}
\item
  \texttt{(double)\ 15\ /\ 4} = \texttt{3.75}
\end{itemize}

This is silent --- there is no error or warning. A blind student may not
catch a calculation that silently discards its decimal portion. Teach
the rule explicitly: \textbf{if you want a decimal result, at least one
operand must be a \texttt{double}.}

\subsubsection{The String Equality Trap}\label{the-string-equality-trap}

\texttt{==} on Strings compares memory addresses, not content. Two
strings with identical characters in different memory locations are
\texttt{==} false but \texttt{.equals()} true. This never appears in
Python; blind students coming from Python are not conditioned to suspect
\texttt{==} as a source of String bugs.

\textbf{Teaching approach:} Show the broken behavior explicitly. Use
\texttt{new\ String("hello")\ ==\ new\ String("hello")} which is
guaranteed to be \texttt{false} even though the strings are identical.
Then show \texttt{.equals()} which is \texttt{true}. Make
\texttt{.equals()} for Strings a non-negotiable rule from day one.

\subsubsection{Null Pointer Exceptions}\label{null-pointer-exceptions}

\texttt{NullPointerException} is the most common Java runtime error. It
occurs when you call a method or access a field on a variable that holds
\texttt{null} (no object reference). Unlike Python's
\texttt{AttributeError} which gives useful context, Java's NPE message
is often terse. Teaching students to recognize and debug NPEs is
essential.

\begin{center}\rule{0.5\linewidth}{0.5pt}\end{center}

\subsection{3. Classroom Setup and
Configuration}\label{classroom-setup-and-configuration}

\subsubsection{IDE Configuration}\label{ide-configuration}

\textbf{VS Code with Extension Pack for Java:} This is the recommended
setup. The Extension Pack includes Language Support for Java (Red Hat),
Debugger for Java, Java Test Runner, Maven for Java, and Project Manager
for Java.

Alternative IDEs to consider:

{\def\LTcaptype{none} % do not increment counter
\begin{longtable}[]{@{}
  >{\raggedright\arraybackslash}p{(\linewidth - 4\tabcolsep) * \real{0.2090}}
  >{\raggedright\arraybackslash}p{(\linewidth - 4\tabcolsep) * \real{0.0970}}
  >{\raggedright\arraybackslash}p{(\linewidth - 4\tabcolsep) * \real{0.6940}}@{}}
\toprule\noalign{}
\begin{minipage}[b]{\linewidth}\raggedright
IDE
\end{minipage} & \begin{minipage}[b]{\linewidth}\raggedright
Accessibility
\end{minipage} & \begin{minipage}[b]{\linewidth}\raggedright
Notes
\end{minipage} \\
\midrule\noalign{}
\endhead
\bottomrule\noalign{}
\endlastfoot
\textbf{VS Code + Extension Pack} & High & Recommended. Full
accessibility via keyboard. \\
\textbf{IntelliJ IDEA Community} & Medium--High & Powerful but complex.
More keyboard shortcuts to learn. JetBrains has an accessibility
guide. \\
\textbf{Eclipse} & Medium & Older UI; can be navigated but more complex
accessibility model. \\
\textbf{BlueJ} & Low--Medium & Designed for teaching; visual design
features may impede accessibility. \\
\textbf{DrJava} & Low--Medium & Simple; not actively maintained. \\
\textbf{jGRASP} & Low & Visualization-focused; primarily visual. \\
\textbf{Greenfoot} & Low & Visual/game-focused; not accessible for
screen reader use. \\
\end{longtable}
}

\textbf{VS Code Java Specific Settings:}

\begin{Shaded}
\begin{Highlighting}[]
\FunctionTok{\{}
  \DataTypeTok{"java.completion.enabled"}\FunctionTok{:} \KeywordTok{true}\FunctionTok{,}
  \DataTypeTok{"java.format.enabled"}\FunctionTok{:} \KeywordTok{true}\FunctionTok{,}
  \DataTypeTok{"editor.formatOnSave"}\FunctionTok{:} \KeywordTok{true}\FunctionTok{,}
  \DataTypeTok{"java.saveActions.organizeImports"}\FunctionTok{:} \KeywordTok{true}\FunctionTok{,}
  \DataTypeTok{"java.completion.importOrder"}\FunctionTok{:} \OtherTok{[}\StringTok{"java"}\OtherTok{,} \StringTok{"javax"}\OtherTok{,} \StringTok{"com"}\OtherTok{,} \StringTok{""}\OtherTok{]}\FunctionTok{,}
  \DataTypeTok{"editor.suggestOnTriggerCharacters"}\FunctionTok{:} \KeywordTok{true}
\FunctionTok{\}}
\end{Highlighting}
\end{Shaded}

\textbf{Key accessibility setting:} Enable ``Trigger suggestions'' on
\texttt{.} so that typing \texttt{myObject.} causes IntelliSense to
announce available methods. This allows audio exploration of an object's
API.

\subsubsection{JDK Installation
Verification}\label{jdk-installation-verification}

\begin{Shaded}
\begin{Highlighting}[]
\ExtensionTok{java} \AttributeTok{{-}{-}version}      \CommentTok{\# Runtime version}
\ExtensionTok{javac} \AttributeTok{{-}{-}version}     \CommentTok{\# Compiler version}
\ExtensionTok{jshell}              \CommentTok{\# Start REPL (Ctrl+D or /exit to quit)}
\end{Highlighting}
\end{Shaded}

All three must work before proceeding. If \texttt{jshell} is not found,
the student may have installed only the JRE (runtime), not the full JDK.
The JDK must be installed.

\begin{center}\rule{0.5\linewidth}{0.5pt}\end{center}

\subsection{4. Spoken Java Conventions}\label{spoken-java-conventions}

Extend the Python symbol conventions with Java-specific additions:

{\def\LTcaptype{none} % do not increment counter
\begin{longtable}[]{@{}ll@{}}
\toprule\noalign{}
Symbol/Pattern & Spoken Form \\
\midrule\noalign{}
\endhead
\bottomrule\noalign{}
\endlastfoot
\texttt{public} & ``public'' \\
\texttt{private} & ``private'' \\
\texttt{static} & ``static'' \\
\texttt{void} & ``void'' (returns nothing) \\
\texttt{new} & ``new'' \\
\texttt{this.} & ``this dot'' \\
\texttt{super.} & ``super dot'' \\
\texttt{@Override} & ``at-override annotation'' \\
\texttt{//\ comment} & ``double-slash comment'' \\
\texttt{/*\ ...\ */} & ``block comment'' \\
\texttt{/**\ ...\ */} & ``JavaDoc comment'' \\
\texttt{String{[}{]}} & ``String array'' \\
\texttt{ArrayList\textless{}String\textgreater{}} & ``ArrayList of
String'' \\
\texttt{HashMap\textless{}String,\ Integer\textgreater{}} & ``HashMap
from String to Integer'' \\
\texttt{throws\ IOException} & ``throws IOException'' \\
\texttt{System.out.println()} & ``sysout'' or ``system-out-println'' \\
\texttt{System.out.printf()} & ``system-out-printf'' \\
\texttt{Integer.parseInt()} & ``Integer-parse-int'' \\
\texttt{(int)\ x} & ``cast to int'' \\
\texttt{null} & ``null'' \\
\texttt{NullPointerException} & ``null pointer exception'' \\
\texttt{ArrayIndexOutOfBoundsException} & ``array index out of
bounds'' \\
\end{longtable}
}

\textbf{Reading a method signature:}

\begin{Shaded}
\begin{Highlighting}[]
\KeywordTok{public} \DataTypeTok{double} \FunctionTok{calculateArea}\OperatorTok{(}\DataTypeTok{double}\NormalTok{ radius}\OperatorTok{)} \KeywordTok{throws} \BuiltInTok{IllegalArgumentException}
\end{Highlighting}
\end{Shaded}

\emph{``public double method calculateArea, takes a double called
radius, may throw IllegalArgumentException''}

\begin{center}\rule{0.5\linewidth}{0.5pt}\end{center}

\subsection{5. Unit 0--1: Setup and
Foundations}\label{unit-01-setup-and-foundations}

\subsubsection{Lesson 0.1: Installing the
JDK}\label{lesson-0.1-installing-the-jdk}

\paragraph{Instructional Cues}\label{instructional-cues}

\begin{itemize}
\tightlist
\item
  \textbf{Use Adoptium (Eclipse Temurin) rather than Oracle JDK.} The
  Oracle JDK requires account creation and license agreement navigation.
  Adoptium is fully open-source and simpler to download.
\item
  \textbf{Verify the download before running.} Adoptium provides SHA
  checksums. Teaching students to verify downloads is a security habit
  worth establishing.
\item
  \textbf{Windows installer tip:} The Adoptium \texttt{.msi} installer
  includes a checkbox for ``Set JAVA\_HOME variable.'' Ensure this is
  checked. Many Java tools (Maven, Gradle, some IDEs) require
  \texttt{JAVA\_HOME} to be set.
\end{itemize}

\paragraph{JShell as the First
Environment}\label{jshell-as-the-first-environment}

Before any file-based Java, spend time in JShell:

\begin{verbatim}
$ jshell
|  Welcome to JShell -- Version 17.0.x
|  For an introduction type: /help intro

jshell> 5 + 3
$1 ==> 8

jshell> "Hello " + "World"
$2 ==> "Hello World"

jshell> int x = 42;
x ==> 42

jshell> System.out.println("Value: " + x);
Value: 42
\end{verbatim}

JShell gives immediate auditory feedback --- type an expression, hear
the result. This is much more accessible for early exploration than the
compile-run-read cycle.

JShell commands to teach:

\begin{itemize}
\tightlist
\item
  \texttt{/list} --- show all code entered so far
\item
  \texttt{/edit} --- open in editor
\item
  \texttt{/vars} --- show all variables
\item
  \texttt{/methods} --- show all methods defined
\item
  \texttt{/exit} --- quit JShell
\end{itemize}

\paragraph{Thought Experiments}\label{thought-experiments}

\begin{enumerate}
\def\labelenumi{\arabic{enumi}.}
\item
  \emph{``Python is interpreted --- each line runs immediately. Java
  compiles to bytecode first, then runs. What does this mean for the
  kind of errors you encounter, and when? If you have a SyntaxError on
  line 50, in Python you might not discover it until the program reaches
  line 50. In Java, when do you discover it?''}
\item
  \emph{``The Java Virtual Machine (JVM) allows the same Java bytecode
  to run on Windows, macOS, and Linux. The JVM is different on each
  platform, but the bytecode is the same. What does this tell you about
  how Java achieves portability? What alternative approaches could a
  language designer take?''}
\item
  \emph{``JShell is a Read-Eval-Print Loop (REPL). Python has one;
  JavaScript has one; many languages have one. But Java didn't have one
  until 2016. Why might a language's designers add a REPL 20 years after
  the language was created? What did programmers do without it?''}
\item
  \emph{``Java was originally designed for embedded systems (like
  televisions), then became dominant for web servers, then for Android
  apps. Each time it succeeded in a new domain, it required extending
  the language. What does this history suggest about the relationship
  between a language's design and the problems it eventually solves?''}
\end{enumerate}

\begin{center}\rule{0.5\linewidth}{0.5pt}\end{center}

\subsubsection{Lesson 1.2: Your First Java
Program}\label{lesson-1.2-your-first-java-program}

\paragraph{Instructional Cues}\label{instructional-cues-1}

\textbf{Anatomy reading exercise --- before typing any code:} Read the
Hello World program aloud, symbol by symbol, using the pronunciation
table:

\begin{Shaded}
\begin{Highlighting}[]
\KeywordTok{public} \KeywordTok{class}\NormalTok{ HelloWorld }\OperatorTok{\{}
    \KeywordTok{public} \DataTypeTok{static} \DataTypeTok{void} \FunctionTok{main}\OperatorTok{(}\BuiltInTok{String}\OperatorTok{[]}\NormalTok{ args}\OperatorTok{)} \OperatorTok{\{}
        \BuiltInTok{System}\OperatorTok{.}\FunctionTok{out}\OperatorTok{.}\FunctionTok{println}\OperatorTok{(}\StringTok{"Hello, World!"}\OperatorTok{);}
    \OperatorTok{\}}
\OperatorTok{\}}
\end{Highlighting}
\end{Shaded}

\emph{``public class HelloWorld open-brace newline four-spaces public
static void main open-paren String array args close-paren open-brace
newline eight-spaces System dot out dot println open-paren quote Hello
comma World! close-quote close-paren semicolon newline four-spaces
close-brace newline close-brace''}

This exercise establishes pattern recognition for the boilerplate that
the student will type hundreds of times. Making it auditorily familiar
early reduces cognitive load later.

\textbf{Three \texttt{System.out} methods:}

\begin{Shaded}
\begin{Highlighting}[]
\BuiltInTok{System}\OperatorTok{.}\FunctionTok{out}\OperatorTok{.}\FunctionTok{println}\OperatorTok{(}\StringTok{"With newline at end"}\OperatorTok{);}     \CommentTok{// Most common}
\BuiltInTok{System}\OperatorTok{.}\FunctionTok{out}\OperatorTok{.}\FunctionTok{print}\OperatorTok{(}\StringTok{"No newline at end"}\OperatorTok{);}          \CommentTok{// For prompts}
\BuiltInTok{System}\OperatorTok{.}\FunctionTok{out}\OperatorTok{.}\FunctionTok{printf}\OperatorTok{(}\StringTok{"Pi is }\SpecialCharTok{\%.2f\%n}\StringTok{"}\OperatorTok{,} \BuiltInTok{Math}\OperatorTok{.}\FunctionTok{PI}\OperatorTok{);}    \CommentTok{// Formatted}
\end{Highlighting}
\end{Shaded}

For accessible programs, \texttt{printf} is extremely valuable --- it
produces consistent, predictable output formatting that reads cleanly
with a screen reader.

\paragraph{Thought Experiments}\label{thought-experiments-1}

\begin{enumerate}
\def\labelenumi{\arabic{enumi}.}
\item
  \emph{``Every Java program requires a \texttt{public\ class} whose
  name matches the file name. Why does Java enforce this one-to-one
  correspondence? What would break if you could put multiple public
  classes in one file?''}
\item
  \emph{``The \texttt{main} method signature is exactly
  \texttt{public\ static\ void\ main(String{[}{]}\ args)}. It cannot be
  \texttt{public\ void\ main()} or
  \texttt{static\ void\ main(String{[}{]}\ args)}. Why is this signature
  fixed? What does each word contribute?''}
\item
  \emph{``Java uses curly braces \texttt{\{\}} to delimit code blocks;
  Python uses indentation. From an accessibility perspective --- for a
  blind programmer --- which approach provides better information about
  structure when read aloud? Consider that braces can be read as `open
  block' and `close block,' while indentation changes require screen
  reader configuration.''}
\item
  \emph{``\texttt{System.out.println()} is long to type but
  descriptively named. Python's \texttt{print()} is short. What naming
  philosophy does each represent? When does conciseness become
  ambiguity?''}
\end{enumerate}

\begin{center}\rule{0.5\linewidth}{0.5pt}\end{center}

\subsubsection{Lesson 1.3: Variables and Primitive Data
Types}\label{lesson-1.3-variables-and-primitive-data-types}

\paragraph{Instructional Cues}\label{instructional-cues-2}

\textbf{Type declaration as specification:} Frame Java's type
declarations positively: \emph{``When you write
\texttt{int\ score\ =\ 85}, you're telling Java: this variable will
always hold a whole number. Java then protects you --- if you
accidentally try to put a name or a decimal into it, the compiler tells
you immediately.''}

\textbf{The eight primitive types --- prioritize four:} In this
curriculum, students need: \texttt{int}, \texttt{double},
\texttt{boolean}, \texttt{char}. The others (\texttt{byte},
\texttt{short}, \texttt{long}, \texttt{float}) exist but are rarely
needed in beginner programs. Introduce the four primary types fully
before acknowledging the others exist.

\textbf{Type casting demonstration:}

\begin{Shaded}
\begin{Highlighting}[]
\CommentTok{// Widening (automatic)}
\DataTypeTok{int}\NormalTok{ i }\OperatorTok{=} \DecValTok{100}\OperatorTok{;}
\DataTypeTok{double}\NormalTok{ d }\OperatorTok{=}\NormalTok{ i}\OperatorTok{;}   \CommentTok{// 100 {-}\textgreater{} 100.0, automatic}

\CommentTok{// Narrowing (explicit, may lose data)}
\DataTypeTok{double}\NormalTok{ pi }\OperatorTok{=} \FloatTok{3.14159}\OperatorTok{;}
\DataTypeTok{int}\NormalTok{ truncated }\OperatorTok{=} \OperatorTok{(}\DataTypeTok{int}\OperatorTok{)}\NormalTok{ pi}\OperatorTok{;}  \CommentTok{// 3.14159 {-}\textgreater{} 3, data LOST}

\BuiltInTok{System}\OperatorTok{.}\FunctionTok{out}\OperatorTok{.}\FunctionTok{println}\OperatorTok{(}\StringTok{"Truncated: "} \OperatorTok{+}\NormalTok{ truncated}\OperatorTok{);}  \CommentTok{// 3, not 4!}
\end{Highlighting}
\end{Shaded}

The truncation behavior (not rounding) is consistently surprising.
Demonstrate it with \texttt{(int)\ 3.9} producing \texttt{3}, and ask:
``Is this rounding or truncating? What is the difference? When would
truncation cause a bug?''

\paragraph{Non-Obvious Accessibility
Strategies}\label{non-obvious-accessibility-strategies}

\begin{itemize}
\item
  \textbf{\texttt{System.out.printf()} for consistent numeric output:}
  \texttt{printf("Balance:\ \$\%.2f\%n",\ balance)} always produces
  exactly 2 decimal places. Without formatting, \texttt{double} values
  may produce varying numbers of decimal places (\texttt{3.0},
  \texttt{3.14}, \texttt{3.1415926535...}) that read differently each
  time. Consistent formatting makes output more predictable.
\item
  \textbf{\texttt{Scanner.hasNext\_\_\_()} methods:} Before calling
  \texttt{scanner.nextInt()}, check \texttt{scanner.hasNextInt()}. If
  the user typed a non-integer, this returns \texttt{false} and prevents
  the exception, allowing graceful re-prompting. For a blind student who
  may not immediately catch an exception in terminal output, this is
  defensive programming with real accessibility value.
\end{itemize}

\paragraph{Thought Experiments}\label{thought-experiments-2}

\begin{enumerate}
\def\labelenumi{\arabic{enumi}.}
\item
  \emph{``Java's \texttt{int} holds values from about -2.1 billion to
  +2.1 billion. What happens if you exceed this limit (integer
  overflow)? The value silently wraps around to the opposite extreme.
  Name a real-world system where integer overflow has caused a
  documented failure.''}
\item
  \emph{``Java has \texttt{int} (primitive) and \texttt{Integer}
  (object). They hold the same values but behave differently.
  \texttt{int} cannot be null; \texttt{Integer} can. \texttt{int} cannot
  be used in a \texttt{List\textless{}int\textgreater{}} --- you need
  \texttt{List\textless{}Integer\textgreater{}}. Why would a language
  designer make this distinction? What does it cost? What does it
  protect against?''}
\item
  \emph{``Java's \texttt{char} holds a single Unicode character.
  \texttt{char\ c\ =\ \textquotesingle{}A\textquotesingle{}} stores the
  letter A. But \texttt{char} is also a numeric type ---
  \texttt{(int)\ \textquotesingle{}A\textquotesingle{}} is \texttt{65}.
  What does this dual nature (character and number) tell you about how
  computers store text internally?''}
\end{enumerate}

\begin{center}\rule{0.5\linewidth}{0.5pt}\end{center}

\subsection{6. Unit 2: I/O and the Math
Class}\label{unit-2-io-and-the-math-class}

\subsubsection{Lesson 2.2: Scanner for User
Input}\label{lesson-2.2-scanner-for-user-input}

\paragraph{Instructional Cues}\label{instructional-cues-3}

\textbf{The Scanner as a wrapper --- explain the layers:}

\begin{verbatim}
Keyboard
  V produces bytes
System.in (InputStream)
  V raw byte stream
new Scanner(System.in)
  V wraps with text-reading methods
scanner.nextLine(), scanner.nextInt(), etc.
\end{verbatim}

This layered model explains why \texttt{Scanner} is needed rather than
reading from \texttt{System.in} directly. It also foreshadows file
reading where \texttt{new\ Scanner(new\ File(...))} replaces
\texttt{new\ Scanner(System.in)}.

\textbf{\texttt{nextInt()} / \texttt{nextLine()} pitfall --- mandatory
demonstration:}

\begin{Shaded}
\begin{Highlighting}[]
\BuiltInTok{Scanner}\NormalTok{ scanner }\OperatorTok{=} \KeywordTok{new} \BuiltInTok{Scanner}\OperatorTok{(}\BuiltInTok{System}\OperatorTok{.}\FunctionTok{in}\OperatorTok{);}
\BuiltInTok{System}\OperatorTok{.}\FunctionTok{out}\OperatorTok{.}\FunctionTok{print}\OperatorTok{(}\StringTok{"Enter your age: "}\OperatorTok{);}
\DataTypeTok{int}\NormalTok{ age }\OperatorTok{=}\NormalTok{ scanner}\OperatorTok{.}\FunctionTok{nextInt}\OperatorTok{();}          \CommentTok{// Reads "25", leaves "\textbackslash{}n" in buffer}
\BuiltInTok{System}\OperatorTok{.}\FunctionTok{out}\OperatorTok{.}\FunctionTok{print}\OperatorTok{(}\StringTok{"Enter your name: "}\OperatorTok{);}
\BuiltInTok{String}\NormalTok{ name }\OperatorTok{=}\NormalTok{ scanner}\OperatorTok{.}\FunctionTok{nextLine}\OperatorTok{();}     \CommentTok{// Reads the leftover "\textbackslash{}n" {-}{-}{-} gets EMPTY string!}
\BuiltInTok{System}\OperatorTok{.}\FunctionTok{out}\OperatorTok{.}\FunctionTok{println}\OperatorTok{(}\StringTok{"Hello "} \OperatorTok{+}\NormalTok{ name }\OperatorTok{+} \StringTok{", age "} \OperatorTok{+}\NormalTok{ age}\OperatorTok{);}
\CommentTok{// Output: "Hello , age 25" {-}{-}{-} name is empty!}
\end{Highlighting}
\end{Shaded}

The fix:

\begin{Shaded}
\begin{Highlighting}[]
\DataTypeTok{int}\NormalTok{ age }\OperatorTok{=}\NormalTok{ scanner}\OperatorTok{.}\FunctionTok{nextInt}\OperatorTok{();}
\NormalTok{scanner}\OperatorTok{.}\FunctionTok{nextLine}\OperatorTok{();}          \CommentTok{// Consume the leftover newline}
\BuiltInTok{String}\NormalTok{ name }\OperatorTok{=}\NormalTok{ scanner}\OperatorTok{.}\FunctionTok{nextLine}\OperatorTok{();}   \CommentTok{// Now reads properly}
\end{Highlighting}
\end{Shaded}

Run the broken version first. The student should hear the incorrect
output and understand why before seeing the fix.

\paragraph{Thought Experiments}\label{thought-experiments-3}

\begin{enumerate}
\def\labelenumi{\arabic{enumi}.}
\item
  \emph{``The \texttt{Scanner} reads from a `stream' of input. A stream
  is a sequence of data that flows from a source to a destination.
  \texttt{System.in} is the keyboard stream. What other streams exist in
  a typical program? (Think about file streams, network streams, the
  output stream.)''}
\item
  \emph{``After \texttt{scanner.nextInt()}, there's a leftover newline
  in the buffer. What IS a buffer? Why does the operating system buffer
  keyboard input rather than delivering each character the moment a key
  is pressed?''}
\item
  \emph{``In GUI (graphical) programs, users fill out forms with text
  fields and click buttons. There's no \texttt{scanner.nextLine()}. What
  mechanism do GUI programs use to collect user input? How does
  event-driven GUI input differ from sequential console input?''}
\end{enumerate}

\begin{center}\rule{0.5\linewidth}{0.5pt}\end{center}

\subsection{7. Unit 3: Control Flow}\label{unit-3-control-flow}

\subsubsection{Lessons 3.1--3.3}\label{lessons-3.13.3}

\paragraph{Instructional Cues}\label{instructional-cues-4}

\textbf{String comparison --- immediate practice:}

\begin{Shaded}
\begin{Highlighting}[]
\BuiltInTok{String}\NormalTok{ day }\OperatorTok{=} \StringTok{"Wednesday"}\OperatorTok{;}

\CommentTok{// Wrong {-}{-}{-} compares memory addresses}
\ControlFlowTok{if} \OperatorTok{(}\NormalTok{day }\OperatorTok{==} \StringTok{"Wednesday"}\OperatorTok{)} \OperatorTok{\{}
    \BuiltInTok{System}\OperatorTok{.}\FunctionTok{out}\OperatorTok{.}\FunctionTok{println}\OperatorTok{(}\StringTok{"Middle of the week"}\OperatorTok{);}  \CommentTok{// May or may not print}
\OperatorTok{\}}

\CommentTok{// Correct {-}{-}{-} compares content}
\ControlFlowTok{if} \OperatorTok{(}\NormalTok{day}\OperatorTok{.}\FunctionTok{equals}\OperatorTok{(}\StringTok{"Wednesday"}\OperatorTok{))} \OperatorTok{\{}
    \BuiltInTok{System}\OperatorTok{.}\FunctionTok{out}\OperatorTok{.}\FunctionTok{println}\OperatorTok{(}\StringTok{"Middle of the week"}\OperatorTok{);}  \CommentTok{// Always correct}
\OperatorTok{\}}

\CommentTok{// Case{-}insensitive}
\ControlFlowTok{if} \OperatorTok{(}\NormalTok{day}\OperatorTok{.}\FunctionTok{equalsIgnoreCase}\OperatorTok{(}\StringTok{"wednesday"}\OperatorTok{))} \OperatorTok{\{}
    \BuiltInTok{System}\OperatorTok{.}\FunctionTok{out}\OperatorTok{.}\FunctionTok{println}\OperatorTok{(}\StringTok{"This works too"}\OperatorTok{);}
\OperatorTok{\}}
\end{Highlighting}
\end{Shaded}

Create a practice file with 10 comparisons, half using \texttt{==} and
half using \texttt{.equals()}. Ask the student to predict which ones
will work correctly before running.

\textbf{Switch expressions (Java 14+):}

Modern Java supports switch expressions that return a value:

\begin{Shaded}
\begin{Highlighting}[]
\BuiltInTok{String}\NormalTok{ dayType }\OperatorTok{=} \ControlFlowTok{switch} \OperatorTok{(}\NormalTok{day}\OperatorTok{)} \OperatorTok{\{}
    \ControlFlowTok{case} \StringTok{"Monday"}\OperatorTok{,} \StringTok{"Tuesday"}\OperatorTok{,} \StringTok{"Wednesday"}\OperatorTok{,} \StringTok{"Thursday"}\OperatorTok{,} \StringTok{"Friday"} \OperatorTok{{-}\textgreater{}} \StringTok{"Weekday"}\OperatorTok{;}
    \ControlFlowTok{case} \StringTok{"Saturday"}\OperatorTok{,} \StringTok{"Sunday"} \OperatorTok{{-}\textgreater{}} \StringTok{"Weekend"}\OperatorTok{;}
    \KeywordTok{default} \OperatorTok{{-}\textgreater{}} \StringTok{"Unknown"}\OperatorTok{;}
\OperatorTok{\};}
\BuiltInTok{System}\OperatorTok{.}\FunctionTok{out}\OperatorTok{.}\FunctionTok{println}\OperatorTok{(}\NormalTok{day }\OperatorTok{+} \StringTok{" is a "} \OperatorTok{+}\NormalTok{ dayType}\OperatorTok{);}
\end{Highlighting}
\end{Shaded}

This eliminates the fallthrough problem and is more concise. Teach both
traditional \texttt{switch} (for reading legacy code) and switch
expressions (for writing new code).

\paragraph{Thought Experiments}\label{thought-experiments-4}

\begin{enumerate}
\def\labelenumi{\arabic{enumi}.}
\item
  \emph{``Java uses \texttt{\&\&} for AND, \texttt{\textbar{}\textbar{}}
  for OR, \texttt{!} for NOT. Python uses the words \texttt{and},
  \texttt{or}, \texttt{not}. Which approach is more readable when heard
  aloud by a screen reader? Is there a meaningful accessibility argument
  for one over the other?''}
\item
  \emph{``The enhanced for loop \texttt{for\ (String\ item\ :\ list)} is
  simpler to write than the traditional
  \texttt{for\ (int\ i\ =\ 0;\ i\ \textless{}\ list.size();\ i++)}. But
  sometimes you need the index. Name three operations you can do with
  the traditional loop that you cannot do (directly) with the enhanced
  loop.''}
\item
  \emph{``Java's \texttt{switch} statement traditionally falls through
  from one case to the next unless you write \texttt{break}. This
  `feature' has caused countless bugs. Modern Java (14+) has switch
  expressions that don't fall through. What does this change tell you
  about how programming language design evolves in response to
  real-world usage patterns?''}
\item
  \emph{``Short-circuit evaluation:
  \texttt{if\ (obj\ !=\ null\ \&\&\ obj.getValue()\ \textgreater{}\ 0)}
  --- if \texttt{obj} is null, the second condition is never evaluated.
  This prevents a NullPointerException. How does this behavior make
  certain coding patterns possible that would otherwise crash?''}
\end{enumerate}

\begin{center}\rule{0.5\linewidth}{0.5pt}\end{center}

\subsection{8. Unit 4: Arrays and
ArrayLists}\label{unit-4-arrays-and-arraylists}

\subsubsection{Lesson 4.1: Arrays}\label{lesson-4.1-arrays}

\paragraph{Instructional Cues}\label{instructional-cues-5}

\textbf{The fixed-size limitation --- demonstrate, don't just describe:}

\begin{Shaded}
\begin{Highlighting}[]
\BuiltInTok{String}\OperatorTok{[]}\NormalTok{ fruits }\OperatorTok{=} \KeywordTok{new} \BuiltInTok{String}\OperatorTok{[}\DecValTok{3}\OperatorTok{];}
\NormalTok{fruits}\OperatorTok{[}\DecValTok{0}\OperatorTok{]} \OperatorTok{=} \StringTok{"apple"}\OperatorTok{;}
\NormalTok{fruits}\OperatorTok{[}\DecValTok{1}\OperatorTok{]} \OperatorTok{=} \StringTok{"banana"}\OperatorTok{;}
\NormalTok{fruits}\OperatorTok{[}\DecValTok{2}\OperatorTok{]} \OperatorTok{=} \StringTok{"cherry"}\OperatorTok{;}
\NormalTok{fruits}\OperatorTok{[}\DecValTok{3}\OperatorTok{]} \OperatorTok{=} \StringTok{"date"}\OperatorTok{;}   \CommentTok{// ArrayIndexOutOfBoundsException!}
\end{Highlighting}
\end{Shaded}

Let the student hear the exception and analyze the message:
\texttt{Index\ 3\ out\ of\ bounds\ for\ length\ 3}. Ask: ``The message
gives you two numbers. What are they? What does each tell you?''

\textbf{\texttt{Arrays.toString()} as mandatory print method:}

The first time \texttt{print(fruits)} produces
\texttt{{[}Ljava.lang.String;@7a5d012b}, stop and explain: ``This is the
memory address of the array, not its contents. Java arrays don't have a
built-in `print nicely' behavior. Always use
\texttt{Arrays.toString(fruits)} or loop through manually.''

\textbf{The \texttt{for} vs \texttt{for-each} tradeoff --- make it
concrete:}

\begin{Shaded}
\begin{Highlighting}[]
\CommentTok{// for{-}each: simple, clear, but no index}
\ControlFlowTok{for} \OperatorTok{(}\BuiltInTok{String}\NormalTok{ fruit }\OperatorTok{:}\NormalTok{ fruits}\OperatorTok{)} \OperatorTok{\{}
    \BuiltInTok{System}\OperatorTok{.}\FunctionTok{out}\OperatorTok{.}\FunctionTok{println}\OperatorTok{(}\NormalTok{fruit}\OperatorTok{);}
\OperatorTok{\}}

\CommentTok{// traditional: more complex, but gives index}
\ControlFlowTok{for} \OperatorTok{(}\DataTypeTok{int}\NormalTok{ i }\OperatorTok{=} \DecValTok{0}\OperatorTok{;}\NormalTok{ i }\OperatorTok{\textless{}}\NormalTok{ fruits}\OperatorTok{.}\FunctionTok{length}\OperatorTok{;}\NormalTok{ i}\OperatorTok{++)} \OperatorTok{\{}
    \BuiltInTok{System}\OperatorTok{.}\FunctionTok{out}\OperatorTok{.}\FunctionTok{println}\OperatorTok{(}\NormalTok{i }\OperatorTok{+} \StringTok{": "} \OperatorTok{+}\NormalTok{ fruits}\OperatorTok{[}\NormalTok{i}\OperatorTok{]);}
\OperatorTok{\}}

\CommentTok{// The numbered version is more accessible!}
\end{Highlighting}
\end{Shaded}

For blind students, numbered output (showing both index and value) is
generally more useful than unnumbered output. Teach the traditional
\texttt{for} loop for array printing specifically because it enables
indexed output.

\paragraph{Thought Experiments}\label{thought-experiments-5}

\begin{enumerate}
\def\labelenumi{\arabic{enumi}.}
\item
  \emph{``Arrays are stored in contiguous memory --- all elements in a
  row, one after another. This is why accessing by index is fast: the
  address of element \texttt{i} is just the start address plus
  \texttt{i\ ×\ element\ size}. Why does this contiguous storage make
  inserting into the middle of an array slow? What must happen to all
  the elements after the insertion point?''}
\item
  \emph{``If you create
  \texttt{String{[}{]}\ arr\ =\ new\ String{[}5{]}}, what are the
  initial values of \texttt{arr{[}0{]}} through \texttt{arr{[}4{]}}? Now
  create \texttt{int{[}{]}\ nums\ =\ new\ int{[}5{]}} --- what are the
  initial values? Why are the defaults different between String arrays
  and int arrays?''}
\item
  \emph{``Two-dimensional arrays in Java are `arrays of arrays.' This
  means each row can have a different length --- a `jagged array.' When
  would a jagged array be more appropriate than a rectangular 2D array?
  Give a real-world data example.''}
\end{enumerate}

\begin{center}\rule{0.5\linewidth}{0.5pt}\end{center}

\subsubsection{Lesson 4.2: ArrayLists}\label{lesson-4.2-arraylists}

\paragraph{Instructional Cues}\label{instructional-cues-6}

\textbf{Autoboxing --- explain the invisibility:}

\begin{Shaded}
\begin{Highlighting}[]
\BuiltInTok{ArrayList}\OperatorTok{\textless{}}\BuiltInTok{Integer}\OperatorTok{\textgreater{}}\NormalTok{ numbers }\OperatorTok{=} \KeywordTok{new} \BuiltInTok{ArrayList}\OperatorTok{\textless{}\textgreater{}();}
\NormalTok{numbers}\OperatorTok{.}\FunctionTok{add}\OperatorTok{(}\DecValTok{42}\OperatorTok{);}        \CommentTok{// 42 (int) is auto{-}boxed to Integer(42)}
\DataTypeTok{int}\NormalTok{ first }\OperatorTok{=}\NormalTok{ numbers}\OperatorTok{.}\FunctionTok{get}\OperatorTok{(}\DecValTok{0}\OperatorTok{);}  \CommentTok{// Integer(42) is auto{-}unboxed to 42}
\end{Highlighting}
\end{Shaded}

This happens invisibly but has a performance cost. For the course, tell
students: ``Use \texttt{ArrayList\textless{}Integer\textgreater{}} when
you need a flexible list of integers. The boxing/unboxing is automatic
but slightly less efficient than an array of \texttt{int}.''

\textbf{Removal by index vs.~by value:}

\begin{Shaded}
\begin{Highlighting}[]
\BuiltInTok{ArrayList}\OperatorTok{\textless{}}\BuiltInTok{String}\OperatorTok{\textgreater{}}\NormalTok{ list }\OperatorTok{=} \KeywordTok{new} \BuiltInTok{ArrayList}\OperatorTok{\textless{}\textgreater{}(}\BuiltInTok{Arrays}\OperatorTok{.}\FunctionTok{asList}\OperatorTok{(}\StringTok{"a"}\OperatorTok{,} \StringTok{"b"}\OperatorTok{,} \StringTok{"c"}\OperatorTok{,} \StringTok{"b"}\OperatorTok{,} \StringTok{"d"}\OperatorTok{));}
\NormalTok{list}\OperatorTok{.}\FunctionTok{remove}\OperatorTok{(}\DecValTok{1}\OperatorTok{);}            \CommentTok{// Removes by INDEX {-}{-}{-} removes "b" at position 1}
\NormalTok{list}\OperatorTok{.}\FunctionTok{remove}\OperatorTok{(}\StringTok{"b"}\OperatorTok{);}          \CommentTok{// Removes by VALUE {-}{-}{-} removes first "b" found}
\NormalTok{list}\OperatorTok{.}\FunctionTok{removeAll}\OperatorTok{(}\BuiltInTok{Arrays}\OperatorTok{.}\FunctionTok{asList}\OperatorTok{(}\StringTok{"b"}\OperatorTok{,} \StringTok{"c"}\OperatorTok{));}  \CommentTok{// Removes all matching values}
\end{Highlighting}
\end{Shaded}

Ask the student: ``If I call \texttt{list.remove(0)} on a list of
integers \texttt{{[}10,\ 20,\ 30{]}}, does this remove the element at
index 0 (which is \texttt{10}) or remove the value \texttt{0} (which
doesn't exist)? How does Java decide?'' (Answer: removes index 0,
because \texttt{0} is an \texttt{int}, not an \texttt{Integer}.)

\paragraph{Thought Experiments}\label{thought-experiments-6}

\begin{enumerate}
\def\labelenumi{\arabic{enumi}.}
\item
  \emph{``An \texttt{ArrayList} doubles its internal capacity when it
  runs out of space. If you start with capacity 16 and keep adding
  elements, when does reallocation occur? Trace the capacity: 16
  -\textgreater{} 32 -\textgreater{} 64 -\textgreater{} 128\ldots{} What
  is the total cost (in element moves) of adding 1000 elements to an
  initially empty ArrayList this way? Is this efficient?''}
\item
  \emph{``Java generics
  \texttt{ArrayList\textless{}String\textgreater{}} enforce type safety
  at compile time. Python lists can hold any types mixed together. When
  is a typed list (Java style) better? When is a mixed list (Python
  style) more flexible? Which approach produces more reliable
  programs?''}
\end{enumerate}

\begin{center}\rule{0.5\linewidth}{0.5pt}\end{center}

\subsection{9. Unit 5: Strings and String
Processing}\label{unit-5-strings-and-string-processing}

\paragraph{Instructional Cues}\label{instructional-cues-7}

\textbf{String immutability --- the Java version:}

In Java, Strings are immutable for the same reason as Python. But Java's
\texttt{String} object is also \textbf{interned} --- the JVM may reuse
the same \texttt{String} object for two variables holding the same
literal value. This is why \texttt{==} sometimes works for String
comparison (when both refer to the same interned object) but is not
reliable.

\textbf{\texttt{StringBuilder} motivation:}

\begin{Shaded}
\begin{Highlighting}[]
\CommentTok{// Inefficient {-}{-}{-} creates 100,000 String objects}
\BuiltInTok{String}\NormalTok{ result }\OperatorTok{=} \StringTok{""}\OperatorTok{;}
\ControlFlowTok{for} \OperatorTok{(}\DataTypeTok{int}\NormalTok{ i }\OperatorTok{=} \DecValTok{0}\OperatorTok{;}\NormalTok{ i }\OperatorTok{\textless{}} \DecValTok{100000}\OperatorTok{;}\NormalTok{ i}\OperatorTok{++)} \OperatorTok{\{}
\NormalTok{    result }\OperatorTok{+=} \StringTok{"x"}\OperatorTok{;}
\OperatorTok{\}}

\CommentTok{// Efficient {-}{-}{-} modifies a single buffer}
\BuiltInTok{StringBuilder}\NormalTok{ sb }\OperatorTok{=} \KeywordTok{new} \BuiltInTok{StringBuilder}\OperatorTok{();}
\ControlFlowTok{for} \OperatorTok{(}\DataTypeTok{int}\NormalTok{ i }\OperatorTok{=} \DecValTok{0}\OperatorTok{;}\NormalTok{ i }\OperatorTok{\textless{}} \DecValTok{100000}\OperatorTok{;}\NormalTok{ i}\OperatorTok{++)} \OperatorTok{\{}
\NormalTok{    sb}\OperatorTok{.}\FunctionTok{append}\OperatorTok{(}\StringTok{"x"}\OperatorTok{);}
\OperatorTok{\}}
\BuiltInTok{String}\NormalTok{ result }\OperatorTok{=}\NormalTok{ sb}\OperatorTok{.}\FunctionTok{toString}\OperatorTok{();}
\end{Highlighting}
\end{Shaded}

The \texttt{+} in a loop creates a new String object on every iteration,
then discards the old one. For small loops this is fine; for large loops
it's a significant performance problem. \texttt{StringBuilder} maintains
a mutable buffer.

\paragraph{Thought Experiments}\label{thought-experiments-7}

\begin{enumerate}
\def\labelenumi{\arabic{enumi}.}
\item
  \emph{``String concatenation with \texttt{+} is convenient but
  expensive in loops. In Python, the equivalent is also discouraged ---
  use \texttt{\textquotesingle{}\textquotesingle{}.join(list)} instead.
  What does this common advice across two languages tell you about how
  immutable strings interact with loops?''}
\item
  \emph{``Java's \texttt{String.format()} and \texttt{printf()} use
  format specifiers like \texttt{\%s}, \texttt{\%d}, \texttt{\%f}.
  Python uses f-strings with \texttt{\{\}}. What are the accessibility
  tradeoffs of each approach? Which reads more naturally when heard
  aloud? Which is easier to type?''}
\item
  \emph{``Regular expressions are a tiny language inside a programming
  language. \texttt{\textbackslash{}\textbackslash{}d+} matches one or
  more digits. \texttt{{[}a-zA-Z{]}+} matches one or more letters. Why
  do programmers use regular expressions despite their difficult-to-read
  syntax? What would the alternative be for complex pattern matching?''}
\end{enumerate}

\begin{center}\rule{0.5\linewidth}{0.5pt}\end{center}

\subsection{10. Unit 6: Methods}\label{unit-6-methods}

\paragraph{Instructional Cues}\label{instructional-cues-8}

\textbf{Method overloading --- auditory navigation:}

\begin{Shaded}
\begin{Highlighting}[]
\KeywordTok{public} \DataTypeTok{static} \DataTypeTok{int} \FunctionTok{add}\OperatorTok{(}\DataTypeTok{int}\NormalTok{ a}\OperatorTok{,} \DataTypeTok{int}\NormalTok{ b}\OperatorTok{)} \OperatorTok{\{} \ControlFlowTok{return}\NormalTok{ a }\OperatorTok{+}\NormalTok{ b}\OperatorTok{;} \OperatorTok{\}}
\KeywordTok{public} \DataTypeTok{static} \DataTypeTok{double} \FunctionTok{add}\OperatorTok{(}\DataTypeTok{double}\NormalTok{ a}\OperatorTok{,} \DataTypeTok{double}\NormalTok{ b}\OperatorTok{)} \OperatorTok{\{} \ControlFlowTok{return}\NormalTok{ a }\OperatorTok{+}\NormalTok{ b}\OperatorTok{;} \OperatorTok{\}}
\KeywordTok{public} \DataTypeTok{static} \DataTypeTok{int} \FunctionTok{add}\OperatorTok{(}\DataTypeTok{int}\NormalTok{ a}\OperatorTok{,} \DataTypeTok{int}\NormalTok{ b}\OperatorTok{,} \DataTypeTok{int}\NormalTok{ c}\OperatorTok{)} \OperatorTok{\{} \ControlFlowTok{return}\NormalTok{ a }\OperatorTok{+}\NormalTok{ b }\OperatorTok{+}\NormalTok{ c}\OperatorTok{;} \OperatorTok{\}}
\end{Highlighting}
\end{Shaded}

When IntelliSense shows overloads in VS Code, pressing
\texttt{Shift+Alt+F2} (Accessible View) lists them all. Teach students
to use this to discover available overloads.

\textbf{\texttt{void} vs.~return type --- make the distinction
tangible:}

\begin{Shaded}
\begin{Highlighting}[]
\CommentTok{// void: for side effects only}
\KeywordTok{public} \DataTypeTok{static} \DataTypeTok{void} \FunctionTok{printGreeting}\OperatorTok{(}\BuiltInTok{String}\NormalTok{ name}\OperatorTok{)} \OperatorTok{\{}
    \BuiltInTok{System}\OperatorTok{.}\FunctionTok{out}\OperatorTok{.}\FunctionTok{println}\OperatorTok{(}\StringTok{"Hello, "} \OperatorTok{+}\NormalTok{ name }\OperatorTok{+} \StringTok{"!"}\OperatorTok{);}
\OperatorTok{\}}

\CommentTok{// Return type: for computation}
\KeywordTok{public} \DataTypeTok{static} \DataTypeTok{double} \FunctionTok{calculateTip}\OperatorTok{(}\DataTypeTok{double}\NormalTok{ billAmount}\OperatorTok{,} \DataTypeTok{double}\NormalTok{ rate}\OperatorTok{)} \OperatorTok{\{}
    \ControlFlowTok{return}\NormalTok{ billAmount }\OperatorTok{*}\NormalTok{ rate}\OperatorTok{;}
\OperatorTok{\}}

\CommentTok{// Using them:}
\FunctionTok{printGreeting}\OperatorTok{(}\StringTok{"Alice"}\OperatorTok{);}                      \CommentTok{// Standalone call}
\DataTypeTok{double}\NormalTok{ tip }\OperatorTok{=} \FunctionTok{calculateTip}\OperatorTok{(}\FloatTok{50.00}\OperatorTok{,} \FloatTok{0.18}\OperatorTok{);}      \CommentTok{// Must capture return value}
\BuiltInTok{System}\OperatorTok{.}\FunctionTok{out}\OperatorTok{.}\FunctionTok{printf}\OperatorTok{(}\StringTok{"Tip: $}\SpecialCharTok{\%.2f\%n}\StringTok{"}\OperatorTok{,}\NormalTok{ tip}\OperatorTok{);}
\end{Highlighting}
\end{Shaded}

\emph{``\texttt{void} methods exist to DO something. Return-type methods
exist to CALCULATE something and give you the result. Ask: `Do I need to
use the output later in my code?' If yes, it needs a return type.''}

\paragraph{Thought Experiments}\label{thought-experiments-8}

\begin{enumerate}
\def\labelenumi{\arabic{enumi}.}
\item
  \emph{``Java methods can throw exceptions:
  \texttt{public\ double\ sqrt(double\ n)\ throws\ IllegalArgumentException}.
  This signature tells callers that the method might fail in a specific
  way. Python has no equivalent syntax. What is the programmer's
  obligation when they see \texttt{throws} in a method signature?''}
\item
  \emph{``Static methods belong to the class, not to any instance:
  \texttt{Math.sqrt(25)} --- you don't need a \texttt{Math} object.
  Instance methods belong to an object: \texttt{myString.length()}. What
  is the conceptual difference? When should a method be static
  vs.~instance?''}
\item
  \emph{``Method parameters in Java are `pass by value' for primitives
  and `pass by reference' for objects. For primitives: the method gets a
  copy, and modifying the parameter doesn't affect the original. For
  objects: the method gets the reference, and modifications to the
  object's fields DO affect the original. Design a test to confirm this
  behavior for both cases.''}
\item
  \_``Recursion: a method that calls itself.
  \texttt{int\ factorial(int\ n)\ \{\ return\ n\ ==\ 0\ ?\ 1\ :\ n\ \_\ factorial(n-1);\ \}}.
  Every recursive solution has an iterative equivalent. When is the
  recursive version clearer? What is the risk of deep recursion (many
  nested calls)?''*
\end{enumerate}

\begin{center}\rule{0.5\linewidth}{0.5pt}\end{center}

\subsection{11. Unit 7: Classes and OOP}\label{unit-7-classes-and-oop}

\subsubsection{Extended Instructional
Cues}\label{extended-instructional-cues}

\textbf{The mandatory \texttt{\_\_str\_\_} equivalent in Java ---
\texttt{toString()}:}

Every student-defined class must override \texttt{toString()}. Without
it, \texttt{System.out.println(myObject)} produces an unreadable string.
The \texttt{@Override} annotation makes this explicit:

\begin{Shaded}
\begin{Highlighting}[]
\AttributeTok{@Override}
\KeywordTok{public} \BuiltInTok{String} \FunctionTok{toString}\OperatorTok{()} \OperatorTok{\{}
    \ControlFlowTok{return} \BuiltInTok{String}\OperatorTok{.}\FunctionTok{format}\OperatorTok{(}\StringTok{"Dog[name=}\SpecialCharTok{\%s}\StringTok{, age=}\SpecialCharTok{\%d}\StringTok{]"}\OperatorTok{,}\NormalTok{ name}\OperatorTok{,}\NormalTok{ age}\OperatorTok{);}
\OperatorTok{\}}
\end{Highlighting}
\end{Shaded}

Require this in every class from Unit 7 onward. Explain: ``When your
screen reader reads \texttt{System.out.println(myDog)}, it reads
whatever \texttt{toString()} returns. If you don't define
\texttt{toString()}, you'll hear a memory address. That's useless.''

\textbf{JavaDoc comments as accessible documentation:}

Teach this format for every method:

\begin{Shaded}
\begin{Highlighting}[]
\CommentTok{/**}
 \CommentTok{*}\NormalTok{ Calculates the area of this circle}\CommentTok{.}
 \CommentTok{*}
 \CommentTok{*} \AnnotationTok{@return}\NormalTok{ the area as a double}\CommentTok{,}\NormalTok{ in square units}
 \CommentTok{*} \AnnotationTok{@throws}\CommentTok{ }\AnnotationTok{IllegalArgumentException}\CommentTok{ }\NormalTok{if radius is negative}
 \CommentTok{*/}
\KeywordTok{public} \DataTypeTok{double} \FunctionTok{getArea}\OperatorTok{()} \OperatorTok{\{}
    \ControlFlowTok{if} \OperatorTok{(}\NormalTok{radius }\OperatorTok{\textless{}} \DecValTok{0}\OperatorTok{)} \ControlFlowTok{throw} \KeywordTok{new} \BuiltInTok{IllegalArgumentException}\OperatorTok{(}\StringTok{"Radius cannot be negative"}\OperatorTok{);}
    \ControlFlowTok{return} \BuiltInTok{Math}\OperatorTok{.}\FunctionTok{PI} \OperatorTok{*}\NormalTok{ radius }\OperatorTok{*}\NormalTok{ radius}\OperatorTok{;}
\OperatorTok{\}}
\end{Highlighting}
\end{Shaded}

When IntelliSense shows this method in VS Code, pressing
\texttt{Shift+Alt+F2} reads the JavaDoc aloud. This makes the student's
own code navigable via audio --- a profound accessibility benefit.

\textbf{The VS Code symbol navigator:} \texttt{Ctrl+Shift+O} opens a
dropdown listing all methods in the current file. Type a few letters to
filter. This allows jumping to any method by name --- much more
efficient than arrow-key navigation through a long file.

\textbf{Access modifiers --- enforce from day one:}

{\def\LTcaptype{none} % do not increment counter
\begin{longtable}[]{@{}
  >{\raggedright\arraybackslash}p{(\linewidth - 4\tabcolsep) * \real{0.3864}}
  >{\raggedright\arraybackslash}p{(\linewidth - 4\tabcolsep) * \real{0.1250}}
  >{\raggedright\arraybackslash}p{(\linewidth - 4\tabcolsep) * \real{0.4886}}@{}}
\toprule\noalign{}
\begin{minipage}[b]{\linewidth}\raggedright
Modifier
\end{minipage} & \begin{minipage}[b]{\linewidth}\raggedright
Rule
\end{minipage} & \begin{minipage}[b]{\linewidth}\raggedright
Reason
\end{minipage} \\
\midrule\noalign{}
\endhead
\bottomrule\noalign{}
\endlastfoot
\texttt{private} on fields & Always & Encapsulation; prevents direct
manipulation \\
\texttt{public} on getters/setters & When needed & Controlled access \\
\texttt{private} on helper methods & Yes & Internal implementation
detail \\
\texttt{public} on methods the user calls & Yes & Interface to the
class \\
\end{longtable}
}

Establish this as a non-negotiable habit, not a preference.
\emph{``Private fields, public interface. Every time.''}

\paragraph{Thought Experiments}\label{thought-experiments-9}

\begin{enumerate}
\def\labelenumi{\arabic{enumi}.}
\item
  \emph{``Encapsulation means the internal state of an object is
  private. Only its public methods can access or modify it. Why would
  you want to prevent direct access to \texttt{account.balance}? What
  bugs could occur if balance were public?''}
\item
  \emph{``Inheritance creates a hierarchy. A \texttt{SavingsAccount}
  extends \texttt{BankAccount}. Every \texttt{SavingsAccount} is-a
  \texttt{BankAccount}, but not every \texttt{BankAccount} is-a
  \texttt{SavingsAccount}. What are the type-safety implications of
  this? If you have a \texttt{BankAccount} variable, can it hold a
  \texttt{SavingsAccount} object?''}
\item
  \emph{``Java interfaces define contracts --- a list of methods that
  any implementing class must provide. A class can implement multiple
  interfaces but can only extend one class. Why this asymmetry? What
  problem does single inheritance solve? What does it prevent?''}
\item
  \emph{``Design Pattern: the Factory Method creates objects without
  specifying the exact class. Instead of \texttt{new\ Dog()}, you'd call
  \texttt{AnimalFactory.create("dog")}. What is the advantage of this
  indirection? When would it be useful?''}
\item
  \emph{``Object-oriented programming was supposed to model the real
  world. A \texttt{Dog} is-a \texttt{Animal}. But real-world
  relationships are complex --- a \texttt{FlyingFish} is both a
  \texttt{Fish} and can fly like a \texttt{Bird}. How does Java's
  single-inheritance model handle this situation? What alternatives
  exist?''}
\item
  \emph{``Polymorphism means one interface, many implementations.
  \texttt{Animal{[}{]}\ animals\ =\ \{new\ Dog(),\ new\ Cat(),\ new\ Bird()\};\ for\ (Animal\ a\ :\ animals)\ \{\ a.speak();\ \}}
  --- each animal speaks differently. Why is this more flexible than
  checking \texttt{instanceof} and calling different methods
  explicitly?''}
\end{enumerate}

\paragraph{Alternative Resources}\label{alternative-resources}

\begin{itemize}
\tightlist
\item
  \textbf{``Head First Design Patterns'' by Freeman, Robson et al.:}
  Visual book, but audio-describable concepts. Covers fundamental OOP
  design patterns. For advanced students.
\item
  \textbf{``Effective Java'' by Joshua Bloch:} The definitive book on
  Java best practices. Dense but invaluable for understanding WHY Java
  is designed the way it is.
\item
  \textbf{Oracle Java Tutorials (docs.oracle.com/javase/tutorial/):}
  Comprehensive, plain HTML, fully accessible. Covers OOP concepts with
  extensive examples.
\end{itemize}

\begin{center}\rule{0.5\linewidth}{0.5pt}\end{center}

\subsection{12. Unit 8: Files and Exception
Handling}\label{unit-8-files-and-exception-handling}

\paragraph{Instructional Cues}\label{instructional-cues-9}

\textbf{Checked vs.~unchecked exceptions --- the core distinction:}

{\def\LTcaptype{none} % do not increment counter
\begin{longtable}[]{@{}
  >{\raggedright\arraybackslash}p{(\linewidth - 6\tabcolsep) * \real{0.1600}}
  >{\raggedright\arraybackslash}p{(\linewidth - 6\tabcolsep) * \real{0.4200}}
  >{\raggedright\arraybackslash}p{(\linewidth - 6\tabcolsep) * \real{0.1250}}
  >{\raggedright\arraybackslash}p{(\linewidth - 6\tabcolsep) * \real{0.2950}}@{}}
\toprule\noalign{}
\begin{minipage}[b]{\linewidth}\raggedright
Type
\end{minipage} & \begin{minipage}[b]{\linewidth}\raggedright
Examples
\end{minipage} & \begin{minipage}[b]{\linewidth}\raggedright
Must Handle?
\end{minipage} & \begin{minipage}[b]{\linewidth}\raggedright
Why?
\end{minipage} \\
\midrule\noalign{}
\endhead
\bottomrule\noalign{}
\endlastfoot
\textbf{Checked} & \texttt{IOException}, \texttt{FileNotFoundException},
\texttt{SQLException} & Yes --- compiler enforces & These are
recoverable; programmer should handle them \\
\textbf{Unchecked (RuntimeException)} & \texttt{NullPointerException},
\texttt{ArrayIndexOutOfBoundsException},
\texttt{IllegalArgumentException} & No --- optional & Usually indicate
programming errors; may not be recoverable \\
\end{longtable}
}

\emph{``Checked exceptions are like warning labels: `This operation can
fail. You must decide what to do if it does.' Unchecked exceptions are
usually bugs --- the program did something wrong.''}

\textbf{Try-with-resources --- teach first, manual close second:}

\begin{Shaded}
\begin{Highlighting}[]
\CommentTok{// Modern: try{-}with{-}resources (always teach this first)}
\ControlFlowTok{try} \OperatorTok{(}\BuiltInTok{BufferedReader}\NormalTok{ reader }\OperatorTok{=} \KeywordTok{new} \BuiltInTok{BufferedReader}\OperatorTok{(}\KeywordTok{new} \BuiltInTok{FileReader}\OperatorTok{(}\StringTok{"data.txt"}\OperatorTok{)))} \OperatorTok{\{}
    \BuiltInTok{String}\NormalTok{ line}\OperatorTok{;}
    \ControlFlowTok{while} \OperatorTok{((}\NormalTok{line }\OperatorTok{=}\NormalTok{ reader}\OperatorTok{.}\FunctionTok{readLine}\OperatorTok{())} \OperatorTok{!=} \KeywordTok{null}\OperatorTok{)} \OperatorTok{\{}
        \BuiltInTok{System}\OperatorTok{.}\FunctionTok{out}\OperatorTok{.}\FunctionTok{println}\OperatorTok{(}\NormalTok{line}\OperatorTok{);}
    \OperatorTok{\}}
\OperatorTok{\}} \ControlFlowTok{catch} \OperatorTok{(}\BuiltInTok{IOException}\NormalTok{ e}\OperatorTok{)} \OperatorTok{\{}
    \BuiltInTok{System}\OperatorTok{.}\FunctionTok{out}\OperatorTok{.}\FunctionTok{println}\OperatorTok{(}\StringTok{"Error reading file: "} \OperatorTok{+}\NormalTok{ e}\OperatorTok{.}\FunctionTok{getMessage}\OperatorTok{());}
\OperatorTok{\}}
\CommentTok{// reader is automatically closed here, even if an exception occurred}

\CommentTok{// Legacy: manual close (show this so students understand what try{-}with{-}resources does)}
\BuiltInTok{BufferedReader}\NormalTok{ reader }\OperatorTok{=} \KeywordTok{null}\OperatorTok{;}
\ControlFlowTok{try} \OperatorTok{\{}
\NormalTok{    reader }\OperatorTok{=} \KeywordTok{new} \BuiltInTok{BufferedReader}\OperatorTok{(}\KeywordTok{new} \BuiltInTok{FileReader}\OperatorTok{(}\StringTok{"data.txt"}\OperatorTok{));}
    \CommentTok{// ... use reader}
\OperatorTok{\}} \ControlFlowTok{catch} \OperatorTok{(}\BuiltInTok{IOException}\NormalTok{ e}\OperatorTok{)} \OperatorTok{\{}
    \BuiltInTok{System}\OperatorTok{.}\FunctionTok{out}\OperatorTok{.}\FunctionTok{println}\OperatorTok{(}\StringTok{"Error: "} \OperatorTok{+}\NormalTok{ e}\OperatorTok{.}\FunctionTok{getMessage}\OperatorTok{());}
\OperatorTok{\}} \ControlFlowTok{finally} \OperatorTok{\{}
    \ControlFlowTok{if} \OperatorTok{(}\NormalTok{reader }\OperatorTok{!=} \KeywordTok{null}\OperatorTok{)} \OperatorTok{\{}
        \ControlFlowTok{try} \OperatorTok{\{}\NormalTok{ reader}\OperatorTok{.}\FunctionTok{close}\OperatorTok{();} \OperatorTok{\}} \ControlFlowTok{catch} \OperatorTok{(}\BuiltInTok{IOException}\NormalTok{ e}\OperatorTok{)} \OperatorTok{\{} \CommentTok{/* ignore */} \OperatorTok{\}}
    \OperatorTok{\}}
\OperatorTok{\}}
\end{Highlighting}
\end{Shaded}

Show the legacy version to explain why try-with-resources exists and
what it simplifies.

\paragraph{Thought Experiments}\label{thought-experiments-10}

\begin{enumerate}
\def\labelenumi{\arabic{enumi}.}
\item
  \emph{``Checked exceptions force you to handle failures at compile
  time. Critics argue this creates `exception chaining hell' --- every
  method signature becomes littered with \texttt{throws} clauses.
  Defenders argue it makes failure modes explicit and programs more
  reliable. Which side do you find more persuasive? Why?''}
\item
  \emph{``Java's \texttt{Path} and \texttt{Files} classes
  (java.nio.file) are a modern replacement for the older \texttt{File}
  class. The older \texttt{File} class's \texttt{delete()} method
  returns a \texttt{boolean} (true/false) instead of throwing an
  exception. What is the danger of reporting failure via a return value?
  What might a programmer accidentally do wrong?''}
\item
  \emph{``A \texttt{NullPointerException} is called `The Billion Dollar
  Mistake' by its inventor Tony Hoare. Explain why null references cause
  so many bugs. Newer languages (Kotlin, Swift) have built-in null
  safety. What would Java code look like if null didn't exist?''}
\end{enumerate}

\begin{center}\rule{0.5\linewidth}{0.5pt}\end{center}

\subsection{13. Unit 9: Collections, Analysis, and
Testing}\label{unit-9-collections-analysis-and-testing}

\subsubsection{JUnit Testing Deep Dive}\label{junit-testing-deep-dive}

\paragraph{Instructional Cues}\label{instructional-cues-10}

\textbf{Test output is inherently accessible:} JUnit produces plain text
output: \texttt{.} for each passing test, \texttt{F} for failures, with
descriptive failure messages. This is one of the most
screen-reader-friendly outputs in the entire course. Let students
experience testing by running a test suite and listening to the results
before writing their own tests.

\textbf{Test naming convention --- sentence style:}

\begin{Shaded}
\begin{Highlighting}[]
\AttributeTok{@Test}
\KeywordTok{public} \DataTypeTok{void} \FunctionTok{calculateArea\_withPositiveRadius\_returnsCorrectValue}\OperatorTok{()} \OperatorTok{\{} \OperatorTok{\}}

\AttributeTok{@Test}
\KeywordTok{public} \DataTypeTok{void} \FunctionTok{calculateArea\_withZeroRadius\_returnsZero}\OperatorTok{()} \OperatorTok{\{} \OperatorTok{\}}

\AttributeTok{@Test}
\KeywordTok{public} \DataTypeTok{void} \FunctionTok{calculateArea\_withNegativeRadius\_throwsException}\OperatorTok{()} \OperatorTok{\{} \OperatorTok{\}}
\end{Highlighting}
\end{Shaded}

The naming convention \texttt{methodName\_condition\_expectedBehavior}
makes test names read as complete sentences when the test runner
announces them. This is both good practice and an accessibility benefit
--- each test announcement tells you exactly what was tested.

\textbf{\texttt{assertThrows} --- test that exceptions work correctly:}

\begin{Shaded}
\begin{Highlighting}[]
\AttributeTok{@Test}
\KeywordTok{public} \DataTypeTok{void} \FunctionTok{withdraw\_withInsufficientFunds\_throwsException}\OperatorTok{()} \OperatorTok{\{}
\NormalTok{    BankAccount account }\OperatorTok{=} \KeywordTok{new} \FunctionTok{BankAccount}\OperatorTok{(}\FloatTok{100.0}\OperatorTok{);}
    \FunctionTok{assertThrows}\OperatorTok{(}\NormalTok{InsufficientFundsException}\OperatorTok{.}\FunctionTok{class}\OperatorTok{,} \OperatorTok{()} \OperatorTok{{-}\textgreater{}}\NormalTok{ account}\OperatorTok{.}\FunctionTok{withdraw}\OperatorTok{(}\FloatTok{200.0}\OperatorTok{));}
\OperatorTok{\}}
\end{Highlighting}
\end{Shaded}

Testing for exceptions is as important as testing for correct outputs.
Many beginners only test the ``happy path.''

\paragraph{Thought Experiments}\label{thought-experiments-11}

\begin{enumerate}
\def\labelenumi{\arabic{enumi}.}
\item
  \emph{``Test-driven development (TDD) says: write the test first, then
  write the implementation. How does writing a test first change the way
  you think about a method? What is forced to be clear before you write
  a line of implementation?''}
\item
  \emph{``Mocking means replacing real dependencies (like a database or
  file system) with fake versions for testing. Why would you want to
  test a method that reads from a database without using a real
  database? What complexity does mocking introduce?''}
\item
  \emph{``A project has 100\% test coverage --- tests execute every
  line. Is the software correct? Name a bug that 100\% coverage would
  not catch. (Hint: think about a bug that only appears with certain
  combinations of inputs.)''}
\end{enumerate}

\begin{center}\rule{0.5\linewidth}{0.5pt}\end{center}

\subsection{14. Units 10--11: Projects and
Capstone}\label{units-1011-projects-and-capstone}

\subsubsection{Java Project
Architecture}\label{java-project-architecture}

For the capstone, encourage a clean package structure:

\begin{verbatim}
src/
|--- com/
|   +--- student/
|       +--- datamanager/
|           |--- model/
|           |   +--- Contact.java
|           |--- storage/
|           |   |--- DataManager.java
|           |   +--- JsonStorage.java
|           |--- ui/
|           |   +--- ConsoleUI.java
|           +--- Main.java
test/
+--- com/
    +--- student/
        +--- datamanager/
            |--- model/
            |   +--- ContactTest.java
            +--- storage/
                +--- DataManagerTest.java
\end{verbatim}

Navigating packages in VS Code: \texttt{Ctrl+Shift+E} for Explorer
(shows package tree), \texttt{Ctrl+P} to open any file by name,
\texttt{Ctrl+Shift+O} for symbols within the current file.

\subsubsection{IDE-Generated Code
Accessibility}\label{ide-generated-code-accessibility}

VS Code and IntelliJ can generate getters, setters, constructors, and
toString. Teach these as time-savers, but require students to read and
understand all generated code:

\begin{itemize}
\tightlist
\item
  \texttt{Ctrl+Shift+P} -\textgreater{} ``Java: Generate Getters and
  Setters''
\item
  \texttt{Ctrl+Shift+P} -\textgreater{} ``Java: Override/Implement
  Methods''
\item
  \texttt{Ctrl+Shift+P} -\textgreater{} ``Java: Generate toString''
\end{itemize}

Generated code is correct but often not styled for accessibility
(missing JavaDoc, generic names). Require students to add JavaDoc and
verify \texttt{toString()} output is meaningful when heard.

\paragraph{Thought Experiments}\label{thought-experiments-12}

\begin{enumerate}
\def\labelenumi{\arabic{enumi}.}
\item
  \emph{``Java in enterprise software is often criticized for having too
  much `boilerplate' --- code that exists for structural reasons rather
  than to express logic. How does this boilerplate affect the experience
  of maintaining a large Java codebase? What tools and patterns does the
  Java ecosystem provide to reduce it?''}
\item
  \emph{``Your capstone application stores data in a file. If your team
  adds a new field to the data class, existing saved files no longer
  match the new format. This is a `schema migration' problem. How would
  you design your storage layer to handle this gracefully?''}
\end{enumerate}

\begin{center}\rule{0.5\linewidth}{0.5pt}\end{center}

\subsection{15. Bridging from Python to
Java}\label{bridging-from-python-to-java}

For courses that teach Python before Java, use these comparison thought
experiments to bridge the languages:

\textbf{Syntax bridge table:}

{\def\LTcaptype{none} % do not increment counter
\begin{longtable}[]{@{}
  >{\raggedright\arraybackslash}p{(\linewidth - 4\tabcolsep) * \real{0.1638}}
  >{\raggedright\arraybackslash}p{(\linewidth - 4\tabcolsep) * \real{0.4310}}
  >{\raggedright\arraybackslash}p{(\linewidth - 4\tabcolsep) * \real{0.4052}}@{}}
\toprule\noalign{}
\begin{minipage}[b]{\linewidth}\raggedright
Python
\end{minipage} & \begin{minipage}[b]{\linewidth}\raggedright
Java
\end{minipage} & \begin{minipage}[b]{\linewidth}\raggedright
Notes
\end{minipage} \\
\midrule\noalign{}
\endhead
\bottomrule\noalign{}
\endlastfoot
\texttt{print("Hello")} & \texttt{System.out.println("Hello")} & \\
\texttt{x\ =\ 5} & \texttt{int\ x\ =\ 5;} & Java requires type
declaration and semicolon \\
\texttt{if\ condition:} & \texttt{if\ (condition)\ \{} & Parens around
condition; brace instead of colon \\
\texttt{elif} & \texttt{else\ if} & Two words in Java \\
\texttt{for\ item\ in\ list:} & \texttt{for\ (Type\ item\ :\ list)\ \{}
& Java requires type \\
\texttt{def\ func(x):} & \texttt{public\ static\ void\ func(int\ x)\ \{}
& Return type and access modifier required \\
\texttt{return\ x} & \texttt{return\ x;} & Semicolon \\
\texttt{\#\ comment} & \texttt{//\ comment} & \\
\texttt{None} & \texttt{null} & \\
\texttt{True}/\texttt{False} & \texttt{true}/\texttt{false} & Lowercase
in Java \\
\texttt{len(list)} & \texttt{list.size()} & Method call, not function \\
\texttt{list.append(x)} & \texttt{list.add(x)} & \\
\texttt{dict{[}"key"{]}} & \texttt{map.get("key")} & \\
\texttt{try/except} & \texttt{try/catch} & \\
\texttt{import\ math} & \texttt{import\ java.lang.Math;} (or use
\texttt{Math.} directly) & \\
\end{longtable}
}

\textbf{Thought experiments for bridging:}

\begin{enumerate}
\def\labelenumi{\arabic{enumi}.}
\item
  \emph{``In Python you write \texttt{x\ =\ 5} and Python figures out
  the type. In Java you write \texttt{int\ x\ =\ 5} and declare the type
  explicitly. Which approach do you prefer after using both? Which one
  would you choose for a large team project? For a quick experiment?''}
\item
  \emph{``Python lists can hold any mixture of types:
  \texttt{{[}1,\ \textquotesingle{}hello\textquotesingle{},\ True,\ 3.14{]}}.
  Java's \texttt{ArrayList\textless{}String\textgreater{}} can only hold
  Strings. What bugs does Python's flexibility enable? What bugs does it
  prevent compared to Java?''}
\item
  \emph{``Python functions are first-class objects --- you can store
  them in variables, pass them as arguments. In Java (before lambdas),
  this was much more complex. Java added lambdas in 2014. What does it
  mean for a language to add a major feature 19 years after it was
  created?''}
\end{enumerate}

\begin{center}\rule{0.5\linewidth}{0.5pt}\end{center}

\subsection{16. Assessment Strategies}\label{assessment-strategies}

\subsubsection{Alternative Assessment
Forms}\label{alternative-assessment-forms}

\textbf{Code review presentations:} The student explains their code out
loud, as if presenting to a team. The teacher asks questions about
design decisions. This tests understanding more authentically than
multiple-choice questions.

\textbf{Bug-hunting assessments:} Provide a working program with 5--8
intentional bugs of different types. The student must find and fix them
using only their screen reader and development tools. This directly
tests real-world debugging skill.

\textbf{Design walkthroughs:} Give the student a problem and ask them to
design the solution --- class structure, methods, data structures ---
without writing any code. Assess the quality of the design, not the
typing speed.

\textbf{Pair programming sessions:} Two students work on one computer.
One dictates code; the other types. Roles swap every 10--15 minutes.
This surfaces communication and reasoning skills that solo assessments
miss.

\begin{center}\rule{0.5\linewidth}{0.5pt}\end{center}

\subsection{17. Supporting Diverse
Needs}\label{supporting-diverse-needs}

\subsubsection{For Braille Display
Users}\label{for-braille-display-users}

Java's verbose syntax can be particularly burdensome on a braille
display with limited width. Strategies:

\begin{itemize}
\tightlist
\item
  Keep lines under 79 characters (PEP 8 applies to style; the Java
  equivalent is generally under 100).
\item
  Use shorter variable names for very long class hierarchies where the
  full name is unwieldy.
\item
  Encourage use of IntelliJ IDEA, which has better braille display
  integration than VS Code in some configurations.
\end{itemize}

\subsubsection{For Students New to Screen
Readers}\label{for-students-new-to-screen-readers}

If a student is still developing screen reader proficiency:

\begin{itemize}
\tightlist
\item
  Begin with JShell instead of file-based programming. Single-line
  interactions provide immediate feedback.
\item
  Use pair work more extensively until independent navigation is fluent.
\item
  Reduce the number of new concepts per session until screen reader
  fluency stabilizes.
\end{itemize}

\begin{center}\rule{0.5\linewidth}{0.5pt}\end{center}

\subsection{18. Alternative and Supplementary
Resources}\label{alternative-and-supplementary-resources}

\subsubsection{Textbooks and Reading}\label{textbooks-and-reading}

{\def\LTcaptype{none} % do not increment counter
\begin{longtable}[]{@{}
  >{\raggedright\arraybackslash}p{(\linewidth - 6\tabcolsep) * \real{0.3789}}
  >{\raggedright\arraybackslash}p{(\linewidth - 6\tabcolsep) * \real{0.1366}}
  >{\raggedright\arraybackslash}p{(\linewidth - 6\tabcolsep) * \real{0.0807}}
  >{\raggedright\arraybackslash}p{(\linewidth - 6\tabcolsep) * \real{0.4037}}@{}}
\toprule\noalign{}
\begin{minipage}[b]{\linewidth}\raggedright
Resource
\end{minipage} & \begin{minipage}[b]{\linewidth}\raggedright
Level
\end{minipage} & \begin{minipage}[b]{\linewidth}\raggedright
Format
\end{minipage} & \begin{minipage}[b]{\linewidth}\raggedright
Accessibility
\end{minipage} \\
\midrule\noalign{}
\endhead
\bottomrule\noalign{}
\endlastfoot
\textbf{Introduction to Programming Using Java} (Eck, free online) &
Beginner--Intermediate & Free HTML/PDF & HTML version excellent;
comprehensive \\
\textbf{Big Java: Early Objects} (Horstmann) & Beginner & Print/ebook &
Standard university text; clear explanations \\
\textbf{Effective Java, 3rd Ed.} (Bloch) & Intermediate--Advanced &
Print/ebook & Essential for understanding Java's design philosophy \\
\textbf{Head First Java} (Sierra, Bates) & Beginner & Print & Very
visual; but text portions accessible; engaging writing style \\
\textbf{Java: The Complete Reference} (Schildt) & Reference &
Print/ebook & Encyclopedic; good for looking up specific features \\
\textbf{Oracle Java Tutorials} (docs.oracle.com) & Beginner--Advanced &
Free HTML & Excellent HTML accessibility; definitive \\
\textbf{Java Documentation} (javadoc) & Reference & Free HTML & Standard
library docs; fully screen-reader-accessible \\
\end{longtable}
}

\subsubsection{Online Courses and
Platforms}\label{online-courses-and-platforms}

{\def\LTcaptype{none} % do not increment counter
\begin{longtable}[]{@{}
  >{\raggedright\arraybackslash}p{(\linewidth - 6\tabcolsep) * \real{0.2791}}
  >{\raggedright\arraybackslash}p{(\linewidth - 6\tabcolsep) * \real{0.2713}}
  >{\raggedright\arraybackslash}p{(\linewidth - 6\tabcolsep) * \real{0.1008}}
  >{\raggedright\arraybackslash}p{(\linewidth - 6\tabcolsep) * \real{0.3488}}@{}}
\toprule\noalign{}
\begin{minipage}[b]{\linewidth}\raggedright
Platform
\end{minipage} & \begin{minipage}[b]{\linewidth}\raggedright
Course
\end{minipage} & \begin{minipage}[b]{\linewidth}\raggedright
Accessibility
\end{minipage} & \begin{minipage}[b]{\linewidth}\raggedright
Notes
\end{minipage} \\
\midrule\noalign{}
\endhead
\bottomrule\noalign{}
\endlastfoot
\textbf{edX / MITx 6.092} & Introduction to Programming in Java & High &
MIT course; video transcripts available \\
\textbf{Coursera} & Java Programming (Duke) & Medium & Free audit
available \\
\textbf{Codecademy} & Learn Java & Partial & Some UI inaccessible;
editor readable \\
\textbf{MOOC.fi} (University of Helsinki) & Java Programming MOOC & High
& Text-based exercises; excellent accessibility \\
\textbf{CodingBat Java} & Java practice problems & High & Plain HTML; no
complex UI \\
\textbf{Exercism Java Track} & Structured exercises with mentoring &
High & Terminal-based; excellent \\
\textbf{LeetCode} & Algorithm problems & Medium & Code editor
accessible; UI moderate \\
\end{longtable}
}

\textbf{MOOC.fi is especially recommended} for blind students --- the
University of Helsinki's Java course has no-frills HTML exercises and
very clear progression. It is one of the most accessible Java learning
platforms available.

\subsubsection{Video and Audio
Resources}\label{video-and-audio-resources}

{\def\LTcaptype{none} % do not increment counter
\begin{longtable}[]{@{}
  >{\raggedright\arraybackslash}p{(\linewidth - 4\tabcolsep) * \real{0.4182}}
  >{\raggedright\arraybackslash}p{(\linewidth - 4\tabcolsep) * \real{0.3545}}
  >{\raggedright\arraybackslash}p{(\linewidth - 4\tabcolsep) * \real{0.2273}}@{}}
\toprule\noalign{}
\begin{minipage}[b]{\linewidth}\raggedright
Resource
\end{minipage} & \begin{minipage}[b]{\linewidth}\raggedright
Description
\end{minipage} & \begin{minipage}[b]{\linewidth}\raggedright
Accessibility
\end{minipage} \\
\midrule\noalign{}
\endhead
\bottomrule\noalign{}
\endlastfoot
\textbf{CS50 AP} (Harvard) & AP Computer Science A course using Java &
Captions and transcripts \\
\textbf{Derek Banas Java Tutorials} (YouTube) & Concise
concept-by-concept tutorials & Captions available \\
\textbf{Cave of Programming Java} (YouTube/website) & Beginner Java
series & Captions; text summaries \\
\textbf{thenewboston Java} (YouTube) & Large beginner series & Captions
available \\
\textbf{Coding with John} (YouTube) & Clear explanations of Java
concepts & Good pacing for beginners \\
\end{longtable}
}

\begin{center}\rule{0.5\linewidth}{0.5pt}\end{center}

\subsection{19. Quick Reference Sheets}\label{quick-reference-sheets}

\subsubsection{Java Compilation and
Running}\label{java-compilation-and-running}

\begin{Shaded}
\begin{Highlighting}[]
\CommentTok{\# Compile a single file}
\ExtensionTok{javac}\NormalTok{ HelloWorld.java    }\CommentTok{\# Creates HelloWorld.class}

\CommentTok{\# Run compiled program}
\ExtensionTok{java}\NormalTok{ HelloWorld          }\CommentTok{\# No .class extension!}

\CommentTok{\# Compile with output directory}
\ExtensionTok{javac} \AttributeTok{{-}d}\NormalTok{ bin src/HelloWorld.java}

\CommentTok{\# Run JShell}
\ExtensionTok{jshell}                   \CommentTok{\# Interactive REPL}
\ExtensionTok{/exit}                    \CommentTok{\# Exit JShell}

\CommentTok{\# Run tests (if using Maven)}
\ExtensionTok{mvn}\NormalTok{ test}

\CommentTok{\# Run tests (if using Gradle)}
\ExtensionTok{./gradlew}\NormalTok{ test}
\end{Highlighting}
\end{Shaded}

\subsubsection{Error Type Quick
Reference}\label{error-type-quick-reference}

{\def\LTcaptype{none} % do not increment counter
\begin{longtable}[]{@{}
  >{\raggedright\arraybackslash}p{(\linewidth - 4\tabcolsep) * \real{0.4276}}
  >{\raggedright\arraybackslash}p{(\linewidth - 4\tabcolsep) * \real{0.3448}}
  >{\raggedright\arraybackslash}p{(\linewidth - 4\tabcolsep) * \real{0.2276}}@{}}
\toprule\noalign{}
\begin{minipage}[b]{\linewidth}\raggedright
Error Type
\end{minipage} & \begin{minipage}[b]{\linewidth}\raggedright
Common Cause
\end{minipage} & \begin{minipage}[b]{\linewidth}\raggedright
First Check
\end{minipage} \\
\midrule\noalign{}
\endhead
\bottomrule\noalign{}
\endlastfoot
\texttt{cannot\ find\ symbol} & Misspelled name or missing import &
Spelling; check imports \\
\texttt{incompatible\ types} & Assigning wrong type or missing cast &
Type of both sides of assignment \\
\texttt{\textquotesingle{};\textquotesingle{}\ expected} & Missing
semicolon & End of the line indicated \\
\texttt{reached\ end\ of\ file\ while\ parsing} & Missing \texttt{\}} &
Count open braces vs close braces \\
\texttt{non-static\ method\ cannot\ be\ referenced\ from\ a\ static\ context}
& Calling instance method from \texttt{main} without object & Need to
create an object first \\
\texttt{NullPointerException} & Calling method on null reference & Which
variable is null \\
\texttt{ArrayIndexOutOfBoundsException} & Index beyond array length &
Print length and the index used \\
\texttt{ClassCastException} & Invalid cast & Check \texttt{instanceof}
before casting \\
\texttt{StackOverflowError} & Infinite recursion & Verify base case \\
\texttt{NumberFormatException} & \texttt{parseInt()} on non-numeric
string & Check input string before parsing \\
\end{longtable}
}

\subsubsection{Scanner Method Reference}\label{scanner-method-reference}

\begin{Shaded}
\begin{Highlighting}[]
\BuiltInTok{Scanner}\NormalTok{ scanner }\OperatorTok{=} \KeywordTok{new} \BuiltInTok{Scanner}\OperatorTok{(}\BuiltInTok{System}\OperatorTok{.}\FunctionTok{in}\OperatorTok{);}

\BuiltInTok{String}\NormalTok{ line   }\OperatorTok{=}\NormalTok{ scanner}\OperatorTok{.}\FunctionTok{nextLine}\OperatorTok{();}     \CommentTok{// Entire line as String}
\BuiltInTok{String}\NormalTok{ word   }\OperatorTok{=}\NormalTok{ scanner}\OperatorTok{.}\FunctionTok{next}\OperatorTok{();}         \CommentTok{// Single token (whitespace{-}delimited)}
\DataTypeTok{int}\NormalTok{ n         }\OperatorTok{=}\NormalTok{ scanner}\OperatorTok{.}\FunctionTok{nextInt}\OperatorTok{();}      \CommentTok{// Integer}
\DataTypeTok{double}\NormalTok{ d      }\OperatorTok{=}\NormalTok{ scanner}\OperatorTok{.}\FunctionTok{nextDouble}\OperatorTok{();}   \CommentTok{// Double}
\DataTypeTok{boolean}\NormalTok{ b     }\OperatorTok{=}\NormalTok{ scanner}\OperatorTok{.}\FunctionTok{nextBoolean}\OperatorTok{();}  \CommentTok{// Boolean (true/false)}

\CommentTok{// Safe versions with type checking}
\ControlFlowTok{if} \OperatorTok{(}\NormalTok{scanner}\OperatorTok{.}\FunctionTok{hasNextInt}\OperatorTok{())} \OperatorTok{\{}
    \DataTypeTok{int}\NormalTok{ n }\OperatorTok{=}\NormalTok{ scanner}\OperatorTok{.}\FunctionTok{nextInt}\OperatorTok{();}
\OperatorTok{\}}

\CommentTok{// After nextInt(), consume leftover newline:}
\NormalTok{scanner}\OperatorTok{.}\FunctionTok{nextLine}\OperatorTok{();}  \CommentTok{// ALWAYS do this before reading a String}

\CommentTok{// From file}
\BuiltInTok{Scanner}\NormalTok{ fileScanner }\OperatorTok{=} \KeywordTok{new} \BuiltInTok{Scanner}\OperatorTok{(}\KeywordTok{new} \BuiltInTok{File}\OperatorTok{(}\StringTok{"data.txt"}\OperatorTok{));}
\CommentTok{// Same methods work; wrap in try{-}catch for FileNotFoundException}

\NormalTok{scanner}\OperatorTok{.}\FunctionTok{close}\OperatorTok{();}  \CommentTok{// Always close when done}
\end{Highlighting}
\end{Shaded}

\subsubsection{Accessible Output
Template}\label{accessible-output-template}

\begin{Shaded}
\begin{Highlighting}[]
\CommentTok{// Numbered list output}
\BuiltInTok{System}\OperatorTok{.}\FunctionTok{out}\OperatorTok{.}\FunctionTok{println}\OperatorTok{(}\StringTok{"{-}{-}{-} Contacts {-}{-}{-}"}\OperatorTok{);}
\ControlFlowTok{for} \OperatorTok{(}\DataTypeTok{int}\NormalTok{ i }\OperatorTok{=} \DecValTok{0}\OperatorTok{;}\NormalTok{ i }\OperatorTok{\textless{}}\NormalTok{ contacts}\OperatorTok{.}\FunctionTok{size}\OperatorTok{();}\NormalTok{ i}\OperatorTok{++)} \OperatorTok{\{}
    \BuiltInTok{System}\OperatorTok{.}\FunctionTok{out}\OperatorTok{.}\FunctionTok{printf}\OperatorTok{(}\StringTok{"}\SpecialCharTok{\%d}\StringTok{. }\SpecialCharTok{\%s\%n}\StringTok{"}\OperatorTok{,}\NormalTok{ i }\OperatorTok{+} \DecValTok{1}\OperatorTok{,}\NormalTok{ contacts}\OperatorTok{.}\FunctionTok{get}\OperatorTok{(}\NormalTok{i}\OperatorTok{));}
\OperatorTok{\}}
\BuiltInTok{System}\OperatorTok{.}\FunctionTok{out}\OperatorTok{.}\FunctionTok{printf}\OperatorTok{(}\StringTok{"Total: }\SpecialCharTok{\%d}\StringTok{ contacts}\SpecialCharTok{\%n}\StringTok{"}\OperatorTok{,}\NormalTok{ contacts}\OperatorTok{.}\FunctionTok{size}\OperatorTok{());}

\CommentTok{// Menu}
\BuiltInTok{System}\OperatorTok{.}\FunctionTok{out}\OperatorTok{.}\FunctionTok{println}\OperatorTok{(}\StringTok{"}\SpecialCharTok{\textbackslash{}n}\StringTok{{-}{-}{-} Main Menu {-}{-}{-}"}\OperatorTok{);}
\BuiltInTok{System}\OperatorTok{.}\FunctionTok{out}\OperatorTok{.}\FunctionTok{println}\OperatorTok{(}\StringTok{"1. Add contact"}\OperatorTok{);}
\BuiltInTok{System}\OperatorTok{.}\FunctionTok{out}\OperatorTok{.}\FunctionTok{println}\OperatorTok{(}\StringTok{"2. View all contacts"}\OperatorTok{);}
\BuiltInTok{System}\OperatorTok{.}\FunctionTok{out}\OperatorTok{.}\FunctionTok{println}\OperatorTok{(}\StringTok{"3. Search contacts"}\OperatorTok{);}
\BuiltInTok{System}\OperatorTok{.}\FunctionTok{out}\OperatorTok{.}\FunctionTok{println}\OperatorTok{(}\StringTok{"4. Delete contact"}\OperatorTok{);}
\BuiltInTok{System}\OperatorTok{.}\FunctionTok{out}\OperatorTok{.}\FunctionTok{println}\OperatorTok{(}\StringTok{"5. Save and quit"}\OperatorTok{);}
\BuiltInTok{System}\OperatorTok{.}\FunctionTok{out}\OperatorTok{.}\FunctionTok{println}\OperatorTok{(}\StringTok{"{-}{-}{-}{-}{-}{-}{-}{-}{-}{-}{-}{-}{-}{-}{-}{-}{-}"}\OperatorTok{);}
\BuiltInTok{System}\OperatorTok{.}\FunctionTok{out}\OperatorTok{.}\FunctionTok{print}\OperatorTok{(}\StringTok{"Enter option (1{-}5): "}\OperatorTok{);}

\CommentTok{// Confirmation messages}
\BuiltInTok{System}\OperatorTok{.}\FunctionTok{out}\OperatorTok{.}\FunctionTok{printf}\OperatorTok{(}\StringTok{"Contact \textquotesingle{}}\SpecialCharTok{\%s}\StringTok{\textquotesingle{} added successfully.}\SpecialCharTok{\%n}\StringTok{"}\OperatorTok{,}\NormalTok{ name}\OperatorTok{);}
\BuiltInTok{System}\OperatorTok{.}\FunctionTok{out}\OperatorTok{.}\FunctionTok{printf}\OperatorTok{(}\StringTok{"Found }\SpecialCharTok{\%d}\StringTok{ matching records.}\SpecialCharTok{\%n}\StringTok{"}\OperatorTok{,}\NormalTok{ count}\OperatorTok{);}
\BuiltInTok{System}\OperatorTok{.}\FunctionTok{out}\OperatorTok{.}\FunctionTok{println}\OperatorTok{(}\StringTok{"Data saved to contacts.json."}\OperatorTok{);}
\BuiltInTok{System}\OperatorTok{.}\FunctionTok{out}\OperatorTok{.}\FunctionTok{println}\OperatorTok{(}\StringTok{"Error: Name field cannot be empty. Please try again."}\OperatorTok{);}
\end{Highlighting}
\end{Shaded}

\newpage

\section{Unit 0: Setting Up Your Java Development
Environment}\label{unit-0-setting-up-your-java-development-environment}

\subsection{Overview}\label{overview}

This unit walks through installing Java and configuring VS Code for
accessible Java development.

\begin{center}\rule{0.5\linewidth}{0.5pt}\end{center}

\subsection{Lesson 0.1: Installing the Java Development Kit
(JDK)}\label{lesson-0.1-installing-the-java-development-kit-jdk}

\subsubsection{What Is the JDK?}\label{what-is-the-jdk}

The \textbf{JDK (Java Development Kit)} includes everything you need to
write, compile, and run Java programs:

\begin{itemize}
\tightlist
\item
  \textbf{JDK} --- Tools for developers (compiler, debugger).
\item
  \textbf{JVM (Java Virtual Machine)} --- Runs compiled Java programs.
\item
  \textbf{JRE (Java Runtime Environment)} --- Libraries and runtime
  support.
\end{itemize}

\subsubsection{Installing the JDK}\label{installing-the-jdk}

\textbf{Recommended:} Download Eclipse Temurin (OpenJDK) from
\url{https://adoptium.net/}.

Choose \textbf{JDK 17 LTS} (Long Term Support) for stability.

\paragraph{Windows}\label{windows}

\begin{enumerate}
\def\labelenumi{\arabic{enumi}.}
\tightlist
\item
  Download the \texttt{.msi} installer.
\item
  \textbf{CRUCIAL STEP:} In the installer, ensure ``Set JAVA\_HOME
  variable'' and ``Add to PATH'' are checked.
\item
  Activate ``Install'' and then ``Finish''.
\end{enumerate}

\paragraph{macOS}\label{macos}

\begin{enumerate}
\def\labelenumi{\arabic{enumi}.}
\tightlist
\item
  Download the \texttt{.pkg} installer.
\item
  Follow the standard macOS installation steps.
\item
  The JDK is automatically added to your PATH.
\end{enumerate}

\paragraph{Linux (Debian/Ubuntu)}\label{linux-debianubuntu}

\begin{Shaded}
\begin{Highlighting}[]
\FunctionTok{sudo}\NormalTok{ apt update}
\FunctionTok{sudo}\NormalTok{ apt install openjdk{-}17{-}jdk}
\end{Highlighting}
\end{Shaded}

\subsubsection{Verify Installation}\label{verify-installation}

Open a terminal and run:

\begin{Shaded}
\begin{Highlighting}[]
\ExtensionTok{java} \AttributeTok{{-}{-}version}
\ExtensionTok{javac} \AttributeTok{{-}{-}version}
\end{Highlighting}
\end{Shaded}

You should hear the version number (e.g., ``openjdk 17.0.x'').

\subsubsection{Skills Checklist}\label{skills-checklist}

\begin{itemize}
\tightlist
\item[$\square$]
  I have installed JDK 17 or later.
\item[$\square$]
  Running \texttt{java\ -\/-version} in the terminal shows the version
  number.
\item[$\square$]
  Running \texttt{javac\ -\/-version} in the terminal shows the compiler
  version.
\end{itemize}

\begin{center}\rule{0.5\linewidth}{0.5pt}\end{center}

\subsection{Lesson 0.2: Installing VS Code and the Java Extension
Pack}\label{lesson-0.2-installing-vs-code-and-the-java-extension-pack}

\begin{enumerate}
\def\labelenumi{\arabic{enumi}.}
\tightlist
\item
  Install VS Code from \url{https://code.visualstudio.com/}.
\item
  Open VS Code and press \texttt{Ctrl\ +\ Shift\ +\ X} to open
  Extensions.
\item
  Search for \textbf{``Extension Pack for Java''} by Microsoft. Install
  it. This installs: Language Support for Java, Debugger for Java, Java
  Test Runner, and Maven support.
\end{enumerate}

\subsubsection{Configure Screen Reader
Mode}\label{configure-screen-reader-mode}

\begin{itemize}
\tightlist
\item
  Press \texttt{Ctrl\ +\ Shift\ +\ P} -\textgreater{} type ``Toggle
  Screen Reader Optimized'' -\textgreater{} press Enter.
\end{itemize}

\subsubsection{VS Code Keyboard Shortcuts for
Java}\label{vs-code-keyboard-shortcuts-for-java}

{\def\LTcaptype{none} % do not increment counter
\begin{longtable}[]{@{}ll@{}}
\toprule\noalign{}
Action & Shortcut \\
\midrule\noalign{}
\endhead
\bottomrule\noalign{}
\endlastfoot
Run Java file & \texttt{Ctrl\ +\ F5} \\
Debug Java file & \texttt{F5} \\
Command Palette & \texttt{Ctrl\ +\ Shift\ +\ P} \\
Problems View & \texttt{Ctrl\ +\ Shift\ +\ M} \\
Next Error & \texttt{F8} \\
Go to Definition & \texttt{F12} \\
Format Code & \texttt{Shift\ +\ Alt\ +\ F} \\
\end{longtable}
}

\subsubsection{Skills Checklist}\label{skills-checklist-1}

\begin{itemize}
\tightlist
\item[$\square$]
  VS Code is installed with the Extension Pack for Java.
\item[$\square$]
  I can toggle Screen Reader Optimized Mode.
\item[$\square$]
  I can navigate errors using \texttt{F8}.
\end{itemize}

\begin{center}\rule{0.5\linewidth}{0.5pt}\end{center}

\subsection{Lesson 0.3: Your First Java
Project}\label{lesson-0.3-your-first-java-project}

\subsubsection{Creating a Project}\label{creating-a-project}

\begin{enumerate}
\def\labelenumi{\arabic{enumi}.}
\tightlist
\item
  Press \texttt{Ctrl\ +\ Shift\ +\ P} -\textgreater{} ``Java: Create
  Java Project''.
\item
  Select ``No build tools''.
\item
  Choose a folder and name the project \texttt{HelloWorld}.
\item
  VS Code will create a \texttt{src} folder with a sample file.
\end{enumerate}

\subsubsection{Writing Hello World}\label{writing-hello-world}

Create a file \texttt{HelloWorld.java} in the \texttt{src} folder:

\begin{Shaded}
\begin{Highlighting}[]
\KeywordTok{public} \KeywordTok{class}\NormalTok{ HelloWorld }\OperatorTok{\{}
    \KeywordTok{public} \DataTypeTok{static} \DataTypeTok{void} \FunctionTok{main}\OperatorTok{(}\BuiltInTok{String}\OperatorTok{[]}\NormalTok{ args}\OperatorTok{)} \OperatorTok{\{}
        \BuiltInTok{System}\OperatorTok{.}\FunctionTok{out}\OperatorTok{.}\FunctionTok{println}\OperatorTok{(}\StringTok{"Hello, World!"}\OperatorTok{);}
    \OperatorTok{\}}
\OperatorTok{\}}
\end{Highlighting}
\end{Shaded}

\subsubsection{Running the Program}\label{running-the-program}

\begin{itemize}
\tightlist
\item
  Press \texttt{Ctrl\ +\ F5} to run without debugging.
\item
  Or open the terminal (`Ctrl + ``) and type:
\end{itemize}

\begin{Shaded}
\begin{Highlighting}[]
\ExtensionTok{javac}\NormalTok{ HelloWorld.java}
\ExtensionTok{java}\NormalTok{ HelloWorld}
\end{Highlighting}
\end{Shaded}

\subsubsection{Key Differences from
Python}\label{key-differences-from-python}

{\def\LTcaptype{none} % do not increment counter
\begin{longtable}[]{@{}ll@{}}
\toprule\noalign{}
Python & Java \\
\midrule\noalign{}
\endhead
\bottomrule\noalign{}
\endlastfoot
No required structure & Every program needs a \texttt{class} \\
No semicolons & Statements end with \texttt{;} \\
Indentation defines blocks & Curly braces \texttt{\{\}} define blocks \\
\texttt{print()} & \texttt{System.out.println()} \\
Interpreted & Compiled (\texttt{.java} -\textgreater{} \texttt{.class}
-\textgreater{} run) \\
Dynamic typing & Static typing (declare variable types) \\
\end{longtable}
}

\subsubsection{Skills Checklist}\label{skills-checklist-2}

\begin{itemize}
\tightlist
\item[$\square$]
  I can create a new Java project in VS Code.
\item[$\square$]
  I can write and run a Hello World program.
\item[$\square$]
  I understand the basic structure of a Java program.
\end{itemize}

\begin{center}\rule{0.5\linewidth}{0.5pt}\end{center}

\section{Unit 1: Foundations of Computing and Getting Started with
Java}\label{unit-1-foundations-of-computing-and-getting-started-with-java}

\subsection{Overview}\label{overview-1}

This unit introduces computer science concepts and the fundamentals of
Java programming.

\textbf{Objectives:}

\begin{itemize}
\tightlist
\item
  Understand hardware vs.~software.
\item
  Write and run basic Java programs.
\item
  Use variables and data types.
\end{itemize}

\begin{center}\rule{0.5\linewidth}{0.5pt}\end{center}

\subsection{Lesson 1.1: Welcome to Computer
Science}\label{lesson-1.1-welcome-to-computer-science}

\emph{(See Python Unit 1.1 for conceptual content --- hardware,
software, algorithms, and problem-solving are language-agnostic.)}

\subsubsection{Java-Specific Notes}\label{java-specific-notes}

Java was created by James Gosling at Sun Microsystems in 1995. Key
principles:

\begin{itemize}
\tightlist
\item
  \textbf{``Write Once, Run Anywhere''} --- Java programs compile to
  bytecode that runs on any JVM.
\item
  \textbf{Strongly typed} --- Every variable must have a declared type.
\item
  \textbf{Object-Oriented} --- Everything lives inside classes.
\end{itemize}

\begin{center}\rule{0.5\linewidth}{0.5pt}\end{center}

\subsection{Lesson 1.2: Your First Java
Program}\label{lesson-1.2-your-first-java-program-1}

\subsubsection{Anatomy of a Java
Program}\label{anatomy-of-a-java-program}

\begin{Shaded}
\begin{Highlighting}[]
\CommentTok{// This is a comment}
\KeywordTok{public} \KeywordTok{class}\NormalTok{ HelloWorld }\OperatorTok{\{}                    \CommentTok{// Class declaration}
    \KeywordTok{public} \DataTypeTok{static} \DataTypeTok{void} \FunctionTok{main}\OperatorTok{(}\BuiltInTok{String}\OperatorTok{[]}\NormalTok{ args}\OperatorTok{)} \OperatorTok{\{} \CommentTok{// Main method {-}{-}{-} program entry point}
        \BuiltInTok{System}\OperatorTok{.}\FunctionTok{out}\OperatorTok{.}\FunctionTok{println}\OperatorTok{(}\StringTok{"Hello, World!"}\OperatorTok{);} \CommentTok{// Print a line to the console}
    \OperatorTok{\}}
\OperatorTok{\}}
\end{Highlighting}
\end{Shaded}

\textbf{Key rules:}

\begin{itemize}
\tightlist
\item
  The file name \textbf{must match} the class name:
  \texttt{HelloWorld.java}.
\item
  Every Java program starts execution from
  \texttt{public\ static\ void\ main(String{[}{]}\ args)}.
\item
  Statements end with a semicolon \texttt{;}.
\item
  Blocks of code are surrounded by curly braces \texttt{\{\}}.
\end{itemize}

\subsubsection{Output Methods}\label{output-methods}

\begin{Shaded}
\begin{Highlighting}[]
\BuiltInTok{System}\OperatorTok{.}\FunctionTok{out}\OperatorTok{.}\FunctionTok{println}\OperatorTok{(}\StringTok{"This prints with a newline at the end."}\OperatorTok{);}
\BuiltInTok{System}\OperatorTok{.}\FunctionTok{out}\OperatorTok{.}\FunctionTok{print}\OperatorTok{(}\StringTok{"This prints without a newline."}\OperatorTok{);}
\BuiltInTok{System}\OperatorTok{.}\FunctionTok{out}\OperatorTok{.}\FunctionTok{printf}\OperatorTok{(}\StringTok{"My name is }\SpecialCharTok{\%s}\StringTok{ and I am }\SpecialCharTok{\%d}\StringTok{ years old.}\SpecialCharTok{\%n}\StringTok{"}\OperatorTok{,} \StringTok{"Alice"}\OperatorTok{,} \DecValTok{15}\OperatorTok{);}
\end{Highlighting}
\end{Shaded}

\subsubsection{Examples}\label{examples}

\textbf{Example 1.2.1: Simple Greeting}

\begin{Shaded}
\begin{Highlighting}[]
\KeywordTok{public} \KeywordTok{class}\NormalTok{ SimpleGreeting }\OperatorTok{\{}
    \KeywordTok{public} \DataTypeTok{static} \DataTypeTok{void} \FunctionTok{main}\OperatorTok{(}\BuiltInTok{String}\OperatorTok{[]}\NormalTok{ args}\OperatorTok{)} \OperatorTok{\{}
        \BuiltInTok{System}\OperatorTok{.}\FunctionTok{out}\OperatorTok{.}\FunctionTok{println}\OperatorTok{(}\StringTok{"Hello, World!"}\OperatorTok{);}
        \BuiltInTok{System}\OperatorTok{.}\FunctionTok{out}\OperatorTok{.}\FunctionTok{println}\OperatorTok{(}\StringTok{"Welcome to Java programming!"}\OperatorTok{);}
    \OperatorTok{\}}
\OperatorTok{\}}
\end{Highlighting}
\end{Shaded}

\textbf{Example 1.2.2: Using printf for Formatted Output}

\begin{Shaded}
\begin{Highlighting}[]
\KeywordTok{public} \KeywordTok{class}\NormalTok{ FormattedOutput }\OperatorTok{\{}
    \KeywordTok{public} \DataTypeTok{static} \DataTypeTok{void} \FunctionTok{main}\OperatorTok{(}\BuiltInTok{String}\OperatorTok{[]}\NormalTok{ args}\OperatorTok{)} \OperatorTok{\{}
        \BuiltInTok{System}\OperatorTok{.}\FunctionTok{out}\OperatorTok{.}\FunctionTok{printf}\OperatorTok{(}\StringTok{"Pi is approximately }\SpecialCharTok{\%.2f\%n}\StringTok{"}\OperatorTok{,} \FloatTok{3.14159}\OperatorTok{);}
        \BuiltInTok{System}\OperatorTok{.}\FunctionTok{out}\OperatorTok{.}\FunctionTok{printf}\OperatorTok{(}\StringTok{"}\SpecialCharTok{\%{-}15s}\StringTok{ }\SpecialCharTok{\%5d\%n}\StringTok{"}\OperatorTok{,} \StringTok{"Apples:"}\OperatorTok{,} \DecValTok{42}\OperatorTok{);}
    \OperatorTok{\}}
\OperatorTok{\}}
\end{Highlighting}
\end{Shaded}

\subsubsection{Skills Checklist}\label{skills-checklist-3}

\begin{itemize}
\tightlist
\item[$\square$]
  I can write a basic Java program with a \texttt{main} method.
\item[$\square$]
  I can use \texttt{System.out.println()} to display output.
\item[$\square$]
  I understand that Java files must match their class name.
\end{itemize}

\begin{center}\rule{0.5\linewidth}{0.5pt}\end{center}

\subsection{Lesson 1.3: Variables and Primitive Data
Types}\label{lesson-1.3-variables-and-primitive-data-types-1}

\subsubsection{Primitive Data Types}\label{primitive-data-types}

{\def\LTcaptype{none} % do not increment counter
\begin{longtable}[]{@{}
  >{\raggedright\arraybackslash}p{(\linewidth - 6\tabcolsep) * \real{0.1304}}
  >{\raggedright\arraybackslash}p{(\linewidth - 6\tabcolsep) * \real{0.3188}}
  >{\raggedright\arraybackslash}p{(\linewidth - 6\tabcolsep) * \real{0.0870}}
  >{\raggedright\arraybackslash}p{(\linewidth - 6\tabcolsep) * \real{0.4638}}@{}}
\toprule\noalign{}
\begin{minipage}[b]{\linewidth}\raggedright
Type
\end{minipage} & \begin{minipage}[b]{\linewidth}\raggedright
Description
\end{minipage} & \begin{minipage}[b]{\linewidth}\raggedright
Size
\end{minipage} & \begin{minipage}[b]{\linewidth}\raggedright
Example
\end{minipage} \\
\midrule\noalign{}
\endhead
\bottomrule\noalign{}
\endlastfoot
\texttt{int} & Whole numbers & 32-bit & \texttt{int\ age\ =\ 25;} \\
\texttt{long} & Large whole numbers & 64-bit &
\texttt{long\ population\ =\ 8000000000L;} \\
\texttt{double} & Decimal numbers & 64-bit &
\texttt{double\ pi\ =\ 3.14159;} \\
\texttt{float} & Decimal (less precise) & 32-bit &
\texttt{float\ price\ =\ 9.99f;} \\
\texttt{boolean} & True or False & 1-bit &
\texttt{boolean\ isRaining\ =\ false;} \\
\texttt{char} & Single character & 16-bit &
\texttt{char\ grade\ =\ \textquotesingle{}A\textquotesingle{};} \\
\end{longtable}
}

\subsubsection{String (Reference Type)}\label{string-reference-type}

\begin{Shaded}
\begin{Highlighting}[]
\BuiltInTok{String}\NormalTok{ name }\OperatorTok{=} \StringTok{"Alice"}\OperatorTok{;}
\BuiltInTok{String}\NormalTok{ message }\OperatorTok{=} \StringTok{"Hello, "} \OperatorTok{+}\NormalTok{ name }\OperatorTok{+} \StringTok{"!"}\OperatorTok{;}  \CommentTok{// Concatenation}
\end{Highlighting}
\end{Shaded}

\subsubsection{Variable Declaration and
Assignment}\label{variable-declaration-and-assignment}

\begin{Shaded}
\begin{Highlighting}[]
\CommentTok{// Declare and assign in one line}
\DataTypeTok{int}\NormalTok{ score }\OperatorTok{=} \DecValTok{100}\OperatorTok{;}
\DataTypeTok{double}\NormalTok{ price }\OperatorTok{=} \FloatTok{4.99}\OperatorTok{;}
\DataTypeTok{boolean}\NormalTok{ isActive }\OperatorTok{=} \KeywordTok{true}\OperatorTok{;}
\BuiltInTok{String}\NormalTok{ studentName }\OperatorTok{=} \StringTok{"Charlie"}\OperatorTok{;}

\CommentTok{// Declare first, assign later}
\DataTypeTok{int}\NormalTok{ total}\OperatorTok{;}
\NormalTok{total }\OperatorTok{=} \DecValTok{50} \OperatorTok{+} \DecValTok{25}\OperatorTok{;}
\end{Highlighting}
\end{Shaded}

\subsubsection{Constants}\label{constants}

\begin{Shaded}
\begin{Highlighting}[]
\DataTypeTok{final} \DataTypeTok{double}\NormalTok{ PI }\OperatorTok{=} \FloatTok{3.14159265358979}\OperatorTok{;}
\DataTypeTok{final} \DataTypeTok{int}\NormalTok{ MAX\_SIZE }\OperatorTok{=} \DecValTok{100}\OperatorTok{;}
\end{Highlighting}
\end{Shaded}

\subsubsection{Type Conversion}\label{type-conversion}

\begin{Shaded}
\begin{Highlighting}[]
\CommentTok{// Widening (automatic {-}{-}{-} safe)}
\DataTypeTok{int}\NormalTok{ intValue }\OperatorTok{=} \DecValTok{42}\OperatorTok{;}
\DataTypeTok{double}\NormalTok{ doubleValue }\OperatorTok{=}\NormalTok{ intValue}\OperatorTok{;}   \CommentTok{// 42 {-}\textgreater{} 42.0}

\CommentTok{// Narrowing (explicit cast {-}{-}{-} may lose data)}
\DataTypeTok{double}\NormalTok{ d }\OperatorTok{=} \FloatTok{9.99}\OperatorTok{;}
\DataTypeTok{int}\NormalTok{ i }\OperatorTok{=} \OperatorTok{(}\DataTypeTok{int}\OperatorTok{)}\NormalTok{ d}\OperatorTok{;}                 \CommentTok{// 9.99 {-}\textgreater{} 9}

\CommentTok{// String to number}
\DataTypeTok{int}\NormalTok{ parsed }\OperatorTok{=} \BuiltInTok{Integer}\OperatorTok{.}\FunctionTok{parseInt}\OperatorTok{(}\StringTok{"123"}\OperatorTok{);}
\DataTypeTok{double}\NormalTok{ parsedD }\OperatorTok{=} \BuiltInTok{Double}\OperatorTok{.}\FunctionTok{parseDouble}\OperatorTok{(}\StringTok{"3.14"}\OperatorTok{);}

\CommentTok{// Number to String}
\BuiltInTok{String}\NormalTok{ s }\OperatorTok{=} \BuiltInTok{String}\OperatorTok{.}\FunctionTok{valueOf}\OperatorTok{(}\DecValTok{42}\OperatorTok{);}
\BuiltInTok{String}\NormalTok{ s2 }\OperatorTok{=} \BuiltInTok{Integer}\OperatorTok{.}\FunctionTok{toString}\OperatorTok{(}\DecValTok{42}\OperatorTok{);}
\end{Highlighting}
\end{Shaded}

\subsubsection{Arithmetic Operators}\label{arithmetic-operators}

{\def\LTcaptype{none} % do not increment counter
\begin{longtable}[]{@{}lll@{}}
\toprule\noalign{}
Operator & Description & Example \\
\midrule\noalign{}
\endhead
\bottomrule\noalign{}
\endlastfoot
\texttt{+} & Addition & \texttt{5\ +\ 3} = \texttt{8} \\
\texttt{-} & Subtraction & \texttt{10\ -\ 4} = \texttt{6} \\
\texttt{*} & Multiplication & \texttt{7\ *\ 6} = \texttt{42} \\
\texttt{/} & Division & \texttt{15\ /\ 4} = \texttt{3} (integer
division!) \\
\texttt{\%} & Modulo & \texttt{10\ \%\ 3} = \texttt{1} \\
\texttt{++} & Increment & \texttt{count++} \\
\texttt{-\/-} & Decrement & \texttt{count-\/-} \\
\end{longtable}
}

\textbf{Important:} Integer division in Java truncates:
\texttt{15\ /\ 4} = \texttt{3}, not \texttt{3.75}. To get a decimal
result, use \texttt{double}: \texttt{15.0\ /\ 4} = \texttt{3.75}.

\subsubsection{Examples}\label{examples-1}

\textbf{Example 1.3.1: Variables and Arithmetic}

\begin{Shaded}
\begin{Highlighting}[]
\KeywordTok{public} \KeywordTok{class}\NormalTok{ VariablesDemo }\OperatorTok{\{}
    \KeywordTok{public} \DataTypeTok{static} \DataTypeTok{void} \FunctionTok{main}\OperatorTok{(}\BuiltInTok{String}\OperatorTok{[]}\NormalTok{ args}\OperatorTok{)} \OperatorTok{\{}
        \DataTypeTok{int}\NormalTok{ appleCount }\OperatorTok{=} \DecValTok{5}\OperatorTok{;}
        \DataTypeTok{int}\NormalTok{ orangeCount }\OperatorTok{=} \DecValTok{3}\OperatorTok{;}
        \DataTypeTok{int}\NormalTok{ totalFruit }\OperatorTok{=}\NormalTok{ appleCount }\OperatorTok{+}\NormalTok{ orangeCount}\OperatorTok{;}
        \BuiltInTok{System}\OperatorTok{.}\FunctionTok{out}\OperatorTok{.}\FunctionTok{println}\OperatorTok{(}\StringTok{"You have "} \OperatorTok{+}\NormalTok{ totalFruit }\OperatorTok{+} \StringTok{" fruits in total."}\OperatorTok{);}

        \DataTypeTok{double}\NormalTok{ pricePerItem }\OperatorTok{=} \FloatTok{2.50}\OperatorTok{;}
        \DataTypeTok{int}\NormalTok{ quantity }\OperatorTok{=} \DecValTok{4}\OperatorTok{;}
        \DataTypeTok{double}\NormalTok{ totalCost }\OperatorTok{=}\NormalTok{ pricePerItem }\OperatorTok{*}\NormalTok{ quantity}\OperatorTok{;}
        \BuiltInTok{System}\OperatorTok{.}\FunctionTok{out}\OperatorTok{.}\FunctionTok{println}\OperatorTok{(}\StringTok{"The total cost is: "} \OperatorTok{+}\NormalTok{ totalCost}\OperatorTok{);}

        \CommentTok{// Integer division}
        \DataTypeTok{int}\NormalTok{ a }\OperatorTok{=} \DecValTok{10}\OperatorTok{,}\NormalTok{ b }\OperatorTok{=} \DecValTok{3}\OperatorTok{;}
        \BuiltInTok{System}\OperatorTok{.}\FunctionTok{out}\OperatorTok{.}\FunctionTok{println}\OperatorTok{(}\StringTok{"10 / 3 = "} \OperatorTok{+} \OperatorTok{(}\NormalTok{a }\OperatorTok{/}\NormalTok{ b}\OperatorTok{));}       \CommentTok{// 3}
        \BuiltInTok{System}\OperatorTok{.}\FunctionTok{out}\OperatorTok{.}\FunctionTok{println}\OperatorTok{(}\StringTok{"10.0 / 3 = "} \OperatorTok{+} \OperatorTok{(}\FloatTok{10.0} \OperatorTok{/}\NormalTok{ b}\OperatorTok{));}  \CommentTok{// 3.3333...}
        \BuiltInTok{System}\OperatorTok{.}\FunctionTok{out}\OperatorTok{.}\FunctionTok{println}\OperatorTok{(}\StringTok{"10 \% 3 = "} \OperatorTok{+} \OperatorTok{(}\NormalTok{a }\OperatorTok{\%}\NormalTok{ b}\OperatorTok{));}        \CommentTok{// 1}
    \OperatorTok{\}}
\OperatorTok{\}}
\end{Highlighting}
\end{Shaded}

\subsubsection{Activities}\label{activities}

\textbf{Activity 1.3.1: Personal Information}

\begin{Shaded}
\begin{Highlighting}[]
\KeywordTok{public} \KeywordTok{class}\NormalTok{ PersonalInfo }\OperatorTok{\{}
    \KeywordTok{public} \DataTypeTok{static} \DataTypeTok{void} \FunctionTok{main}\OperatorTok{(}\BuiltInTok{String}\OperatorTok{[]}\NormalTok{ args}\OperatorTok{)} \OperatorTok{\{}
        \BuiltInTok{String}\NormalTok{ myName }\OperatorTok{=} \StringTok{"Jamie"}\OperatorTok{;}
        \DataTypeTok{int}\NormalTok{ favoriteNumber }\OperatorTok{=} \DecValTok{7}\OperatorTok{;}
        \DataTypeTok{double}\NormalTok{ heightInMeters }\OperatorTok{=} \FloatTok{1.75}\OperatorTok{;}
        \DataTypeTok{boolean}\NormalTok{ likesPizza }\OperatorTok{=} \KeywordTok{true}\OperatorTok{;}

        \BuiltInTok{System}\OperatorTok{.}\FunctionTok{out}\OperatorTok{.}\FunctionTok{println}\OperatorTok{(}\StringTok{"My name is: "} \OperatorTok{+}\NormalTok{ myName}\OperatorTok{);}
        \BuiltInTok{System}\OperatorTok{.}\FunctionTok{out}\OperatorTok{.}\FunctionTok{println}\OperatorTok{(}\StringTok{"Favorite number: "} \OperatorTok{+}\NormalTok{ favoriteNumber}\OperatorTok{);}
        \BuiltInTok{System}\OperatorTok{.}\FunctionTok{out}\OperatorTok{.}\FunctionTok{printf}\OperatorTok{(}\StringTok{"Height: }\SpecialCharTok{\%.2f}\StringTok{ meters}\SpecialCharTok{\%n}\StringTok{"}\OperatorTok{,}\NormalTok{ heightInMeters}\OperatorTok{);}
        \BuiltInTok{System}\OperatorTok{.}\FunctionTok{out}\OperatorTok{.}\FunctionTok{println}\OperatorTok{(}\StringTok{"Likes pizza: "} \OperatorTok{+}\NormalTok{ likesPizza}\OperatorTok{);}
    \OperatorTok{\}}
\OperatorTok{\}}
\end{Highlighting}
\end{Shaded}

\textbf{Activity 1.3.2: Simple Math}

\begin{Shaded}
\begin{Highlighting}[]
\KeywordTok{public} \KeywordTok{class}\NormalTok{ SimpleMath }\OperatorTok{\{}
    \KeywordTok{public} \DataTypeTok{static} \DataTypeTok{void} \FunctionTok{main}\OperatorTok{(}\BuiltInTok{String}\OperatorTok{[]}\NormalTok{ args}\OperatorTok{)} \OperatorTok{\{}
        \DataTypeTok{int}\NormalTok{ num1 }\OperatorTok{=} \DecValTok{15}\OperatorTok{;}
        \DataTypeTok{int}\NormalTok{ num2 }\OperatorTok{=} \DecValTok{7}\OperatorTok{;}

        \BuiltInTok{System}\OperatorTok{.}\FunctionTok{out}\OperatorTok{.}\FunctionTok{println}\OperatorTok{(}\StringTok{"Sum: "} \OperatorTok{+} \OperatorTok{(}\NormalTok{num1 }\OperatorTok{+}\NormalTok{ num2}\OperatorTok{));}
        \BuiltInTok{System}\OperatorTok{.}\FunctionTok{out}\OperatorTok{.}\FunctionTok{println}\OperatorTok{(}\StringTok{"Difference: "} \OperatorTok{+} \OperatorTok{(}\NormalTok{num1 }\OperatorTok{{-}}\NormalTok{ num2}\OperatorTok{));}
        \BuiltInTok{System}\OperatorTok{.}\FunctionTok{out}\OperatorTok{.}\FunctionTok{println}\OperatorTok{(}\StringTok{"Product: "} \OperatorTok{+} \OperatorTok{(}\NormalTok{num1 }\OperatorTok{*}\NormalTok{ num2}\OperatorTok{));}
        \BuiltInTok{System}\OperatorTok{.}\FunctionTok{out}\OperatorTok{.}\FunctionTok{printf}\OperatorTok{(}\StringTok{"Quotient: }\SpecialCharTok{\%.2f\%n}\StringTok{"}\OperatorTok{,} \OperatorTok{(}\DataTypeTok{double}\OperatorTok{)}\NormalTok{ num1 }\OperatorTok{/}\NormalTok{ num2}\OperatorTok{);}
        \BuiltInTok{System}\OperatorTok{.}\FunctionTok{out}\OperatorTok{.}\FunctionTok{println}\OperatorTok{(}\StringTok{"Remainder: "} \OperatorTok{+} \OperatorTok{(}\NormalTok{num1 }\OperatorTok{\%}\NormalTok{ num2}\OperatorTok{));}
    \OperatorTok{\}}
\OperatorTok{\}}
\end{Highlighting}
\end{Shaded}

\subsubsection{Skills Checklist}\label{skills-checklist-4}

\begin{itemize}
\tightlist
\item[$\square$]
  I can declare variables with the correct Java types.
\item[$\square$]
  I can perform basic arithmetic operations in Java.
\item[$\square$]
  I understand the difference between integer and floating-point
  division.
\item[$\square$]
  I can convert between types using casting and parse methods.
\item[$\square$]
  I can use \texttt{final} to declare constants.
\end{itemize}

\newpage

\section{Unit 2: Writing Simple Programs and
Input/Output}\label{unit-2-writing-simple-programs-and-inputoutput}

\subsection{Overview}\label{overview-2}

This unit covers getting user input, type conversions, and working with
the Math library in Java.

\begin{center}\rule{0.5\linewidth}{0.5pt}\end{center}

\subsection{Lesson 2.1: The Software Development
Process}\label{lesson-2.1-the-software-development-process}

\emph{(See Python Unit 2.1 --- the same design principles apply:
analyze, design, implement, test.)}

\begin{center}\rule{0.5\linewidth}{0.5pt}\end{center}

\subsection{Lesson 2.2: Getting User Input with
Scanner}\label{lesson-2.2-getting-user-input-with-scanner}

\subsubsection{\texorpdfstring{The \texttt{Scanner}
Class}{The Scanner Class}}\label{the-scanner-class}

Java uses the \texttt{Scanner} class to read input from the user:

\begin{Shaded}
\begin{Highlighting}[]
\KeywordTok{import} \ImportTok{java}\OperatorTok{.}\ImportTok{util}\OperatorTok{.}\ImportTok{Scanner}\OperatorTok{;}

\KeywordTok{public} \KeywordTok{class}\NormalTok{ UserInput }\OperatorTok{\{}
    \KeywordTok{public} \DataTypeTok{static} \DataTypeTok{void} \FunctionTok{main}\OperatorTok{(}\BuiltInTok{String}\OperatorTok{[]}\NormalTok{ args}\OperatorTok{)} \OperatorTok{\{}
        \BuiltInTok{Scanner}\NormalTok{ scanner }\OperatorTok{=} \KeywordTok{new} \BuiltInTok{Scanner}\OperatorTok{(}\BuiltInTok{System}\OperatorTok{.}\FunctionTok{in}\OperatorTok{);}

        \BuiltInTok{System}\OperatorTok{.}\FunctionTok{out}\OperatorTok{.}\FunctionTok{print}\OperatorTok{(}\StringTok{"What is your name? "}\OperatorTok{);}
        \BuiltInTok{String}\NormalTok{ name }\OperatorTok{=}\NormalTok{ scanner}\OperatorTok{.}\FunctionTok{nextLine}\OperatorTok{();}

        \BuiltInTok{System}\OperatorTok{.}\FunctionTok{out}\OperatorTok{.}\FunctionTok{print}\OperatorTok{(}\StringTok{"How old are you? "}\OperatorTok{);}
        \DataTypeTok{int}\NormalTok{ age }\OperatorTok{=}\NormalTok{ scanner}\OperatorTok{.}\FunctionTok{nextInt}\OperatorTok{();}

        \BuiltInTok{System}\OperatorTok{.}\FunctionTok{out}\OperatorTok{.}\FunctionTok{println}\OperatorTok{(}\StringTok{"Hello, "} \OperatorTok{+}\NormalTok{ name }\OperatorTok{+} \StringTok{"!"}\OperatorTok{);}
        \BuiltInTok{System}\OperatorTok{.}\FunctionTok{out}\OperatorTok{.}\FunctionTok{println}\OperatorTok{(}\StringTok{"Next year you will be "} \OperatorTok{+} \OperatorTok{(}\NormalTok{age }\OperatorTok{+} \DecValTok{1}\OperatorTok{)} \OperatorTok{+} \StringTok{" years old."}\OperatorTok{);}

\NormalTok{        scanner}\OperatorTok{.}\FunctionTok{close}\OperatorTok{();}
    \OperatorTok{\}}
\OperatorTok{\}}
\end{Highlighting}
\end{Shaded}

\subsubsection{Scanner Methods}\label{scanner-methods}

{\def\LTcaptype{none} % do not increment counter
\begin{longtable}[]{@{}
  >{\raggedright\arraybackslash}p{(\linewidth - 4\tabcolsep) * \real{0.1852}}
  >{\raggedright\arraybackslash}p{(\linewidth - 4\tabcolsep) * \real{0.3704}}
  >{\raggedright\arraybackslash}p{(\linewidth - 4\tabcolsep) * \real{0.4444}}@{}}
\toprule\noalign{}
\begin{minipage}[b]{\linewidth}\raggedright
Method
\end{minipage} & \begin{minipage}[b]{\linewidth}\raggedright
Reads
\end{minipage} & \begin{minipage}[b]{\linewidth}\raggedright
Example
\end{minipage} \\
\midrule\noalign{}
\endhead
\bottomrule\noalign{}
\endlastfoot
\texttt{nextLine()} & Entire line as String &
\texttt{String\ name\ =\ scanner.nextLine();} \\
\texttt{next()} & Single word (up to whitespace) &
\texttt{String\ word\ =\ scanner.next();} \\
\texttt{nextInt()} & Integer & \texttt{int\ n\ =\ scanner.nextInt();} \\
\texttt{nextDouble()} & Double &
\texttt{double\ d\ =\ scanner.nextDouble();} \\
\texttt{nextBoolean()} & Boolean &
\texttt{boolean\ b\ =\ scanner.nextBoolean();} \\
\end{longtable}
}

\subsubsection{\texorpdfstring{Common Pitfall: \texttt{nextInt()} +
\texttt{nextLine()}}{Common Pitfall: nextInt() + nextLine()}}\label{common-pitfall-nextint-nextline}

After calling \texttt{nextInt()}, there is a leftover newline character.
Call \texttt{scanner.nextLine()} to consume it before reading another
string:

\begin{Shaded}
\begin{Highlighting}[]
\DataTypeTok{int}\NormalTok{ age }\OperatorTok{=}\NormalTok{ scanner}\OperatorTok{.}\FunctionTok{nextInt}\OperatorTok{();}
\NormalTok{scanner}\OperatorTok{.}\FunctionTok{nextLine}\OperatorTok{();}            \CommentTok{// Consume leftover newline}
\BuiltInTok{String}\NormalTok{ name }\OperatorTok{=}\NormalTok{ scanner}\OperatorTok{.}\FunctionTok{nextLine}\OperatorTok{();}  \CommentTok{// Now reads correctly}
\end{Highlighting}
\end{Shaded}

\subsubsection{Examples}\label{examples-2}

\textbf{Example 2.2.1: Temperature Converter}

\begin{Shaded}
\begin{Highlighting}[]
\KeywordTok{import} \ImportTok{java}\OperatorTok{.}\ImportTok{util}\OperatorTok{.}\ImportTok{Scanner}\OperatorTok{;}

\KeywordTok{public} \KeywordTok{class}\NormalTok{ TempConverter }\OperatorTok{\{}
    \KeywordTok{public} \DataTypeTok{static} \DataTypeTok{void} \FunctionTok{main}\OperatorTok{(}\BuiltInTok{String}\OperatorTok{[]}\NormalTok{ args}\OperatorTok{)} \OperatorTok{\{}
        \BuiltInTok{Scanner}\NormalTok{ scanner }\OperatorTok{=} \KeywordTok{new} \BuiltInTok{Scanner}\OperatorTok{(}\BuiltInTok{System}\OperatorTok{.}\FunctionTok{in}\OperatorTok{);}

        \BuiltInTok{System}\OperatorTok{.}\FunctionTok{out}\OperatorTok{.}\FunctionTok{print}\OperatorTok{(}\StringTok{"Enter temperature in Fahrenheit: "}\OperatorTok{);}
        \DataTypeTok{double}\NormalTok{ fahrenheit }\OperatorTok{=}\NormalTok{ scanner}\OperatorTok{.}\FunctionTok{nextDouble}\OperatorTok{();}

        \DataTypeTok{double}\NormalTok{ celsius }\OperatorTok{=} \OperatorTok{(}\NormalTok{fahrenheit }\OperatorTok{{-}} \DecValTok{32}\OperatorTok{)} \OperatorTok{*} \FloatTok{5.0} \OperatorTok{/} \FloatTok{9.0}\OperatorTok{;}

        \BuiltInTok{System}\OperatorTok{.}\FunctionTok{out}\OperatorTok{.}\FunctionTok{printf}\OperatorTok{(}\StringTok{"}\SpecialCharTok{\%.1f}\StringTok{°F is }\SpecialCharTok{\%.1f}\StringTok{°C}\SpecialCharTok{\%n}\StringTok{"}\OperatorTok{,}\NormalTok{ fahrenheit}\OperatorTok{,}\NormalTok{ celsius}\OperatorTok{);}
\NormalTok{        scanner}\OperatorTok{.}\FunctionTok{close}\OperatorTok{();}
    \OperatorTok{\}}
\OperatorTok{\}}
\end{Highlighting}
\end{Shaded}

\textbf{Example 2.2.2: Circle Calculator}

\begin{Shaded}
\begin{Highlighting}[]
\KeywordTok{import} \ImportTok{java}\OperatorTok{.}\ImportTok{util}\OperatorTok{.}\ImportTok{Scanner}\OperatorTok{;}

\KeywordTok{public} \KeywordTok{class}\NormalTok{ CircleCalculator }\OperatorTok{\{}
    \KeywordTok{public} \DataTypeTok{static} \DataTypeTok{void} \FunctionTok{main}\OperatorTok{(}\BuiltInTok{String}\OperatorTok{[]}\NormalTok{ args}\OperatorTok{)} \OperatorTok{\{}
        \BuiltInTok{Scanner}\NormalTok{ scanner }\OperatorTok{=} \KeywordTok{new} \BuiltInTok{Scanner}\OperatorTok{(}\BuiltInTok{System}\OperatorTok{.}\FunctionTok{in}\OperatorTok{);}

        \BuiltInTok{System}\OperatorTok{.}\FunctionTok{out}\OperatorTok{.}\FunctionTok{print}\OperatorTok{(}\StringTok{"Enter the radius: "}\OperatorTok{);}
        \DataTypeTok{double}\NormalTok{ radius }\OperatorTok{=}\NormalTok{ scanner}\OperatorTok{.}\FunctionTok{nextDouble}\OperatorTok{();}

        \DataTypeTok{double}\NormalTok{ area }\OperatorTok{=} \BuiltInTok{Math}\OperatorTok{.}\FunctionTok{PI} \OperatorTok{*}\NormalTok{ radius }\OperatorTok{*}\NormalTok{ radius}\OperatorTok{;}
        \DataTypeTok{double}\NormalTok{ circumference }\OperatorTok{=} \DecValTok{2} \OperatorTok{*} \BuiltInTok{Math}\OperatorTok{.}\FunctionTok{PI} \OperatorTok{*}\NormalTok{ radius}\OperatorTok{;}

        \BuiltInTok{System}\OperatorTok{.}\FunctionTok{out}\OperatorTok{.}\FunctionTok{printf}\OperatorTok{(}\StringTok{"Area: }\SpecialCharTok{\%.2f\%n}\StringTok{"}\OperatorTok{,}\NormalTok{ area}\OperatorTok{);}
        \BuiltInTok{System}\OperatorTok{.}\FunctionTok{out}\OperatorTok{.}\FunctionTok{printf}\OperatorTok{(}\StringTok{"Circumference: }\SpecialCharTok{\%.2f\%n}\StringTok{"}\OperatorTok{,}\NormalTok{ circumference}\OperatorTok{);}
\NormalTok{        scanner}\OperatorTok{.}\FunctionTok{close}\OperatorTok{();}
    \OperatorTok{\}}
\OperatorTok{\}}
\end{Highlighting}
\end{Shaded}

\subsubsection{Skills Checklist}\label{skills-checklist-5}

\begin{itemize}
\tightlist
\item[$\square$]
  I can import and create a \texttt{Scanner} object.
\item[$\square$]
  I can read different types of input from the user.
\item[$\square$]
  I can handle the \texttt{nextInt()}/\texttt{nextLine()} newline issue.
\end{itemize}

\begin{center}\rule{0.5\linewidth}{0.5pt}\end{center}

\subsection{Lesson 2.3: The Math Class}\label{lesson-2.3-the-math-class}

\begin{Shaded}
\begin{Highlighting}[]
\CommentTok{// Math class provides static methods {-}{-}{-} no import needed}
\DataTypeTok{double}\NormalTok{ squareRoot }\OperatorTok{=} \BuiltInTok{Math}\OperatorTok{.}\FunctionTok{sqrt}\OperatorTok{(}\FloatTok{25.0}\OperatorTok{);}     \CommentTok{// 5.0}
\DataTypeTok{double}\NormalTok{ power }\OperatorTok{=} \BuiltInTok{Math}\OperatorTok{.}\FunctionTok{pow}\OperatorTok{(}\DecValTok{2}\OperatorTok{,} \DecValTok{10}\OperatorTok{);}           \CommentTok{// 1024.0}
\DataTypeTok{double}\NormalTok{ absolute }\OperatorTok{=} \BuiltInTok{Math}\OperatorTok{.}\FunctionTok{abs}\OperatorTok{({-}}\FloatTok{7.5}\OperatorTok{);}         \CommentTok{// 7.5}
\DataTypeTok{int}\NormalTok{ ceiling }\OperatorTok{=} \OperatorTok{(}\DataTypeTok{int}\OperatorTok{)} \BuiltInTok{Math}\OperatorTok{.}\FunctionTok{ceil}\OperatorTok{(}\FloatTok{4.2}\OperatorTok{);}       \CommentTok{// 5}
\DataTypeTok{int}\NormalTok{ floor }\OperatorTok{=} \OperatorTok{(}\DataTypeTok{int}\OperatorTok{)} \BuiltInTok{Math}\OperatorTok{.}\FunctionTok{floor}\OperatorTok{(}\FloatTok{4.8}\OperatorTok{);}        \CommentTok{// 4}
\DataTypeTok{double}\NormalTok{ rounded }\OperatorTok{=} \BuiltInTok{Math}\OperatorTok{.}\FunctionTok{round}\OperatorTok{(}\FloatTok{4.6}\OperatorTok{);}         \CommentTok{// 5}
\DataTypeTok{double}\NormalTok{ maxVal }\OperatorTok{=} \BuiltInTok{Math}\OperatorTok{.}\FunctionTok{max}\OperatorTok{(}\DecValTok{10}\OperatorTok{,} \DecValTok{20}\OperatorTok{);}         \CommentTok{// 20}
\DataTypeTok{double}\NormalTok{ minVal }\OperatorTok{=} \BuiltInTok{Math}\OperatorTok{.}\FunctionTok{min}\OperatorTok{(}\DecValTok{10}\OperatorTok{,} \DecValTok{20}\OperatorTok{);}         \CommentTok{// 10}
\DataTypeTok{double}\NormalTok{ piValue }\OperatorTok{=} \BuiltInTok{Math}\OperatorTok{.}\FunctionTok{PI}\OperatorTok{;}                 \CommentTok{// 3.14159...}
\DataTypeTok{double}\NormalTok{ eValue }\OperatorTok{=} \BuiltInTok{Math}\OperatorTok{.}\FunctionTok{E}\OperatorTok{;}                   \CommentTok{// 2.71828...}
\DataTypeTok{double}\NormalTok{ sinVal }\OperatorTok{=} \BuiltInTok{Math}\OperatorTok{.}\FunctionTok{sin}\OperatorTok{(}\BuiltInTok{Math}\OperatorTok{.}\FunctionTok{PI} \OperatorTok{/} \DecValTok{2}\OperatorTok{);}   \CommentTok{// 1.0}
\DataTypeTok{double}\NormalTok{ logVal }\OperatorTok{=} \BuiltInTok{Math}\OperatorTok{.}\FunctionTok{log}\OperatorTok{(}\BuiltInTok{Math}\OperatorTok{.}\FunctionTok{E}\OperatorTok{);}        \CommentTok{// 1.0  (natural log)}
\DataTypeTok{double}\NormalTok{ log10 }\OperatorTok{=} \BuiltInTok{Math}\OperatorTok{.}\FunctionTok{log10}\OperatorTok{(}\DecValTok{100}\OperatorTok{);}          \CommentTok{// 2.0}
\end{Highlighting}
\end{Shaded}

\subsubsection{Skills Checklist}\label{skills-checklist-6}

\begin{itemize}
\tightlist
\item[$\square$]
  I can use common \texttt{Math} class methods.
\item[$\square$]
  I can use \texttt{Math.PI} for circle calculations.
\end{itemize}

\begin{center}\rule{0.5\linewidth}{0.5pt}\end{center}

\section{Unit 3: Control Flow --- Decisions and
Loops}\label{unit-3-control-flow-decisions-and-loops}

\subsection{Overview}\label{overview-3}

This unit introduces conditional statements and loops in Java.

\textbf{Objectives:}

\begin{itemize}
\tightlist
\item
  Use \texttt{if}, \texttt{else\ if}, and \texttt{else} for
  decision-making.
\item
  Use \texttt{for}, \texttt{while}, and \texttt{do-while} loops.
\item
  Use \texttt{break} and \texttt{continue}.
\end{itemize}

\begin{center}\rule{0.5\linewidth}{0.5pt}\end{center}

\subsection{\texorpdfstring{Lesson 3.1: Boolean Expressions and
\texttt{if}
Statements}{Lesson 3.1: Boolean Expressions and if Statements}}\label{lesson-3.1-boolean-expressions-and-if-statements}

\subsubsection{Comparison Operators}\label{comparison-operators}

{\def\LTcaptype{none} % do not increment counter
\begin{longtable}[]{@{}lll@{}}
\toprule\noalign{}
Operator & Meaning & Example \\
\midrule\noalign{}
\endhead
\bottomrule\noalign{}
\endlastfoot
\texttt{==} & Equal to & \texttt{age\ ==\ 18} \\
\texttt{!=} & Not equal to & \texttt{score\ !=\ 0} \\
\texttt{\textless{}} & Less than & \texttt{temp\ \textless{}\ 0} \\
\texttt{\textgreater{}} & Greater than &
\texttt{score\ \textgreater{}\ 90} \\
\texttt{\textless{}=} & Less than or equal &
\texttt{age\ \textless{}=\ 17} \\
\texttt{\textgreater{}=} & Greater than or equal &
\texttt{score\ \textgreater{}=\ 60} \\
\end{longtable}
}

\subsubsection{Logical Operators}\label{logical-operators}

{\def\LTcaptype{none} % do not increment counter
\begin{longtable}[]{@{}
  >{\raggedright\arraybackslash}p{(\linewidth - 4\tabcolsep) * \real{0.1111}}
  >{\raggedright\arraybackslash}p{(\linewidth - 4\tabcolsep) * \real{0.3056}}
  >{\raggedright\arraybackslash}p{(\linewidth - 4\tabcolsep) * \real{0.5833}}@{}}
\toprule\noalign{}
\begin{minipage}[b]{\linewidth}\raggedright
Operator
\end{minipage} & \begin{minipage}[b]{\linewidth}\raggedright
Meaning
\end{minipage} & \begin{minipage}[b]{\linewidth}\raggedright
Example
\end{minipage} \\
\midrule\noalign{}
\endhead
\bottomrule\noalign{}
\endlastfoot
\texttt{\&\&} & AND (both true) &
\texttt{age\ \textgreater{}=\ 18\ \&\&\ hasId\ ==\ true} \\
\texttt{\textbackslash{}\textbar{}\textbackslash{}\textbar{}} & OR (at
least one true) &
\texttt{day.equals("Sat")\ \textbackslash{}\textbar{}\textbackslash{}\textbar{}\ day.equals("Sun")} \\
\texttt{!} & NOT (negates) & \texttt{!isRaining} \\
\end{longtable}
}

\subsubsection{\texorpdfstring{\texttt{if-else\ if-else}}{if-else if-else}}\label{if-else-if-else}

\begin{Shaded}
\begin{Highlighting}[]
\DataTypeTok{int}\NormalTok{ score }\OperatorTok{=} \DecValTok{85}\OperatorTok{;}

\ControlFlowTok{if} \OperatorTok{(}\NormalTok{score }\OperatorTok{\textgreater{}=} \DecValTok{90}\OperatorTok{)} \OperatorTok{\{}
    \BuiltInTok{System}\OperatorTok{.}\FunctionTok{out}\OperatorTok{.}\FunctionTok{println}\OperatorTok{(}\StringTok{"Grade: A"}\OperatorTok{);}
\OperatorTok{\}} \ControlFlowTok{else} \ControlFlowTok{if} \OperatorTok{(}\NormalTok{score }\OperatorTok{\textgreater{}=} \DecValTok{80}\OperatorTok{)} \OperatorTok{\{}
    \BuiltInTok{System}\OperatorTok{.}\FunctionTok{out}\OperatorTok{.}\FunctionTok{println}\OperatorTok{(}\StringTok{"Grade: B"}\OperatorTok{);}
\OperatorTok{\}} \ControlFlowTok{else} \ControlFlowTok{if} \OperatorTok{(}\NormalTok{score }\OperatorTok{\textgreater{}=} \DecValTok{70}\OperatorTok{)} \OperatorTok{\{}
    \BuiltInTok{System}\OperatorTok{.}\FunctionTok{out}\OperatorTok{.}\FunctionTok{println}\OperatorTok{(}\StringTok{"Grade: C"}\OperatorTok{);}
\OperatorTok{\}} \ControlFlowTok{else} \OperatorTok{\{}
    \BuiltInTok{System}\OperatorTok{.}\FunctionTok{out}\OperatorTok{.}\FunctionTok{println}\OperatorTok{(}\StringTok{"Grade: F"}\OperatorTok{);}
\OperatorTok{\}}
\end{Highlighting}
\end{Shaded}

\subsubsection{String Comparison}\label{string-comparison}

\textbf{Important:} Never use \texttt{==} to compare Strings. Use
\texttt{.equals()}:

\begin{Shaded}
\begin{Highlighting}[]
\BuiltInTok{String}\NormalTok{ color }\OperatorTok{=} \StringTok{"blue"}\OperatorTok{;}
\ControlFlowTok{if} \OperatorTok{(}\NormalTok{color}\OperatorTok{.}\FunctionTok{equals}\OperatorTok{(}\StringTok{"blue"}\OperatorTok{))} \OperatorTok{\{}
    \BuiltInTok{System}\OperatorTok{.}\FunctionTok{out}\OperatorTok{.}\FunctionTok{println}\OperatorTok{(}\StringTok{"It\textquotesingle{}s blue!"}\OperatorTok{);}
\OperatorTok{\}}
\ControlFlowTok{if} \OperatorTok{(}\NormalTok{color}\OperatorTok{.}\FunctionTok{equalsIgnoreCase}\OperatorTok{(}\StringTok{"BLUE"}\OperatorTok{))} \OperatorTok{\{}
    \BuiltInTok{System}\OperatorTok{.}\FunctionTok{out}\OperatorTok{.}\FunctionTok{println}\OperatorTok{(}\StringTok{"Case{-}insensitive match!"}\OperatorTok{);}
\OperatorTok{\}}
\end{Highlighting}
\end{Shaded}

\subsubsection{\texorpdfstring{\texttt{switch}
Statement}{switch Statement}}\label{switch-statement}

\begin{Shaded}
\begin{Highlighting}[]
\DataTypeTok{int}\NormalTok{ day }\OperatorTok{=} \DecValTok{3}\OperatorTok{;}
\ControlFlowTok{switch} \OperatorTok{(}\NormalTok{day}\OperatorTok{)} \OperatorTok{\{}
    \ControlFlowTok{case} \DecValTok{1}\OperatorTok{:}
        \BuiltInTok{System}\OperatorTok{.}\FunctionTok{out}\OperatorTok{.}\FunctionTok{println}\OperatorTok{(}\StringTok{"Monday"}\OperatorTok{);}
        \ControlFlowTok{break}\OperatorTok{;}
    \ControlFlowTok{case} \DecValTok{2}\OperatorTok{:}
        \BuiltInTok{System}\OperatorTok{.}\FunctionTok{out}\OperatorTok{.}\FunctionTok{println}\OperatorTok{(}\StringTok{"Tuesday"}\OperatorTok{);}
        \ControlFlowTok{break}\OperatorTok{;}
    \ControlFlowTok{case} \DecValTok{3}\OperatorTok{:}
        \BuiltInTok{System}\OperatorTok{.}\FunctionTok{out}\OperatorTok{.}\FunctionTok{println}\OperatorTok{(}\StringTok{"Wednesday"}\OperatorTok{);}
        \ControlFlowTok{break}\OperatorTok{;}
    \KeywordTok{default}\OperatorTok{:}
        \BuiltInTok{System}\OperatorTok{.}\FunctionTok{out}\OperatorTok{.}\FunctionTok{println}\OperatorTok{(}\StringTok{"Other day"}\OperatorTok{);}
\OperatorTok{\}}
\end{Highlighting}
\end{Shaded}

\subsubsection{Examples}\label{examples-3}

\textbf{Example 3.1.1: Age Check}

\begin{Shaded}
\begin{Highlighting}[]
\KeywordTok{import} \ImportTok{java}\OperatorTok{.}\ImportTok{util}\OperatorTok{.}\ImportTok{Scanner}\OperatorTok{;}

\KeywordTok{public} \KeywordTok{class}\NormalTok{ AgeCheck }\OperatorTok{\{}
    \KeywordTok{public} \DataTypeTok{static} \DataTypeTok{void} \FunctionTok{main}\OperatorTok{(}\BuiltInTok{String}\OperatorTok{[]}\NormalTok{ args}\OperatorTok{)} \OperatorTok{\{}
        \BuiltInTok{Scanner}\NormalTok{ scanner }\OperatorTok{=} \KeywordTok{new} \BuiltInTok{Scanner}\OperatorTok{(}\BuiltInTok{System}\OperatorTok{.}\FunctionTok{in}\OperatorTok{);}
        \BuiltInTok{System}\OperatorTok{.}\FunctionTok{out}\OperatorTok{.}\FunctionTok{print}\OperatorTok{(}\StringTok{"Enter your age: "}\OperatorTok{);}
        \DataTypeTok{int}\NormalTok{ age }\OperatorTok{=}\NormalTok{ scanner}\OperatorTok{.}\FunctionTok{nextInt}\OperatorTok{();}

        \ControlFlowTok{if} \OperatorTok{(}\NormalTok{age }\OperatorTok{\textgreater{}=} \DecValTok{18}\OperatorTok{)} \OperatorTok{\{}
            \BuiltInTok{System}\OperatorTok{.}\FunctionTok{out}\OperatorTok{.}\FunctionTok{println}\OperatorTok{(}\StringTok{"You are an adult."}\OperatorTok{);}
        \OperatorTok{\}} \ControlFlowTok{else} \OperatorTok{\{}
            \BuiltInTok{System}\OperatorTok{.}\FunctionTok{out}\OperatorTok{.}\FunctionTok{println}\OperatorTok{(}\StringTok{"You are a minor."}\OperatorTok{);}
        \OperatorTok{\}}
\NormalTok{        scanner}\OperatorTok{.}\FunctionTok{close}\OperatorTok{();}
    \OperatorTok{\}}
\OperatorTok{\}}
\end{Highlighting}
\end{Shaded}

\textbf{Example 3.1.2: Grade Calculator}

\begin{Shaded}
\begin{Highlighting}[]
\KeywordTok{import} \ImportTok{java}\OperatorTok{.}\ImportTok{util}\OperatorTok{.}\ImportTok{Scanner}\OperatorTok{;}

\KeywordTok{public} \KeywordTok{class}\NormalTok{ GradeCalculator }\OperatorTok{\{}
    \KeywordTok{public} \DataTypeTok{static} \DataTypeTok{void} \FunctionTok{main}\OperatorTok{(}\BuiltInTok{String}\OperatorTok{[]}\NormalTok{ args}\OperatorTok{)} \OperatorTok{\{}
        \BuiltInTok{Scanner}\NormalTok{ scanner }\OperatorTok{=} \KeywordTok{new} \BuiltInTok{Scanner}\OperatorTok{(}\BuiltInTok{System}\OperatorTok{.}\FunctionTok{in}\OperatorTok{);}
        \BuiltInTok{System}\OperatorTok{.}\FunctionTok{out}\OperatorTok{.}\FunctionTok{print}\OperatorTok{(}\StringTok{"Enter your score (0{-}100): "}\OperatorTok{);}
        \DataTypeTok{int}\NormalTok{ score }\OperatorTok{=}\NormalTok{ scanner}\OperatorTok{.}\FunctionTok{nextInt}\OperatorTok{();}

        \BuiltInTok{String}\NormalTok{ grade}\OperatorTok{;}
        \ControlFlowTok{if} \OperatorTok{(}\NormalTok{score }\OperatorTok{\textgreater{}=} \DecValTok{90}\OperatorTok{)} \OperatorTok{\{}
\NormalTok{            grade }\OperatorTok{=} \StringTok{"A"}\OperatorTok{;}
        \OperatorTok{\}} \ControlFlowTok{else} \ControlFlowTok{if} \OperatorTok{(}\NormalTok{score }\OperatorTok{\textgreater{}=} \DecValTok{80}\OperatorTok{)} \OperatorTok{\{}
\NormalTok{            grade }\OperatorTok{=} \StringTok{"B"}\OperatorTok{;}
        \OperatorTok{\}} \ControlFlowTok{else} \ControlFlowTok{if} \OperatorTok{(}\NormalTok{score }\OperatorTok{\textgreater{}=} \DecValTok{70}\OperatorTok{)} \OperatorTok{\{}
\NormalTok{            grade }\OperatorTok{=} \StringTok{"C"}\OperatorTok{;}
        \OperatorTok{\}} \ControlFlowTok{else} \ControlFlowTok{if} \OperatorTok{(}\NormalTok{score }\OperatorTok{\textgreater{}=} \DecValTok{60}\OperatorTok{)} \OperatorTok{\{}
\NormalTok{            grade }\OperatorTok{=} \StringTok{"D"}\OperatorTok{;}
        \OperatorTok{\}} \ControlFlowTok{else} \OperatorTok{\{}
\NormalTok{            grade }\OperatorTok{=} \StringTok{"F"}\OperatorTok{;}
        \OperatorTok{\}}

        \BuiltInTok{System}\OperatorTok{.}\FunctionTok{out}\OperatorTok{.}\FunctionTok{println}\OperatorTok{(}\StringTok{"Your grade: "} \OperatorTok{+}\NormalTok{ grade}\OperatorTok{);}
\NormalTok{        scanner}\OperatorTok{.}\FunctionTok{close}\OperatorTok{();}
    \OperatorTok{\}}
\OperatorTok{\}}
\end{Highlighting}
\end{Shaded}

\subsubsection{Skills Checklist}\label{skills-checklist-7}

\begin{itemize}
\tightlist
\item[$\square$]
  I can write \texttt{if}, \texttt{else\ if}, and \texttt{else}
  statements in Java.
\item[$\square$]
  I can use \texttt{\&\&}, \texttt{\textbar{}\textbar{}}, and \texttt{!}
  to combine conditions.
\item[$\square$]
  I use \texttt{.equals()} (not \texttt{==}) to compare Strings.
\item[$\square$]
  I can write a \texttt{switch} statement.
\end{itemize}

\begin{center}\rule{0.5\linewidth}{0.5pt}\end{center}

\subsection{\texorpdfstring{Lesson 3.2: \texttt{for}
Loops}{Lesson 3.2: for Loops}}\label{lesson-3.2-for-loops}

\begin{Shaded}
\begin{Highlighting}[]
\CommentTok{// Basic for loop}
\ControlFlowTok{for} \OperatorTok{(}\DataTypeTok{int}\NormalTok{ i }\OperatorTok{=} \DecValTok{0}\OperatorTok{;}\NormalTok{ i }\OperatorTok{\textless{}} \DecValTok{5}\OperatorTok{;}\NormalTok{ i}\OperatorTok{++)} \OperatorTok{\{}
    \BuiltInTok{System}\OperatorTok{.}\FunctionTok{out}\OperatorTok{.}\FunctionTok{println}\OperatorTok{(}\StringTok{"i = "} \OperatorTok{+}\NormalTok{ i}\OperatorTok{);}
\OperatorTok{\}}

\CommentTok{// Counting down}
\ControlFlowTok{for} \OperatorTok{(}\DataTypeTok{int}\NormalTok{ i }\OperatorTok{=} \DecValTok{10}\OperatorTok{;}\NormalTok{ i }\OperatorTok{\textgreater{}} \DecValTok{0}\OperatorTok{;}\NormalTok{ i}\OperatorTok{{-}{-})} \OperatorTok{\{}
    \BuiltInTok{System}\OperatorTok{.}\FunctionTok{out}\OperatorTok{.}\FunctionTok{println}\OperatorTok{(}\NormalTok{i}\OperatorTok{);}
\OperatorTok{\}}
\BuiltInTok{System}\OperatorTok{.}\FunctionTok{out}\OperatorTok{.}\FunctionTok{println}\OperatorTok{(}\StringTok{"Blast off!"}\OperatorTok{);}

\CommentTok{// Stepping by 2}
\ControlFlowTok{for} \OperatorTok{(}\DataTypeTok{int}\NormalTok{ i }\OperatorTok{=} \DecValTok{0}\OperatorTok{;}\NormalTok{ i }\OperatorTok{\textless{}=} \DecValTok{10}\OperatorTok{;}\NormalTok{ i }\OperatorTok{+=} \DecValTok{2}\OperatorTok{)} \OperatorTok{\{}
    \BuiltInTok{System}\OperatorTok{.}\FunctionTok{out}\OperatorTok{.}\FunctionTok{println}\OperatorTok{(}\NormalTok{i}\OperatorTok{);}
\OperatorTok{\}}
\end{Highlighting}
\end{Shaded}

\subsubsection{\texorpdfstring{Enhanced \texttt{for} Loop
(for-each)}{Enhanced for Loop (for-each)}}\label{enhanced-for-loop-for-each}

\begin{Shaded}
\begin{Highlighting}[]
\BuiltInTok{String}\OperatorTok{[]}\NormalTok{ fruits }\OperatorTok{=} \OperatorTok{\{}\StringTok{"apple"}\OperatorTok{,} \StringTok{"banana"}\OperatorTok{,} \StringTok{"cherry"}\OperatorTok{\};}
\ControlFlowTok{for} \OperatorTok{(}\BuiltInTok{String}\NormalTok{ fruit }\OperatorTok{:}\NormalTok{ fruits}\OperatorTok{)} \OperatorTok{\{}
    \BuiltInTok{System}\OperatorTok{.}\FunctionTok{out}\OperatorTok{.}\FunctionTok{println}\OperatorTok{(}\StringTok{"I like "} \OperatorTok{+}\NormalTok{ fruit }\OperatorTok{+} \StringTok{"."}\OperatorTok{);}
\OperatorTok{\}}

\DataTypeTok{int}\OperatorTok{[]}\NormalTok{ numbers }\OperatorTok{=} \OperatorTok{\{}\DecValTok{1}\OperatorTok{,} \DecValTok{2}\OperatorTok{,} \DecValTok{3}\OperatorTok{,} \DecValTok{4}\OperatorTok{,} \DecValTok{5}\OperatorTok{\};}
\DataTypeTok{int}\NormalTok{ sum }\OperatorTok{=} \DecValTok{0}\OperatorTok{;}
\ControlFlowTok{for} \OperatorTok{(}\DataTypeTok{int}\NormalTok{ num }\OperatorTok{:}\NormalTok{ numbers}\OperatorTok{)} \OperatorTok{\{}
\NormalTok{    sum }\OperatorTok{+=}\NormalTok{ num}\OperatorTok{;}
\OperatorTok{\}}
\BuiltInTok{System}\OperatorTok{.}\FunctionTok{out}\OperatorTok{.}\FunctionTok{println}\OperatorTok{(}\StringTok{"Sum: "} \OperatorTok{+}\NormalTok{ sum}\OperatorTok{);}
\end{Highlighting}
\end{Shaded}

\subsubsection{Examples}\label{examples-4}

\textbf{Example 3.2.1: Multiplication Table}

\begin{Shaded}
\begin{Highlighting}[]
\KeywordTok{public} \KeywordTok{class}\NormalTok{ MultiplicationTable }\OperatorTok{\{}
    \KeywordTok{public} \DataTypeTok{static} \DataTypeTok{void} \FunctionTok{main}\OperatorTok{(}\BuiltInTok{String}\OperatorTok{[]}\NormalTok{ args}\OperatorTok{)} \OperatorTok{\{}
        \DataTypeTok{int}\NormalTok{ n }\OperatorTok{=} \DecValTok{5}\OperatorTok{;}
        \BuiltInTok{System}\OperatorTok{.}\FunctionTok{out}\OperatorTok{.}\FunctionTok{println}\OperatorTok{(}\StringTok{"Multiplication table for "} \OperatorTok{+}\NormalTok{ n }\OperatorTok{+} \StringTok{":"}\OperatorTok{);}
        \ControlFlowTok{for} \OperatorTok{(}\DataTypeTok{int}\NormalTok{ i }\OperatorTok{=} \DecValTok{1}\OperatorTok{;}\NormalTok{ i }\OperatorTok{\textless{}=} \DecValTok{10}\OperatorTok{;}\NormalTok{ i}\OperatorTok{++)} \OperatorTok{\{}
            \BuiltInTok{System}\OperatorTok{.}\FunctionTok{out}\OperatorTok{.}\FunctionTok{println}\OperatorTok{(}\NormalTok{n }\OperatorTok{+} \StringTok{" x "} \OperatorTok{+}\NormalTok{ i }\OperatorTok{+} \StringTok{" = "} \OperatorTok{+} \OperatorTok{(}\NormalTok{n }\OperatorTok{*}\NormalTok{ i}\OperatorTok{));}
        \OperatorTok{\}}
    \OperatorTok{\}}
\OperatorTok{\}}
\end{Highlighting}
\end{Shaded}

\subsubsection{Skills Checklist}\label{skills-checklist-8}

\begin{itemize}
\tightlist
\item[$\square$]
  I can write a basic \texttt{for} loop with initialization, condition,
  and update.
\item[$\square$]
  I can use an enhanced \texttt{for} loop to iterate over an array.
\end{itemize}

\begin{center}\rule{0.5\linewidth}{0.5pt}\end{center}

\subsection{\texorpdfstring{Lesson 3.3: \texttt{while} and
\texttt{do-while}
Loops}{Lesson 3.3: while and do-while Loops}}\label{lesson-3.3-while-and-do-while-loops}

\subsubsection{\texorpdfstring{\texttt{while}
Loop}{while Loop}}\label{while-loop}

\begin{Shaded}
\begin{Highlighting}[]
\DataTypeTok{int}\NormalTok{ count }\OperatorTok{=} \DecValTok{1}\OperatorTok{;}
\ControlFlowTok{while} \OperatorTok{(}\NormalTok{count }\OperatorTok{\textless{}=} \DecValTok{5}\OperatorTok{)} \OperatorTok{\{}
    \BuiltInTok{System}\OperatorTok{.}\FunctionTok{out}\OperatorTok{.}\FunctionTok{println}\OperatorTok{(}\StringTok{"Count: "} \OperatorTok{+}\NormalTok{ count}\OperatorTok{);}
\NormalTok{    count}\OperatorTok{++;}
\OperatorTok{\}}
\end{Highlighting}
\end{Shaded}

\subsubsection{\texorpdfstring{\texttt{do-while}
Loop}{do-while Loop}}\label{do-while-loop}

The \texttt{do-while} loop runs at least once before checking the
condition:

\begin{Shaded}
\begin{Highlighting}[]
\BuiltInTok{Scanner}\NormalTok{ scanner }\OperatorTok{=} \KeywordTok{new} \BuiltInTok{Scanner}\OperatorTok{(}\BuiltInTok{System}\OperatorTok{.}\FunctionTok{in}\OperatorTok{);}
\BuiltInTok{String}\NormalTok{ input}\OperatorTok{;}
\ControlFlowTok{do} \OperatorTok{\{}
    \BuiltInTok{System}\OperatorTok{.}\FunctionTok{out}\OperatorTok{.}\FunctionTok{print}\OperatorTok{(}\StringTok{"Enter a positive number: "}\OperatorTok{);}
    \DataTypeTok{int}\NormalTok{ number }\OperatorTok{=}\NormalTok{ scanner}\OperatorTok{.}\FunctionTok{nextInt}\OperatorTok{();}
    \BuiltInTok{System}\OperatorTok{.}\FunctionTok{out}\OperatorTok{.}\FunctionTok{println}\OperatorTok{(}\StringTok{"You entered: "} \OperatorTok{+}\NormalTok{ number}\OperatorTok{);}
    \BuiltInTok{System}\OperatorTok{.}\FunctionTok{out}\OperatorTok{.}\FunctionTok{print}\OperatorTok{(}\StringTok{"Continue? (yes/no): "}\OperatorTok{);}
\NormalTok{    input }\OperatorTok{=}\NormalTok{ scanner}\OperatorTok{.}\FunctionTok{next}\OperatorTok{();}
\OperatorTok{\}} \ControlFlowTok{while} \OperatorTok{(}\NormalTok{input}\OperatorTok{.}\FunctionTok{equals}\OperatorTok{(}\StringTok{"yes"}\OperatorTok{));}
\end{Highlighting}
\end{Shaded}

\subsubsection{\texorpdfstring{\texttt{break} and
\texttt{continue}}{break and continue}}\label{break-and-continue}

\begin{Shaded}
\begin{Highlighting}[]
\CommentTok{// break {-}{-}{-} exit the loop entirely}
\ControlFlowTok{for} \OperatorTok{(}\DataTypeTok{int}\NormalTok{ i }\OperatorTok{=} \DecValTok{0}\OperatorTok{;}\NormalTok{ i }\OperatorTok{\textless{}} \DecValTok{10}\OperatorTok{;}\NormalTok{ i}\OperatorTok{++)} \OperatorTok{\{}
    \ControlFlowTok{if} \OperatorTok{(}\NormalTok{i }\OperatorTok{==} \DecValTok{5}\OperatorTok{)} \ControlFlowTok{break}\OperatorTok{;}
    \BuiltInTok{System}\OperatorTok{.}\FunctionTok{out}\OperatorTok{.}\FunctionTok{println}\OperatorTok{(}\NormalTok{i}\OperatorTok{);}   \CommentTok{// Prints 0, 1, 2, 3, 4}
\OperatorTok{\}}

\CommentTok{// continue {-}{-}{-} skip the current iteration}
\ControlFlowTok{for} \OperatorTok{(}\DataTypeTok{int}\NormalTok{ i }\OperatorTok{=} \DecValTok{0}\OperatorTok{;}\NormalTok{ i }\OperatorTok{\textless{}} \DecValTok{10}\OperatorTok{;}\NormalTok{ i}\OperatorTok{++)} \OperatorTok{\{}
    \ControlFlowTok{if} \OperatorTok{(}\NormalTok{i }\OperatorTok{\%} \DecValTok{2} \OperatorTok{==} \DecValTok{0}\OperatorTok{)} \ControlFlowTok{continue}\OperatorTok{;}
    \BuiltInTok{System}\OperatorTok{.}\FunctionTok{out}\OperatorTok{.}\FunctionTok{println}\OperatorTok{(}\NormalTok{i}\OperatorTok{);}   \CommentTok{// Prints 1, 3, 5, 7, 9}
\OperatorTok{\}}
\end{Highlighting}
\end{Shaded}

\subsubsection{Skills Checklist}\label{skills-checklist-9}

\begin{itemize}
\tightlist
\item[$\square$]
  I can write a \texttt{while} loop.
\item[$\square$]
  I can write a \texttt{do-while} loop.
\item[$\square$]
  I can use \texttt{break} to exit a loop early.
\item[$\square$]
  I can use \texttt{continue} to skip an iteration.
\end{itemize}

\begin{center}\rule{0.5\linewidth}{0.5pt}\end{center}

\section{Unit 4: Data Structures --- Arrays and
ArrayLists}\label{unit-4-data-structures-arrays-and-arraylists}

\subsection{Overview}\label{overview-4}

This unit introduces Java's built-in array type and the flexible
\texttt{ArrayList} class.

\begin{center}\rule{0.5\linewidth}{0.5pt}\end{center}

\subsection{Lesson 4.1: Arrays}\label{lesson-4.1-arrays-1}

\subsubsection{Declaring and Initializing
Arrays}\label{declaring-and-initializing-arrays}

\begin{Shaded}
\begin{Highlighting}[]
\CommentTok{// Declare and allocate}
\DataTypeTok{int}\OperatorTok{[]}\NormalTok{ numbers }\OperatorTok{=} \KeywordTok{new} \DataTypeTok{int}\OperatorTok{[}\DecValTok{5}\OperatorTok{];}            \CommentTok{// [0, 0, 0, 0, 0]}

\CommentTok{// Declare and initialize with values}
\BuiltInTok{String}\OperatorTok{[]}\NormalTok{ colors }\OperatorTok{=} \OperatorTok{\{}\StringTok{"red"}\OperatorTok{,} \StringTok{"blue"}\OperatorTok{,} \StringTok{"green"}\OperatorTok{,} \StringTok{"yellow"}\OperatorTok{\};}

\CommentTok{// Access elements (zero{-}based index)}
\BuiltInTok{System}\OperatorTok{.}\FunctionTok{out}\OperatorTok{.}\FunctionTok{println}\OperatorTok{(}\NormalTok{colors}\OperatorTok{[}\DecValTok{0}\OperatorTok{]);}         \CommentTok{// "red"}
\BuiltInTok{System}\OperatorTok{.}\FunctionTok{out}\OperatorTok{.}\FunctionTok{println}\OperatorTok{(}\NormalTok{colors}\OperatorTok{[}\NormalTok{colors}\OperatorTok{.}\FunctionTok{length} \OperatorTok{{-}} \DecValTok{1}\OperatorTok{]);}  \CommentTok{// "yellow"}

\CommentTok{// Modify an element}
\NormalTok{colors}\OperatorTok{[}\DecValTok{1}\OperatorTok{]} \OperatorTok{=} \StringTok{"purple"}\OperatorTok{;}

\CommentTok{// Loop through array}
\ControlFlowTok{for} \OperatorTok{(}\DataTypeTok{int}\NormalTok{ i }\OperatorTok{=} \DecValTok{0}\OperatorTok{;}\NormalTok{ i }\OperatorTok{\textless{}}\NormalTok{ numbers}\OperatorTok{.}\FunctionTok{length}\OperatorTok{;}\NormalTok{ i}\OperatorTok{++)} \OperatorTok{\{}
\NormalTok{    numbers}\OperatorTok{[}\NormalTok{i}\OperatorTok{]} \OperatorTok{=}\NormalTok{ i }\OperatorTok{*} \DecValTok{2}\OperatorTok{;}
    \BuiltInTok{System}\OperatorTok{.}\FunctionTok{out}\OperatorTok{.}\FunctionTok{println}\OperatorTok{(}\StringTok{"numbers["} \OperatorTok{+}\NormalTok{ i }\OperatorTok{+} \StringTok{"] = "} \OperatorTok{+}\NormalTok{ numbers}\OperatorTok{[}\NormalTok{i}\OperatorTok{]);}
\OperatorTok{\}}
\end{Highlighting}
\end{Shaded}

\subsubsection{Common Array Operations}\label{common-array-operations}

\begin{Shaded}
\begin{Highlighting}[]
\KeywordTok{import} \ImportTok{java}\OperatorTok{.}\ImportTok{util}\OperatorTok{.}\ImportTok{Arrays}\OperatorTok{;}

\DataTypeTok{int}\OperatorTok{[]}\NormalTok{ data }\OperatorTok{=} \OperatorTok{\{}\DecValTok{5}\OperatorTok{,} \DecValTok{2}\OperatorTok{,} \DecValTok{9}\OperatorTok{,} \DecValTok{1}\OperatorTok{,} \DecValTok{7}\OperatorTok{,} \DecValTok{3}\OperatorTok{\};}

\CommentTok{// Sort}
\BuiltInTok{Arrays}\OperatorTok{.}\FunctionTok{sort}\OperatorTok{(}\NormalTok{data}\OperatorTok{);}
\BuiltInTok{System}\OperatorTok{.}\FunctionTok{out}\OperatorTok{.}\FunctionTok{println}\OperatorTok{(}\BuiltInTok{Arrays}\OperatorTok{.}\FunctionTok{toString}\OperatorTok{(}\NormalTok{data}\OperatorTok{));}   \CommentTok{// [1, 2, 3, 5, 7, 9]}

\CommentTok{// Search (array must be sorted first)}
\DataTypeTok{int}\NormalTok{ index }\OperatorTok{=} \BuiltInTok{Arrays}\OperatorTok{.}\FunctionTok{binarySearch}\OperatorTok{(}\NormalTok{data}\OperatorTok{,} \DecValTok{7}\OperatorTok{);}
\BuiltInTok{System}\OperatorTok{.}\FunctionTok{out}\OperatorTok{.}\FunctionTok{println}\OperatorTok{(}\StringTok{"Found 7 at index: "} \OperatorTok{+}\NormalTok{ index}\OperatorTok{);}

\CommentTok{// Copy}
\DataTypeTok{int}\OperatorTok{[]}\NormalTok{ copy }\OperatorTok{=} \BuiltInTok{Arrays}\OperatorTok{.}\FunctionTok{copyOf}\OperatorTok{(}\NormalTok{data}\OperatorTok{,}\NormalTok{ data}\OperatorTok{.}\FunctionTok{length}\OperatorTok{);}
\DataTypeTok{int}\OperatorTok{[]}\NormalTok{ partial }\OperatorTok{=} \BuiltInTok{Arrays}\OperatorTok{.}\FunctionTok{copyOfRange}\OperatorTok{(}\NormalTok{data}\OperatorTok{,} \DecValTok{1}\OperatorTok{,} \DecValTok{4}\OperatorTok{);}  \CommentTok{// [2, 3, 5]}
\end{Highlighting}
\end{Shaded}

\subsubsection{Skills Checklist}\label{skills-checklist-10}

\begin{itemize}
\tightlist
\item[$\square$]
  I can declare and initialize an array in Java.
\item[$\square$]
  I can access and modify array elements by index.
\item[$\square$]
  I can iterate over an array using a \texttt{for} loop.
\item[$\square$]
  I can use \texttt{Arrays.sort()} and \texttt{Arrays.toString()}.
\end{itemize}

\begin{center}\rule{0.5\linewidth}{0.5pt}\end{center}

\subsection{Lesson 4.2: ArrayLists}\label{lesson-4.2-arraylists-1}

\texttt{ArrayList} is a resizable array that can grow or shrink
dynamically.

\begin{Shaded}
\begin{Highlighting}[]
\KeywordTok{import} \ImportTok{java}\OperatorTok{.}\ImportTok{util}\OperatorTok{.}\ImportTok{ArrayList}\OperatorTok{;}

\BuiltInTok{ArrayList}\OperatorTok{\textless{}}\BuiltInTok{String}\OperatorTok{\textgreater{}}\NormalTok{ shoppingList }\OperatorTok{=} \KeywordTok{new} \BuiltInTok{ArrayList}\OperatorTok{\textless{}\textgreater{}();}

\CommentTok{// Add items}
\NormalTok{shoppingList}\OperatorTok{.}\FunctionTok{add}\OperatorTok{(}\StringTok{"milk"}\OperatorTok{);}
\NormalTok{shoppingList}\OperatorTok{.}\FunctionTok{add}\OperatorTok{(}\StringTok{"bread"}\OperatorTok{);}
\NormalTok{shoppingList}\OperatorTok{.}\FunctionTok{add}\OperatorTok{(}\StringTok{"eggs"}\OperatorTok{);}
\NormalTok{shoppingList}\OperatorTok{.}\FunctionTok{add}\OperatorTok{(}\DecValTok{1}\OperatorTok{,} \StringTok{"butter"}\OperatorTok{);}   \CommentTok{// Insert at index 1}

\CommentTok{// Access items}
\BuiltInTok{System}\OperatorTok{.}\FunctionTok{out}\OperatorTok{.}\FunctionTok{println}\OperatorTok{(}\NormalTok{shoppingList}\OperatorTok{.}\FunctionTok{get}\OperatorTok{(}\DecValTok{0}\OperatorTok{));}          \CommentTok{// "milk"}
\BuiltInTok{System}\OperatorTok{.}\FunctionTok{out}\OperatorTok{.}\FunctionTok{println}\OperatorTok{(}\NormalTok{shoppingList}\OperatorTok{.}\FunctionTok{size}\OperatorTok{());}           \CommentTok{// 4}

\CommentTok{// Remove items}
\NormalTok{shoppingList}\OperatorTok{.}\FunctionTok{remove}\OperatorTok{(}\StringTok{"bread"}\OperatorTok{);}                     \CommentTok{// Remove by value}
\NormalTok{shoppingList}\OperatorTok{.}\FunctionTok{remove}\OperatorTok{(}\DecValTok{0}\OperatorTok{);}                           \CommentTok{// Remove by index}

\CommentTok{// Check membership}
\DataTypeTok{boolean}\NormalTok{ hasMilk }\OperatorTok{=}\NormalTok{ shoppingList}\OperatorTok{.}\FunctionTok{contains}\OperatorTok{(}\StringTok{"milk"}\OperatorTok{);}

\CommentTok{// Loop through}
\ControlFlowTok{for} \OperatorTok{(}\BuiltInTok{String}\NormalTok{ item }\OperatorTok{:}\NormalTok{ shoppingList}\OperatorTok{)} \OperatorTok{\{}
    \BuiltInTok{System}\OperatorTok{.}\FunctionTok{out}\OperatorTok{.}\FunctionTok{println}\OperatorTok{(}\StringTok{"{-} "} \OperatorTok{+}\NormalTok{ item}\OperatorTok{);}
\OperatorTok{\}}
\end{Highlighting}
\end{Shaded}

\subsubsection{ArrayList vs.~Array}\label{arraylist-vs.-array}

{\def\LTcaptype{none} % do not increment counter
\begin{longtable}[]{@{}lll@{}}
\toprule\noalign{}
Feature & \texttt{Array} & \texttt{ArrayList} \\
\midrule\noalign{}
\endhead
\bottomrule\noalign{}
\endlastfoot
Size & Fixed & Dynamic \\
Type & Primitive or Object & Objects only \\
Syntax & \texttt{int{[}{]}} &
\texttt{ArrayList\textless{}Integer\textgreater{}} \\
\texttt{.length} / \texttt{.size()} & \texttt{.length} &
\texttt{.size()} \\
Performance & Slightly faster & More flexible \\
\end{longtable}
}

\subsubsection{Skills Checklist}\label{skills-checklist-11}

\begin{itemize}
\tightlist
\item[$\square$]
  I can create an \texttt{ArrayList} with a generic type.
\item[$\square$]
  I can add, access, remove, and search items in an \texttt{ArrayList}.
\item[$\square$]
  I can explain when to use an array vs.~an \texttt{ArrayList}.
\end{itemize}

\begin{center}\rule{0.5\linewidth}{0.5pt}\end{center}

\subsection{Lesson 4.3: Multidimensional
Arrays}\label{lesson-4.3-multidimensional-arrays}

\begin{Shaded}
\begin{Highlighting}[]
\CommentTok{// 2D array (matrix)}
\DataTypeTok{int}\OperatorTok{[][]}\NormalTok{ matrix }\OperatorTok{=} \OperatorTok{\{}
    \OperatorTok{\{}\DecValTok{1}\OperatorTok{,} \DecValTok{2}\OperatorTok{,} \DecValTok{3}\OperatorTok{\},}
    \OperatorTok{\{}\DecValTok{4}\OperatorTok{,} \DecValTok{5}\OperatorTok{,} \DecValTok{6}\OperatorTok{\},}
    \OperatorTok{\{}\DecValTok{7}\OperatorTok{,} \DecValTok{8}\OperatorTok{,} \DecValTok{9}\OperatorTok{\}}
\OperatorTok{\};}

\CommentTok{// Access element at row 1, column 2}
\BuiltInTok{System}\OperatorTok{.}\FunctionTok{out}\OperatorTok{.}\FunctionTok{println}\OperatorTok{(}\NormalTok{matrix}\OperatorTok{[}\DecValTok{1}\OperatorTok{][}\DecValTok{2}\OperatorTok{]);}   \CommentTok{// 6}

\CommentTok{// Loop through 2D array}
\ControlFlowTok{for} \OperatorTok{(}\DataTypeTok{int}\NormalTok{ row }\OperatorTok{=} \DecValTok{0}\OperatorTok{;}\NormalTok{ row }\OperatorTok{\textless{}}\NormalTok{ matrix}\OperatorTok{.}\FunctionTok{length}\OperatorTok{;}\NormalTok{ row}\OperatorTok{++)} \OperatorTok{\{}
    \ControlFlowTok{for} \OperatorTok{(}\DataTypeTok{int}\NormalTok{ col }\OperatorTok{=} \DecValTok{0}\OperatorTok{;}\NormalTok{ col }\OperatorTok{\textless{}}\NormalTok{ matrix}\OperatorTok{[}\NormalTok{row}\OperatorTok{].}\FunctionTok{length}\OperatorTok{;}\NormalTok{ col}\OperatorTok{++)} \OperatorTok{\{}
        \BuiltInTok{System}\OperatorTok{.}\FunctionTok{out}\OperatorTok{.}\FunctionTok{print}\OperatorTok{(}\NormalTok{matrix}\OperatorTok{[}\NormalTok{row}\OperatorTok{][}\NormalTok{col}\OperatorTok{]} \OperatorTok{+} \StringTok{" "}\OperatorTok{);}
    \OperatorTok{\}}
    \BuiltInTok{System}\OperatorTok{.}\FunctionTok{out}\OperatorTok{.}\FunctionTok{println}\OperatorTok{();}
\OperatorTok{\}}
\end{Highlighting}
\end{Shaded}

\subsubsection{Skills Checklist}\label{skills-checklist-12}

\begin{itemize}
\tightlist
\item[$\square$]
  I can declare and initialize a 2D array.
\item[$\square$]
  I can access elements using row and column indices.
\item[$\square$]
  I can iterate through a 2D array using nested loops.
\end{itemize}

\newpage

\section{Unit 5: Strings and String
Processing}\label{unit-5-strings-and-string-processing-1}

\subsection{Overview}\label{overview-5}

Java's \texttt{String} class is rich with methods for text manipulation.

\begin{center}\rule{0.5\linewidth}{0.5pt}\end{center}

\subsection{Lesson 5.1: Working with
Strings}\label{lesson-5.1-working-with-strings}

\subsubsection{String Basics}\label{string-basics}

\begin{Shaded}
\begin{Highlighting}[]
\BuiltInTok{String}\NormalTok{ name }\OperatorTok{=} \StringTok{"Alice Smith"}\OperatorTok{;}

\CommentTok{// Length}
\BuiltInTok{System}\OperatorTok{.}\FunctionTok{out}\OperatorTok{.}\FunctionTok{println}\OperatorTok{(}\NormalTok{name}\OperatorTok{.}\FunctionTok{length}\OperatorTok{());}           \CommentTok{// 11}

\CommentTok{// Accessing characters}
\BuiltInTok{System}\OperatorTok{.}\FunctionTok{out}\OperatorTok{.}\FunctionTok{println}\OperatorTok{(}\NormalTok{name}\OperatorTok{.}\FunctionTok{charAt}\OperatorTok{(}\DecValTok{0}\OperatorTok{));}          \CommentTok{// \textquotesingle{}A\textquotesingle{}}
\BuiltInTok{System}\OperatorTok{.}\FunctionTok{out}\OperatorTok{.}\FunctionTok{println}\OperatorTok{(}\NormalTok{name}\OperatorTok{.}\FunctionTok{charAt}\OperatorTok{(}\NormalTok{name}\OperatorTok{.}\FunctionTok{length}\OperatorTok{()} \OperatorTok{{-}} \DecValTok{1}\OperatorTok{));}  \CommentTok{// \textquotesingle{}h\textquotesingle{}}

\CommentTok{// Substrings}
\BuiltInTok{System}\OperatorTok{.}\FunctionTok{out}\OperatorTok{.}\FunctionTok{println}\OperatorTok{(}\NormalTok{name}\OperatorTok{.}\FunctionTok{substring}\OperatorTok{(}\DecValTok{6}\OperatorTok{));}       \CommentTok{// "Smith"}
\BuiltInTok{System}\OperatorTok{.}\FunctionTok{out}\OperatorTok{.}\FunctionTok{println}\OperatorTok{(}\NormalTok{name}\OperatorTok{.}\FunctionTok{substring}\OperatorTok{(}\DecValTok{0}\OperatorTok{,} \DecValTok{5}\OperatorTok{));}    \CommentTok{// "Alice"}

\CommentTok{// Case conversion}
\BuiltInTok{System}\OperatorTok{.}\FunctionTok{out}\OperatorTok{.}\FunctionTok{println}\OperatorTok{(}\NormalTok{name}\OperatorTok{.}\FunctionTok{toUpperCase}\OperatorTok{());}      \CommentTok{// "ALICE SMITH"}
\BuiltInTok{System}\OperatorTok{.}\FunctionTok{out}\OperatorTok{.}\FunctionTok{println}\OperatorTok{(}\NormalTok{name}\OperatorTok{.}\FunctionTok{toLowerCase}\OperatorTok{());}      \CommentTok{// "alice smith"}

\CommentTok{// Trimming whitespace}
\BuiltInTok{String}\NormalTok{ padded }\OperatorTok{=} \StringTok{"  hello  "}\OperatorTok{;}
\BuiltInTok{System}\OperatorTok{.}\FunctionTok{out}\OperatorTok{.}\FunctionTok{println}\OperatorTok{(}\NormalTok{padded}\OperatorTok{.}\FunctionTok{trim}\OperatorTok{());}           \CommentTok{// "hello"}
\BuiltInTok{System}\OperatorTok{.}\FunctionTok{out}\OperatorTok{.}\FunctionTok{println}\OperatorTok{(}\NormalTok{padded}\OperatorTok{.}\FunctionTok{strip}\OperatorTok{());}          \CommentTok{// "hello" (better Unicode support)}

\CommentTok{// Searching}
\BuiltInTok{System}\OperatorTok{.}\FunctionTok{out}\OperatorTok{.}\FunctionTok{println}\OperatorTok{(}\NormalTok{name}\OperatorTok{.}\FunctionTok{indexOf}\OperatorTok{(}\StringTok{"Smith"}\OperatorTok{));}   \CommentTok{// 6}
\BuiltInTok{System}\OperatorTok{.}\FunctionTok{out}\OperatorTok{.}\FunctionTok{println}\OperatorTok{(}\NormalTok{name}\OperatorTok{.}\FunctionTok{contains}\OperatorTok{(}\StringTok{"Alice"}\OperatorTok{));}  \CommentTok{// true}
\BuiltInTok{System}\OperatorTok{.}\FunctionTok{out}\OperatorTok{.}\FunctionTok{println}\OperatorTok{(}\NormalTok{name}\OperatorTok{.}\FunctionTok{startsWith}\OperatorTok{(}\StringTok{"Ali"}\OperatorTok{));}  \CommentTok{// true}
\BuiltInTok{System}\OperatorTok{.}\FunctionTok{out}\OperatorTok{.}\FunctionTok{println}\OperatorTok{(}\NormalTok{name}\OperatorTok{.}\FunctionTok{endsWith}\OperatorTok{(}\StringTok{"ith"}\OperatorTok{));}    \CommentTok{// true}

\CommentTok{// Replacing}
\BuiltInTok{String}\NormalTok{ result }\OperatorTok{=}\NormalTok{ name}\OperatorTok{.}\FunctionTok{replace}\OperatorTok{(}\StringTok{"Alice"}\OperatorTok{,} \StringTok{"Bob"}\OperatorTok{);}   \CommentTok{// "Bob Smith"}
\BuiltInTok{String}\NormalTok{ clean }\OperatorTok{=} \StringTok{"a,b,c"}\OperatorTok{.}\FunctionTok{replace}\OperatorTok{(}\StringTok{","}\OperatorTok{,} \StringTok{" "}\OperatorTok{);}       \CommentTok{// "a b c"}

\CommentTok{// Splitting}
\BuiltInTok{String}\NormalTok{ sentence }\OperatorTok{=} \StringTok{"the quick brown fox"}\OperatorTok{;}
\BuiltInTok{String}\OperatorTok{[]}\NormalTok{ words }\OperatorTok{=}\NormalTok{ sentence}\OperatorTok{.}\FunctionTok{split}\OperatorTok{(}\StringTok{" "}\OperatorTok{);}
\ControlFlowTok{for} \OperatorTok{(}\BuiltInTok{String}\NormalTok{ word }\OperatorTok{:}\NormalTok{ words}\OperatorTok{)} \OperatorTok{\{}
    \BuiltInTok{System}\OperatorTok{.}\FunctionTok{out}\OperatorTok{.}\FunctionTok{println}\OperatorTok{(}\NormalTok{word}\OperatorTok{);}
\OperatorTok{\}}

\CommentTok{// Checking equality}
\BuiltInTok{String}\NormalTok{ s1 }\OperatorTok{=} \StringTok{"hello"}\OperatorTok{;}
\BuiltInTok{String}\NormalTok{ s2 }\OperatorTok{=} \StringTok{"hello"}\OperatorTok{;}
\BuiltInTok{System}\OperatorTok{.}\FunctionTok{out}\OperatorTok{.}\FunctionTok{println}\OperatorTok{(}\NormalTok{s1}\OperatorTok{.}\FunctionTok{equals}\OperatorTok{(}\NormalTok{s2}\OperatorTok{));}           \CommentTok{// true (correct)}
\BuiltInTok{System}\OperatorTok{.}\FunctionTok{out}\OperatorTok{.}\FunctionTok{println}\OperatorTok{(}\NormalTok{s1}\OperatorTok{.}\FunctionTok{equalsIgnoreCase}\OperatorTok{(}\StringTok{"HELLO"}\OperatorTok{));} \CommentTok{// true}

\CommentTok{// Converting other types to String}
\BuiltInTok{String}\NormalTok{ fromInt }\OperatorTok{=} \BuiltInTok{String}\OperatorTok{.}\FunctionTok{valueOf}\OperatorTok{(}\DecValTok{42}\OperatorTok{);}
\BuiltInTok{String}\NormalTok{ fromDouble }\OperatorTok{=} \BuiltInTok{String}\OperatorTok{.}\FunctionTok{valueOf}\OperatorTok{(}\FloatTok{3.14}\OperatorTok{);}
\BuiltInTok{String}\NormalTok{ formatted }\OperatorTok{=} \BuiltInTok{String}\OperatorTok{.}\FunctionTok{format}\OperatorTok{(}\StringTok{"Name: }\SpecialCharTok{\%{-}10s}\StringTok{ Age: }\SpecialCharTok{\%3d}\StringTok{"}\OperatorTok{,} \StringTok{"Alice"}\OperatorTok{,} \DecValTok{15}\OperatorTok{);}
\end{Highlighting}
\end{Shaded}

\subsubsection{\texorpdfstring{\texttt{StringBuilder} for Efficient
Concatenation}{StringBuilder for Efficient Concatenation}}\label{stringbuilder-for-efficient-concatenation}

Use \texttt{StringBuilder} when building strings in a loop (more
efficient than \texttt{+}):

\begin{Shaded}
\begin{Highlighting}[]
\BuiltInTok{StringBuilder}\NormalTok{ sb }\OperatorTok{=} \KeywordTok{new} \BuiltInTok{StringBuilder}\OperatorTok{();}
\ControlFlowTok{for} \OperatorTok{(}\DataTypeTok{int}\NormalTok{ i }\OperatorTok{=} \DecValTok{1}\OperatorTok{;}\NormalTok{ i }\OperatorTok{\textless{}=} \DecValTok{5}\OperatorTok{;}\NormalTok{ i}\OperatorTok{++)} \OperatorTok{\{}
\NormalTok{    sb}\OperatorTok{.}\FunctionTok{append}\OperatorTok{(}\StringTok{"Item "}\OperatorTok{).}\FunctionTok{append}\OperatorTok{(}\NormalTok{i}\OperatorTok{).}\FunctionTok{append}\OperatorTok{(}\StringTok{"}\SpecialCharTok{\textbackslash{}n}\StringTok{"}\OperatorTok{);}
\OperatorTok{\}}
\BuiltInTok{System}\OperatorTok{.}\FunctionTok{out}\OperatorTok{.}\FunctionTok{println}\OperatorTok{(}\NormalTok{sb}\OperatorTok{.}\FunctionTok{toString}\OperatorTok{());}
\end{Highlighting}
\end{Shaded}

\subsubsection{Examples}\label{examples-5}

\textbf{Example 5.1.1: String Analysis}

\begin{Shaded}
\begin{Highlighting}[]
\KeywordTok{import} \ImportTok{java}\OperatorTok{.}\ImportTok{util}\OperatorTok{.}\ImportTok{Scanner}\OperatorTok{;}

\KeywordTok{public} \KeywordTok{class}\NormalTok{ StringAnalysis }\OperatorTok{\{}
    \KeywordTok{public} \DataTypeTok{static} \DataTypeTok{void} \FunctionTok{main}\OperatorTok{(}\BuiltInTok{String}\OperatorTok{[]}\NormalTok{ args}\OperatorTok{)} \OperatorTok{\{}
        \BuiltInTok{Scanner}\NormalTok{ scanner }\OperatorTok{=} \KeywordTok{new} \BuiltInTok{Scanner}\OperatorTok{(}\BuiltInTok{System}\OperatorTok{.}\FunctionTok{in}\OperatorTok{);}
        \BuiltInTok{System}\OperatorTok{.}\FunctionTok{out}\OperatorTok{.}\FunctionTok{print}\OperatorTok{(}\StringTok{"Enter a sentence: "}\OperatorTok{);}
        \BuiltInTok{String}\NormalTok{ sentence }\OperatorTok{=}\NormalTok{ scanner}\OperatorTok{.}\FunctionTok{nextLine}\OperatorTok{();}

        \BuiltInTok{String}\OperatorTok{[]}\NormalTok{ words }\OperatorTok{=}\NormalTok{ sentence}\OperatorTok{.}\FunctionTok{trim}\OperatorTok{().}\FunctionTok{split}\OperatorTok{(}\StringTok{"}\SpecialCharTok{\textbackslash{}\textbackslash{}}\StringTok{s+"}\OperatorTok{);}

        \BuiltInTok{System}\OperatorTok{.}\FunctionTok{out}\OperatorTok{.}\FunctionTok{println}\OperatorTok{(}\StringTok{"Original: "} \OperatorTok{+}\NormalTok{ sentence}\OperatorTok{);}
        \BuiltInTok{System}\OperatorTok{.}\FunctionTok{out}\OperatorTok{.}\FunctionTok{println}\OperatorTok{(}\StringTok{"Character count: "} \OperatorTok{+}\NormalTok{ sentence}\OperatorTok{.}\FunctionTok{length}\OperatorTok{());}
        \BuiltInTok{System}\OperatorTok{.}\FunctionTok{out}\OperatorTok{.}\FunctionTok{println}\OperatorTok{(}\StringTok{"Word count: "} \OperatorTok{+}\NormalTok{ words}\OperatorTok{.}\FunctionTok{length}\OperatorTok{);}
        \BuiltInTok{System}\OperatorTok{.}\FunctionTok{out}\OperatorTok{.}\FunctionTok{println}\OperatorTok{(}\StringTok{"Uppercase: "} \OperatorTok{+}\NormalTok{ sentence}\OperatorTok{.}\FunctionTok{toUpperCase}\OperatorTok{());}
        \BuiltInTok{System}\OperatorTok{.}\FunctionTok{out}\OperatorTok{.}\FunctionTok{println}\OperatorTok{(}\StringTok{"Lowercase: "} \OperatorTok{+}\NormalTok{ sentence}\OperatorTok{.}\FunctionTok{toLowerCase}\OperatorTok{());}

\NormalTok{        scanner}\OperatorTok{.}\FunctionTok{close}\OperatorTok{();}
    \OperatorTok{\}}
\OperatorTok{\}}
\end{Highlighting}
\end{Shaded}

\textbf{Example 5.1.2: Word Frequency Counter}

\begin{Shaded}
\begin{Highlighting}[]
\KeywordTok{import} \ImportTok{java}\OperatorTok{.}\ImportTok{util}\OperatorTok{.}\ImportTok{HashMap}\OperatorTok{;}
\KeywordTok{import} \ImportTok{java}\OperatorTok{.}\ImportTok{util}\OperatorTok{.}\ImportTok{Map}\OperatorTok{;}

\KeywordTok{public} \KeywordTok{class}\NormalTok{ WordFrequency }\OperatorTok{\{}
    \KeywordTok{public} \DataTypeTok{static} \DataTypeTok{void} \FunctionTok{main}\OperatorTok{(}\BuiltInTok{String}\OperatorTok{[]}\NormalTok{ args}\OperatorTok{)} \OperatorTok{\{}
        \BuiltInTok{String}\NormalTok{ text }\OperatorTok{=} \StringTok{"the quick brown fox jumps over the lazy dog the quick brown fox"}\OperatorTok{;}
        \BuiltInTok{String}\OperatorTok{[]}\NormalTok{ words }\OperatorTok{=}\NormalTok{ text}\OperatorTok{.}\FunctionTok{toLowerCase}\OperatorTok{().}\FunctionTok{split}\OperatorTok{(}\StringTok{"}\SpecialCharTok{\textbackslash{}\textbackslash{}}\StringTok{s+"}\OperatorTok{);}

        \BuiltInTok{Map}\OperatorTok{\textless{}}\BuiltInTok{String}\OperatorTok{,} \BuiltInTok{Integer}\OperatorTok{\textgreater{}}\NormalTok{ frequency }\OperatorTok{=} \KeywordTok{new} \BuiltInTok{HashMap}\OperatorTok{\textless{}\textgreater{}();}
        \ControlFlowTok{for} \OperatorTok{(}\BuiltInTok{String}\NormalTok{ word }\OperatorTok{:}\NormalTok{ words}\OperatorTok{)} \OperatorTok{\{}
\NormalTok{            frequency}\OperatorTok{.}\FunctionTok{put}\OperatorTok{(}\NormalTok{word}\OperatorTok{,}\NormalTok{ frequency}\OperatorTok{.}\FunctionTok{getOrDefault}\OperatorTok{(}\NormalTok{word}\OperatorTok{,} \DecValTok{0}\OperatorTok{)} \OperatorTok{+} \DecValTok{1}\OperatorTok{);}
        \OperatorTok{\}}

        \BuiltInTok{System}\OperatorTok{.}\FunctionTok{out}\OperatorTok{.}\FunctionTok{println}\OperatorTok{(}\StringTok{"Word Frequencies:"}\OperatorTok{);}
        \ControlFlowTok{for} \OperatorTok{(}\BuiltInTok{Map}\OperatorTok{.}\FunctionTok{Entry}\OperatorTok{\textless{}}\BuiltInTok{String}\OperatorTok{,} \BuiltInTok{Integer}\OperatorTok{\textgreater{}}\NormalTok{ entry }\OperatorTok{:}\NormalTok{ frequency}\OperatorTok{.}\FunctionTok{entrySet}\OperatorTok{())} \OperatorTok{\{}
            \BuiltInTok{System}\OperatorTok{.}\FunctionTok{out}\OperatorTok{.}\FunctionTok{println}\OperatorTok{(}\StringTok{"\textquotesingle{}"} \OperatorTok{+}\NormalTok{ entry}\OperatorTok{.}\FunctionTok{getKey}\OperatorTok{()} \OperatorTok{+} \StringTok{"\textquotesingle{}: "} \OperatorTok{+}\NormalTok{ entry}\OperatorTok{.}\FunctionTok{getValue}\OperatorTok{());}
        \OperatorTok{\}}
    \OperatorTok{\}}
\OperatorTok{\}}
\end{Highlighting}
\end{Shaded}

\subsubsection{Skills Checklist}\label{skills-checklist-13}

\begin{itemize}
\tightlist
\item[$\square$]
  I can use common \texttt{String} methods in Java.
\item[$\square$]
  I can compare strings using \texttt{.equals()}.
\item[$\square$]
  I can split a string into an array of words.
\item[$\square$]
  I can use \texttt{StringBuilder} for efficient string building.
\end{itemize}

\begin{center}\rule{0.5\linewidth}{0.5pt}\end{center}

\section{Unit 6: Methods (Functions)}\label{unit-6-methods-functions}

\subsection{Overview}\label{overview-6}

Java organizes reusable code into \textbf{methods} (the Java equivalent
of Python functions).

\begin{center}\rule{0.5\linewidth}{0.5pt}\end{center}

\subsection{Lesson 6.1: Defining and Calling
Methods}\label{lesson-6.1-defining-and-calling-methods}

\subsubsection{Method Syntax}\label{method-syntax}

\begin{Shaded}
\begin{Highlighting}[]
\NormalTok{accessModifier returnType }\FunctionTok{methodName}\OperatorTok{(}\NormalTok{paramType paramName}\OperatorTok{,} \KeywordTok{...}\OperatorTok{)} \OperatorTok{\{}
    \CommentTok{// method body}
    \ControlFlowTok{return}\NormalTok{ value}\OperatorTok{;}  \CommentTok{// if returnType is not void}
\OperatorTok{\}}
\end{Highlighting}
\end{Shaded}

\subsubsection{Examples}\label{examples-6}

\textbf{Example 6.1.1: Simple Greeting Method}

\begin{Shaded}
\begin{Highlighting}[]
\KeywordTok{public} \KeywordTok{class}\NormalTok{ GreetingDemo }\OperatorTok{\{}

    \KeywordTok{public} \DataTypeTok{static} \DataTypeTok{void} \FunctionTok{greetUser}\OperatorTok{()} \OperatorTok{\{}
        \BuiltInTok{System}\OperatorTok{.}\FunctionTok{out}\OperatorTok{.}\FunctionTok{println}\OperatorTok{(}\StringTok{"Hello, Java programmer!"}\OperatorTok{);}
    \OperatorTok{\}}

    \KeywordTok{public} \DataTypeTok{static} \DataTypeTok{void} \FunctionTok{greetByName}\OperatorTok{(}\BuiltInTok{String}\NormalTok{ name}\OperatorTok{)} \OperatorTok{\{}
        \BuiltInTok{System}\OperatorTok{.}\FunctionTok{out}\OperatorTok{.}\FunctionTok{println}\OperatorTok{(}\StringTok{"Hello, "} \OperatorTok{+}\NormalTok{ name }\OperatorTok{+} \StringTok{"!"}\OperatorTok{);}
    \OperatorTok{\}}

    \KeywordTok{public} \DataTypeTok{static} \DataTypeTok{void} \FunctionTok{main}\OperatorTok{(}\BuiltInTok{String}\OperatorTok{[]}\NormalTok{ args}\OperatorTok{)} \OperatorTok{\{}
        \FunctionTok{greetUser}\OperatorTok{();}
        \FunctionTok{greetByName}\OperatorTok{(}\StringTok{"Alice"}\OperatorTok{);}
        \FunctionTok{greetByName}\OperatorTok{(}\StringTok{"Bob"}\OperatorTok{);}
    \OperatorTok{\}}
\OperatorTok{\}}
\end{Highlighting}
\end{Shaded}

\textbf{Example 6.1.2: Return Values}

\begin{Shaded}
\begin{Highlighting}[]
\KeywordTok{public} \KeywordTok{class}\NormalTok{ Calculator }\OperatorTok{\{}

    \KeywordTok{public} \DataTypeTok{static} \DataTypeTok{int} \FunctionTok{addNumbers}\OperatorTok{(}\DataTypeTok{int}\NormalTok{ a}\OperatorTok{,} \DataTypeTok{int}\NormalTok{ b}\OperatorTok{)} \OperatorTok{\{}
        \ControlFlowTok{return}\NormalTok{ a }\OperatorTok{+}\NormalTok{ b}\OperatorTok{;}
    \OperatorTok{\}}

    \KeywordTok{public} \DataTypeTok{static} \DataTypeTok{double} \FunctionTok{calculateCircleArea}\OperatorTok{(}\DataTypeTok{double}\NormalTok{ radius}\OperatorTok{)} \OperatorTok{\{}
        \ControlFlowTok{return} \BuiltInTok{Math}\OperatorTok{.}\FunctionTok{PI} \OperatorTok{*}\NormalTok{ radius }\OperatorTok{*}\NormalTok{ radius}\OperatorTok{;}
    \OperatorTok{\}}

    \KeywordTok{public} \DataTypeTok{static} \BuiltInTok{String} \FunctionTok{getGrade}\OperatorTok{(}\DataTypeTok{int}\NormalTok{ score}\OperatorTok{)} \OperatorTok{\{}
        \ControlFlowTok{if} \OperatorTok{(}\NormalTok{score }\OperatorTok{\textgreater{}=} \DecValTok{90}\OperatorTok{)} \ControlFlowTok{return} \StringTok{"A"}\OperatorTok{;}
        \ControlFlowTok{else} \ControlFlowTok{if} \OperatorTok{(}\NormalTok{score }\OperatorTok{\textgreater{}=} \DecValTok{80}\OperatorTok{)} \ControlFlowTok{return} \StringTok{"B"}\OperatorTok{;}
        \ControlFlowTok{else} \ControlFlowTok{if} \OperatorTok{(}\NormalTok{score }\OperatorTok{\textgreater{}=} \DecValTok{70}\OperatorTok{)} \ControlFlowTok{return} \StringTok{"C"}\OperatorTok{;}
        \ControlFlowTok{else} \ControlFlowTok{if} \OperatorTok{(}\NormalTok{score }\OperatorTok{\textgreater{}=} \DecValTok{60}\OperatorTok{)} \ControlFlowTok{return} \StringTok{"D"}\OperatorTok{;}
        \ControlFlowTok{else} \ControlFlowTok{return} \StringTok{"F"}\OperatorTok{;}
    \OperatorTok{\}}

    \KeywordTok{public} \DataTypeTok{static} \DataTypeTok{void} \FunctionTok{main}\OperatorTok{(}\BuiltInTok{String}\OperatorTok{[]}\NormalTok{ args}\OperatorTok{)} \OperatorTok{\{}
        \DataTypeTok{int}\NormalTok{ sum }\OperatorTok{=} \FunctionTok{addNumbers}\OperatorTok{(}\DecValTok{5}\OperatorTok{,} \DecValTok{3}\OperatorTok{);}
        \BuiltInTok{System}\OperatorTok{.}\FunctionTok{out}\OperatorTok{.}\FunctionTok{println}\OperatorTok{(}\StringTok{"Sum: "} \OperatorTok{+}\NormalTok{ sum}\OperatorTok{);}

        \DataTypeTok{double}\NormalTok{ area }\OperatorTok{=} \FunctionTok{calculateCircleArea}\OperatorTok{(}\FloatTok{5.0}\OperatorTok{);}
        \BuiltInTok{System}\OperatorTok{.}\FunctionTok{out}\OperatorTok{.}\FunctionTok{printf}\OperatorTok{(}\StringTok{"Circle area: }\SpecialCharTok{\%.2f\%n}\StringTok{"}\OperatorTok{,}\NormalTok{ area}\OperatorTok{);}

        \BuiltInTok{System}\OperatorTok{.}\FunctionTok{out}\OperatorTok{.}\FunctionTok{println}\OperatorTok{(}\StringTok{"Grade for 85: "} \OperatorTok{+} \FunctionTok{getGrade}\OperatorTok{(}\DecValTok{85}\OperatorTok{));}
    \OperatorTok{\}}
\OperatorTok{\}}
\end{Highlighting}
\end{Shaded}

\subsubsection{Method Overloading}\label{method-overloading}

Java allows multiple methods with the same name but different
parameters:

\begin{Shaded}
\begin{Highlighting}[]
\KeywordTok{public} \DataTypeTok{static} \DataTypeTok{int} \FunctionTok{add}\OperatorTok{(}\DataTypeTok{int}\NormalTok{ a}\OperatorTok{,} \DataTypeTok{int}\NormalTok{ b}\OperatorTok{)} \OperatorTok{\{} \ControlFlowTok{return}\NormalTok{ a }\OperatorTok{+}\NormalTok{ b}\OperatorTok{;} \OperatorTok{\}}
\KeywordTok{public} \DataTypeTok{static} \DataTypeTok{double} \FunctionTok{add}\OperatorTok{(}\DataTypeTok{double}\NormalTok{ a}\OperatorTok{,} \DataTypeTok{double}\NormalTok{ b}\OperatorTok{)} \OperatorTok{\{} \ControlFlowTok{return}\NormalTok{ a }\OperatorTok{+}\NormalTok{ b}\OperatorTok{;} \OperatorTok{\}}
\KeywordTok{public} \DataTypeTok{static} \DataTypeTok{int} \FunctionTok{add}\OperatorTok{(}\DataTypeTok{int}\NormalTok{ a}\OperatorTok{,} \DataTypeTok{int}\NormalTok{ b}\OperatorTok{,} \DataTypeTok{int}\NormalTok{ c}\OperatorTok{)} \OperatorTok{\{} \ControlFlowTok{return}\NormalTok{ a }\OperatorTok{+}\NormalTok{ b }\OperatorTok{+}\NormalTok{ c}\OperatorTok{;} \OperatorTok{\}}
\end{Highlighting}
\end{Shaded}

\subsubsection{Passing Arrays to
Methods}\label{passing-arrays-to-methods}

\begin{Shaded}
\begin{Highlighting}[]
\KeywordTok{public} \DataTypeTok{static} \DataTypeTok{int} \FunctionTok{sum}\OperatorTok{(}\DataTypeTok{int}\OperatorTok{[]}\NormalTok{ numbers}\OperatorTok{)} \OperatorTok{\{}
    \DataTypeTok{int}\NormalTok{ total }\OperatorTok{=} \DecValTok{0}\OperatorTok{;}
    \ControlFlowTok{for} \OperatorTok{(}\DataTypeTok{int}\NormalTok{ num }\OperatorTok{:}\NormalTok{ numbers}\OperatorTok{)} \OperatorTok{\{}
\NormalTok{        total }\OperatorTok{+=}\NormalTok{ num}\OperatorTok{;}
    \OperatorTok{\}}
    \ControlFlowTok{return}\NormalTok{ total}\OperatorTok{;}
\OperatorTok{\}}

\CommentTok{// Calling:}
\DataTypeTok{int}\OperatorTok{[]}\NormalTok{ data }\OperatorTok{=} \OperatorTok{\{}\DecValTok{1}\OperatorTok{,} \DecValTok{2}\OperatorTok{,} \DecValTok{3}\OperatorTok{,} \DecValTok{4}\OperatorTok{,} \DecValTok{5}\OperatorTok{\};}
\BuiltInTok{System}\OperatorTok{.}\FunctionTok{out}\OperatorTok{.}\FunctionTok{println}\OperatorTok{(}\StringTok{"Sum: "} \OperatorTok{+} \FunctionTok{sum}\OperatorTok{(}\NormalTok{data}\OperatorTok{));}
\end{Highlighting}
\end{Shaded}

\subsubsection{Varargs (Variable
Arguments)}\label{varargs-variable-arguments}

\begin{Shaded}
\begin{Highlighting}[]
\KeywordTok{public} \DataTypeTok{static} \DataTypeTok{int} \FunctionTok{sumAll}\OperatorTok{(}\DataTypeTok{int}\KeywordTok{...}\NormalTok{ numbers}\OperatorTok{)} \OperatorTok{\{}
    \DataTypeTok{int}\NormalTok{ total }\OperatorTok{=} \DecValTok{0}\OperatorTok{;}
    \ControlFlowTok{for} \OperatorTok{(}\DataTypeTok{int}\NormalTok{ num }\OperatorTok{:}\NormalTok{ numbers}\OperatorTok{)} \OperatorTok{\{}
\NormalTok{        total }\OperatorTok{+=}\NormalTok{ num}\OperatorTok{;}
    \OperatorTok{\}}
    \ControlFlowTok{return}\NormalTok{ total}\OperatorTok{;}
\OperatorTok{\}}

\CommentTok{// Calling:}
\BuiltInTok{System}\OperatorTok{.}\FunctionTok{out}\OperatorTok{.}\FunctionTok{println}\OperatorTok{(}\FunctionTok{sumAll}\OperatorTok{(}\DecValTok{1}\OperatorTok{,} \DecValTok{2}\OperatorTok{,} \DecValTok{3}\OperatorTok{));}          \CommentTok{// 6}
\BuiltInTok{System}\OperatorTok{.}\FunctionTok{out}\OperatorTok{.}\FunctionTok{println}\OperatorTok{(}\FunctionTok{sumAll}\OperatorTok{(}\DecValTok{10}\OperatorTok{,} \DecValTok{20}\OperatorTok{,} \DecValTok{30}\OperatorTok{,} \DecValTok{40}\OperatorTok{));}    \CommentTok{// 100}
\end{Highlighting}
\end{Shaded}

\subsubsection{Skills Checklist}\label{skills-checklist-14}

\begin{itemize}
\tightlist
\item[$\square$]
  I can define a method with a return type and parameters.
\item[$\square$]
  I can call a method and use its return value.
\item[$\square$]
  I understand the difference between \texttt{void} and
  non-\texttt{void} methods.
\item[$\square$]
  I can overload methods with different parameter lists.
\item[$\square$]
  I can pass arrays to methods.
\end{itemize}

\begin{center}\rule{0.5\linewidth}{0.5pt}\end{center}

\section{Unit 7: Classes and Object-Oriented
Programming}\label{unit-7-classes-and-object-oriented-programming}

\subsection{Overview}\label{overview-7}

Java is fundamentally object-oriented. This unit covers defining and
using classes.

\begin{center}\rule{0.5\linewidth}{0.5pt}\end{center}

\subsection{Lesson 7.1: Why Learn About
Classes?}\label{lesson-7.1-why-learn-about-classes}

\emph{(See Python Unit 7.1 for conceptual content --- OOP concepts are
language-agnostic.)}

Java is a \textbf{class-based} language --- every Java program is a
class. Understanding OOP is essential for Java.

\begin{center}\rule{0.5\linewidth}{0.5pt}\end{center}

\subsection{Lesson 7.2: Creating and Using a
Class}\label{lesson-7.2-creating-and-using-a-class}

\subsubsection{Class Syntax}\label{class-syntax}

\begin{Shaded}
\begin{Highlighting}[]
\KeywordTok{public} \KeywordTok{class}\NormalTok{ Dog }\OperatorTok{\{}
    \CommentTok{// Instance variables (fields)}
    \KeywordTok{private} \BuiltInTok{String}\NormalTok{ name}\OperatorTok{;}
    \KeywordTok{private} \DataTypeTok{int}\NormalTok{ age}\OperatorTok{;}

    \CommentTok{// Constructor}
    \KeywordTok{public} \FunctionTok{Dog}\OperatorTok{(}\BuiltInTok{String}\NormalTok{ name}\OperatorTok{,} \DataTypeTok{int}\NormalTok{ age}\OperatorTok{)} \OperatorTok{\{}
        \KeywordTok{this}\OperatorTok{.}\FunctionTok{name} \OperatorTok{=}\NormalTok{ name}\OperatorTok{;}
        \KeywordTok{this}\OperatorTok{.}\FunctionTok{age} \OperatorTok{=}\NormalTok{ age}\OperatorTok{;}
    \OperatorTok{\}}

    \CommentTok{// Methods}
    \KeywordTok{public} \DataTypeTok{void} \FunctionTok{sit}\OperatorTok{()} \OperatorTok{\{}
        \BuiltInTok{System}\OperatorTok{.}\FunctionTok{out}\OperatorTok{.}\FunctionTok{println}\OperatorTok{(}\NormalTok{name }\OperatorTok{+} \StringTok{" is now sitting."}\OperatorTok{);}
    \OperatorTok{\}}

    \KeywordTok{public} \DataTypeTok{void} \FunctionTok{rollOver}\OperatorTok{()} \OperatorTok{\{}
        \BuiltInTok{System}\OperatorTok{.}\FunctionTok{out}\OperatorTok{.}\FunctionTok{println}\OperatorTok{(}\NormalTok{name }\OperatorTok{+} \StringTok{" rolled over!"}\OperatorTok{);}
    \OperatorTok{\}}

    \KeywordTok{public} \DataTypeTok{void} \FunctionTok{getOlder}\OperatorTok{(}\DataTypeTok{int}\NormalTok{ years}\OperatorTok{)} \OperatorTok{\{}
        \KeywordTok{this}\OperatorTok{.}\FunctionTok{age} \OperatorTok{+=}\NormalTok{ years}\OperatorTok{;}
        \BuiltInTok{System}\OperatorTok{.}\FunctionTok{out}\OperatorTok{.}\FunctionTok{println}\OperatorTok{(}\NormalTok{name }\OperatorTok{+} \StringTok{" is now "} \OperatorTok{+}\NormalTok{ age }\OperatorTok{+} \StringTok{" years old."}\OperatorTok{);}
    \OperatorTok{\}}

    \CommentTok{// Getters and Setters}
    \KeywordTok{public} \BuiltInTok{String} \FunctionTok{getName}\OperatorTok{()} \OperatorTok{\{} \ControlFlowTok{return}\NormalTok{ name}\OperatorTok{;} \OperatorTok{\}}
    \KeywordTok{public} \DataTypeTok{int} \FunctionTok{getAge}\OperatorTok{()} \OperatorTok{\{} \ControlFlowTok{return}\NormalTok{ age}\OperatorTok{;} \OperatorTok{\}}
    \KeywordTok{public} \DataTypeTok{void} \FunctionTok{setAge}\OperatorTok{(}\DataTypeTok{int}\NormalTok{ age}\OperatorTok{)} \OperatorTok{\{} \KeywordTok{this}\OperatorTok{.}\FunctionTok{age} \OperatorTok{=}\NormalTok{ age}\OperatorTok{;} \OperatorTok{\}}

    \AttributeTok{@Override}
    \KeywordTok{public} \BuiltInTok{String} \FunctionTok{toString}\OperatorTok{()} \OperatorTok{\{}
        \ControlFlowTok{return} \StringTok{"Dog\{name=\textquotesingle{}"} \OperatorTok{+}\NormalTok{ name }\OperatorTok{+} \StringTok{"\textquotesingle{}, age="} \OperatorTok{+}\NormalTok{ age }\OperatorTok{+} \StringTok{"\}"}\OperatorTok{;}
    \OperatorTok{\}}
\OperatorTok{\}}
\end{Highlighting}
\end{Shaded}

\subsubsection{Using the Class}\label{using-the-class}

\begin{Shaded}
\begin{Highlighting}[]
\KeywordTok{public} \KeywordTok{class}\NormalTok{ DogDemo }\OperatorTok{\{}
    \KeywordTok{public} \DataTypeTok{static} \DataTypeTok{void} \FunctionTok{main}\OperatorTok{(}\BuiltInTok{String}\OperatorTok{[]}\NormalTok{ args}\OperatorTok{)} \OperatorTok{\{}
\NormalTok{        Dog myDog }\OperatorTok{=} \KeywordTok{new} \FunctionTok{Dog}\OperatorTok{(}\StringTok{"Buddy"}\OperatorTok{,} \DecValTok{3}\OperatorTok{);}
\NormalTok{        Dog yourDog }\OperatorTok{=} \KeywordTok{new} \FunctionTok{Dog}\OperatorTok{(}\StringTok{"Lucy"}\OperatorTok{,} \DecValTok{1}\OperatorTok{);}

        \BuiltInTok{System}\OperatorTok{.}\FunctionTok{out}\OperatorTok{.}\FunctionTok{println}\OperatorTok{(}\StringTok{"My dog\textquotesingle{}s name is "} \OperatorTok{+}\NormalTok{ myDog}\OperatorTok{.}\FunctionTok{getName}\OperatorTok{()} \OperatorTok{+} \StringTok{"."}\OperatorTok{);}
\NormalTok{        myDog}\OperatorTok{.}\FunctionTok{sit}\OperatorTok{();}
\NormalTok{        myDog}\OperatorTok{.}\FunctionTok{rollOver}\OperatorTok{();}
\NormalTok{        myDog}\OperatorTok{.}\FunctionTok{getOlder}\OperatorTok{(}\DecValTok{2}\OperatorTok{);}

        \BuiltInTok{System}\OperatorTok{.}\FunctionTok{out}\OperatorTok{.}\FunctionTok{println}\OperatorTok{(}\NormalTok{myDog}\OperatorTok{);}
    \OperatorTok{\}}
\OperatorTok{\}}
\end{Highlighting}
\end{Shaded}

\subsubsection{Access Modifiers}\label{access-modifiers}

{\def\LTcaptype{none} % do not increment counter
\begin{longtable}[]{@{}ll@{}}
\toprule\noalign{}
Modifier & Access Level \\
\midrule\noalign{}
\endhead
\bottomrule\noalign{}
\endlastfoot
\texttt{public} & Accessible from anywhere \\
\texttt{private} & Only within the same class \\
\texttt{protected} & Within the class and subclasses \\
(default) & Within the same package \\
\end{longtable}
}

\textbf{Best practice:} Make instance variables \texttt{private} and
provide \texttt{public} getter/setter methods --- this is
\textbf{encapsulation}.

\subsubsection{Skills Checklist}\label{skills-checklist-15}

\begin{itemize}
\tightlist
\item[$\square$]
  I can define a Java class with fields, a constructor, and methods.
\item[$\square$]
  I can create instances of a class using \texttt{new}.
\item[$\square$]
  I understand the purpose of \texttt{this} in Java.
\item[$\square$]
  I can write getter and setter methods.
\item[$\square$]
  I understand access modifiers (\texttt{public}, \texttt{private}).
\end{itemize}

\begin{center}\rule{0.5\linewidth}{0.5pt}\end{center}

\subsection{Lesson 7.3: Inheritance}\label{lesson-7.3-inheritance}

\begin{Shaded}
\begin{Highlighting}[]
\KeywordTok{public} \KeywordTok{class}\NormalTok{ Vehicle }\OperatorTok{\{}
    \KeywordTok{protected} \BuiltInTok{String}\NormalTok{ make}\OperatorTok{;}
    \KeywordTok{protected} \BuiltInTok{String}\NormalTok{ model}\OperatorTok{;}
    \KeywordTok{protected} \DataTypeTok{int}\NormalTok{ year}\OperatorTok{;}

    \KeywordTok{public} \FunctionTok{Vehicle}\OperatorTok{(}\BuiltInTok{String}\NormalTok{ make}\OperatorTok{,} \BuiltInTok{String}\NormalTok{ model}\OperatorTok{,} \DataTypeTok{int}\NormalTok{ year}\OperatorTok{)} \OperatorTok{\{}
        \KeywordTok{this}\OperatorTok{.}\FunctionTok{make} \OperatorTok{=}\NormalTok{ make}\OperatorTok{;}
        \KeywordTok{this}\OperatorTok{.}\FunctionTok{model} \OperatorTok{=}\NormalTok{ model}\OperatorTok{;}
        \KeywordTok{this}\OperatorTok{.}\FunctionTok{year} \OperatorTok{=}\NormalTok{ year}\OperatorTok{;}
    \OperatorTok{\}}

    \KeywordTok{public} \BuiltInTok{String} \FunctionTok{getDescriptiveName}\OperatorTok{()} \OperatorTok{\{}
        \ControlFlowTok{return}\NormalTok{ year }\OperatorTok{+} \StringTok{" "} \OperatorTok{+}\NormalTok{ make }\OperatorTok{+} \StringTok{" "} \OperatorTok{+}\NormalTok{ model}\OperatorTok{;}
    \OperatorTok{\}}

    \KeywordTok{public} \DataTypeTok{void} \FunctionTok{readOdometer}\OperatorTok{()} \OperatorTok{\{}
        \BuiltInTok{System}\OperatorTok{.}\FunctionTok{out}\OperatorTok{.}\FunctionTok{println}\OperatorTok{(}\StringTok{"No odometer reading available."}\OperatorTok{);}
    \OperatorTok{\}}
\OperatorTok{\}}

\KeywordTok{public} \KeywordTok{class}\NormalTok{ Car }\KeywordTok{extends}\NormalTok{ Vehicle }\OperatorTok{\{}
    \KeywordTok{private} \DataTypeTok{int}\NormalTok{ odometer}\OperatorTok{;}

    \KeywordTok{public} \FunctionTok{Car}\OperatorTok{(}\BuiltInTok{String}\NormalTok{ make}\OperatorTok{,} \BuiltInTok{String}\NormalTok{ model}\OperatorTok{,} \DataTypeTok{int}\NormalTok{ year}\OperatorTok{)} \OperatorTok{\{}
        \KeywordTok{super}\OperatorTok{(}\NormalTok{make}\OperatorTok{,}\NormalTok{ model}\OperatorTok{,}\NormalTok{ year}\OperatorTok{);}   \CommentTok{// Call parent constructor}
        \KeywordTok{this}\OperatorTok{.}\FunctionTok{odometer} \OperatorTok{=} \DecValTok{0}\OperatorTok{;}
    \OperatorTok{\}}

    \AttributeTok{@Override}
    \KeywordTok{public} \DataTypeTok{void} \FunctionTok{readOdometer}\OperatorTok{()} \OperatorTok{\{}
        \BuiltInTok{System}\OperatorTok{.}\FunctionTok{out}\OperatorTok{.}\FunctionTok{println}\OperatorTok{(}\StringTok{"This car has "} \OperatorTok{+}\NormalTok{ odometer }\OperatorTok{+} \StringTok{" miles on it."}\OperatorTok{);}
    \OperatorTok{\}}

    \KeywordTok{public} \DataTypeTok{void} \FunctionTok{updateOdometer}\OperatorTok{(}\DataTypeTok{int}\NormalTok{ miles}\OperatorTok{)} \OperatorTok{\{}
        \ControlFlowTok{if} \OperatorTok{(}\NormalTok{miles }\OperatorTok{\textgreater{}=}\NormalTok{ odometer}\OperatorTok{)} \OperatorTok{\{}
\NormalTok{            odometer }\OperatorTok{=}\NormalTok{ miles}\OperatorTok{;}
        \OperatorTok{\}} \ControlFlowTok{else} \OperatorTok{\{}
            \BuiltInTok{System}\OperatorTok{.}\FunctionTok{out}\OperatorTok{.}\FunctionTok{println}\OperatorTok{(}\StringTok{"Error: Cannot roll back odometer."}\OperatorTok{);}
        \OperatorTok{\}}
    \OperatorTok{\}}
\OperatorTok{\}}

\KeywordTok{public} \KeywordTok{class}\NormalTok{ ElectricCar }\KeywordTok{extends}\NormalTok{ Car }\OperatorTok{\{}
    \KeywordTok{private} \DataTypeTok{int}\NormalTok{ batterySize}\OperatorTok{;}

    \KeywordTok{public} \FunctionTok{ElectricCar}\OperatorTok{(}\BuiltInTok{String}\NormalTok{ make}\OperatorTok{,} \BuiltInTok{String}\NormalTok{ model}\OperatorTok{,} \DataTypeTok{int}\NormalTok{ year}\OperatorTok{)} \OperatorTok{\{}
        \KeywordTok{super}\OperatorTok{(}\NormalTok{make}\OperatorTok{,}\NormalTok{ model}\OperatorTok{,}\NormalTok{ year}\OperatorTok{);}
        \KeywordTok{this}\OperatorTok{.}\FunctionTok{batterySize} \OperatorTok{=} \DecValTok{75}\OperatorTok{;}
    \OperatorTok{\}}

    \KeywordTok{public} \DataTypeTok{void} \FunctionTok{describeBattery}\OperatorTok{()} \OperatorTok{\{}
        \BuiltInTok{System}\OperatorTok{.}\FunctionTok{out}\OperatorTok{.}\FunctionTok{println}\OperatorTok{(}\StringTok{"This car has a "} \OperatorTok{+}\NormalTok{ batterySize }\OperatorTok{+} \StringTok{"{-}kWh battery."}\OperatorTok{);}
    \OperatorTok{\}}
\OperatorTok{\}}
\end{Highlighting}
\end{Shaded}

\subsubsection{Abstract Classes and
Interfaces}\label{abstract-classes-and-interfaces}

\begin{Shaded}
\begin{Highlighting}[]
\CommentTok{// Abstract class {-}{-}{-} cannot be instantiated}
\KeywordTok{public} \KeywordTok{abstract} \KeywordTok{class} \BuiltInTok{Shape} \OperatorTok{\{}
    \KeywordTok{public} \KeywordTok{abstract} \DataTypeTok{double} \FunctionTok{area}\OperatorTok{();}      \CommentTok{// Abstract method {-}{-}{-} no body}
    \KeywordTok{public} \KeywordTok{abstract} \DataTypeTok{double} \FunctionTok{perimeter}\OperatorTok{();} \CommentTok{// Must be implemented by subclasses}

    \KeywordTok{public} \DataTypeTok{void} \FunctionTok{printInfo}\OperatorTok{()} \OperatorTok{\{}           \CommentTok{// Concrete method {-}{-}{-} has body}
        \BuiltInTok{System}\OperatorTok{.}\FunctionTok{out}\OperatorTok{.}\FunctionTok{printf}\OperatorTok{(}\StringTok{"Area: }\SpecialCharTok{\%.2f}\StringTok{, Perimeter: }\SpecialCharTok{\%.2f\%n}\StringTok{"}\OperatorTok{,} \FunctionTok{area}\OperatorTok{(),} \FunctionTok{perimeter}\OperatorTok{());}
    \OperatorTok{\}}
\OperatorTok{\}}

\KeywordTok{public} \KeywordTok{class}\NormalTok{ Circle }\KeywordTok{extends} \BuiltInTok{Shape} \OperatorTok{\{}
    \KeywordTok{private} \DataTypeTok{double}\NormalTok{ radius}\OperatorTok{;}

    \KeywordTok{public} \FunctionTok{Circle}\OperatorTok{(}\DataTypeTok{double}\NormalTok{ radius}\OperatorTok{)} \OperatorTok{\{} \KeywordTok{this}\OperatorTok{.}\FunctionTok{radius} \OperatorTok{=}\NormalTok{ radius}\OperatorTok{;} \OperatorTok{\}}

    \AttributeTok{@Override}
    \KeywordTok{public} \DataTypeTok{double} \FunctionTok{area}\OperatorTok{()} \OperatorTok{\{} \ControlFlowTok{return} \BuiltInTok{Math}\OperatorTok{.}\FunctionTok{PI} \OperatorTok{*}\NormalTok{ radius }\OperatorTok{*}\NormalTok{ radius}\OperatorTok{;} \OperatorTok{\}}

    \AttributeTok{@Override}
    \KeywordTok{public} \DataTypeTok{double} \FunctionTok{perimeter}\OperatorTok{()} \OperatorTok{\{} \ControlFlowTok{return} \DecValTok{2} \OperatorTok{*} \BuiltInTok{Math}\OperatorTok{.}\FunctionTok{PI} \OperatorTok{*}\NormalTok{ radius}\OperatorTok{;} \OperatorTok{\}}
\OperatorTok{\}}

\CommentTok{// Interface {-}{-}{-} defines a contract}
\KeywordTok{public} \KeywordTok{interface}\NormalTok{ Drawable }\OperatorTok{\{}
    \DataTypeTok{void} \FunctionTok{draw}\OperatorTok{();}
    \KeywordTok{default} \DataTypeTok{void} \FunctionTok{resize}\OperatorTok{(}\DataTypeTok{double}\NormalTok{ factor}\OperatorTok{)} \OperatorTok{\{}
        \BuiltInTok{System}\OperatorTok{.}\FunctionTok{out}\OperatorTok{.}\FunctionTok{println}\OperatorTok{(}\StringTok{"Resizing by factor "} \OperatorTok{+}\NormalTok{ factor}\OperatorTok{);}
    \OperatorTok{\}}
\OperatorTok{\}}
\end{Highlighting}
\end{Shaded}

\subsubsection{Skills Checklist}\label{skills-checklist-16}

\begin{itemize}
\tightlist
\item[$\square$]
  I can create a subclass that extends a parent class.
\item[$\square$]
  I can use \texttt{super()} to call the parent constructor.
\item[$\square$]
  I can override methods using \texttt{@Override}.
\item[$\square$]
  I understand abstract classes and interfaces.
\end{itemize}

\begin{center}\rule{0.5\linewidth}{0.5pt}\end{center}

\section{Unit 8: Files and Exception
Handling}\label{unit-8-files-and-exception-handling-1}

\subsection{Overview}\label{overview-8}

This unit covers reading and writing files in Java and handling
exceptions.

\begin{center}\rule{0.5\linewidth}{0.5pt}\end{center}

\subsection{Lesson 8.1: Reading from
Files}\label{lesson-8.1-reading-from-files}

\begin{Shaded}
\begin{Highlighting}[]
\KeywordTok{import} \ImportTok{java}\OperatorTok{.}\ImportTok{io}\OperatorTok{.*;}
\KeywordTok{import} \ImportTok{java}\OperatorTok{.}\ImportTok{util}\OperatorTok{.*;}

\KeywordTok{public} \KeywordTok{class}\NormalTok{ FileReaderDemo }\OperatorTok{\{}
    \KeywordTok{public} \DataTypeTok{static} \DataTypeTok{void} \FunctionTok{main}\OperatorTok{(}\BuiltInTok{String}\OperatorTok{[]}\NormalTok{ args}\OperatorTok{)} \OperatorTok{\{}
        \CommentTok{// Reading entire file using Scanner}
        \ControlFlowTok{try} \OperatorTok{(}\BuiltInTok{Scanner}\NormalTok{ fileScanner }\OperatorTok{=} \KeywordTok{new} \BuiltInTok{Scanner}\OperatorTok{(}\KeywordTok{new} \BuiltInTok{File}\OperatorTok{(}\StringTok{"my\_notes.txt"}\OperatorTok{)))} \OperatorTok{\{}
            \ControlFlowTok{while} \OperatorTok{(}\NormalTok{fileScanner}\OperatorTok{.}\FunctionTok{hasNextLine}\OperatorTok{())} \OperatorTok{\{}
                \BuiltInTok{String}\NormalTok{ line }\OperatorTok{=}\NormalTok{ fileScanner}\OperatorTok{.}\FunctionTok{nextLine}\OperatorTok{();}
                \BuiltInTok{System}\OperatorTok{.}\FunctionTok{out}\OperatorTok{.}\FunctionTok{println}\OperatorTok{(}\NormalTok{line}\OperatorTok{);}
            \OperatorTok{\}}
        \OperatorTok{\}} \ControlFlowTok{catch} \OperatorTok{(}\BuiltInTok{FileNotFoundException}\NormalTok{ e}\OperatorTok{)} \OperatorTok{\{}
            \BuiltInTok{System}\OperatorTok{.}\FunctionTok{out}\OperatorTok{.}\FunctionTok{println}\OperatorTok{(}\StringTok{"Error: File not found {-} "} \OperatorTok{+}\NormalTok{ e}\OperatorTok{.}\FunctionTok{getMessage}\OperatorTok{());}
        \OperatorTok{\}}

        \CommentTok{// Reading with BufferedReader (more efficient for large files)}
        \ControlFlowTok{try} \OperatorTok{(}\BuiltInTok{BufferedReader}\NormalTok{ reader }\OperatorTok{=} \KeywordTok{new} \BuiltInTok{BufferedReader}\OperatorTok{(}\KeywordTok{new} \BuiltInTok{FileReader}\OperatorTok{(}\StringTok{"my\_notes.txt"}\OperatorTok{)))} \OperatorTok{\{}
            \BuiltInTok{String}\NormalTok{ line}\OperatorTok{;}
            \ControlFlowTok{while} \OperatorTok{((}\NormalTok{line }\OperatorTok{=}\NormalTok{ reader}\OperatorTok{.}\FunctionTok{readLine}\OperatorTok{())} \OperatorTok{!=} \KeywordTok{null}\OperatorTok{)} \OperatorTok{\{}
                \BuiltInTok{System}\OperatorTok{.}\FunctionTok{out}\OperatorTok{.}\FunctionTok{println}\OperatorTok{(}\NormalTok{line}\OperatorTok{);}
            \OperatorTok{\}}
        \OperatorTok{\}} \ControlFlowTok{catch} \OperatorTok{(}\BuiltInTok{IOException}\NormalTok{ e}\OperatorTok{)} \OperatorTok{\{}
            \BuiltInTok{System}\OperatorTok{.}\FunctionTok{out}\OperatorTok{.}\FunctionTok{println}\OperatorTok{(}\StringTok{"Error reading file: "} \OperatorTok{+}\NormalTok{ e}\OperatorTok{.}\FunctionTok{getMessage}\OperatorTok{());}
        \OperatorTok{\}}

        \CommentTok{// Reading all lines into a List}
        \ControlFlowTok{try} \OperatorTok{\{}
            \BuiltInTok{List}\OperatorTok{\textless{}}\BuiltInTok{String}\OperatorTok{\textgreater{}}\NormalTok{ lines }\OperatorTok{=} \KeywordTok{new}\NormalTok{ java}\OperatorTok{.}\FunctionTok{nio}\OperatorTok{.}\FunctionTok{file}\OperatorTok{.}\FunctionTok{Files}\OperatorTok{.}\FunctionTok{readAllLines}\OperatorTok{(}
\NormalTok{                java}\OperatorTok{.}\FunctionTok{nio}\OperatorTok{.}\FunctionTok{file}\OperatorTok{.}\FunctionTok{Paths}\OperatorTok{.}\FunctionTok{get}\OperatorTok{(}\StringTok{"my\_notes.txt"}\OperatorTok{)}
            \OperatorTok{);}
            \ControlFlowTok{for} \OperatorTok{(}\BuiltInTok{String}\NormalTok{ line }\OperatorTok{:}\NormalTok{ lines}\OperatorTok{)} \OperatorTok{\{}
                \BuiltInTok{System}\OperatorTok{.}\FunctionTok{out}\OperatorTok{.}\FunctionTok{println}\OperatorTok{(}\NormalTok{line}\OperatorTok{);}
            \OperatorTok{\}}
        \OperatorTok{\}} \ControlFlowTok{catch} \OperatorTok{(}\BuiltInTok{IOException}\NormalTok{ e}\OperatorTok{)} \OperatorTok{\{}
            \BuiltInTok{System}\OperatorTok{.}\FunctionTok{out}\OperatorTok{.}\FunctionTok{println}\OperatorTok{(}\StringTok{"Error: "} \OperatorTok{+}\NormalTok{ e}\OperatorTok{.}\FunctionTok{getMessage}\OperatorTok{());}
        \OperatorTok{\}}
    \OperatorTok{\}}
\OperatorTok{\}}
\end{Highlighting}
\end{Shaded}

\subsubsection{\texorpdfstring{Simpler approach using
\texttt{java.nio.file}}{Simpler approach using java.nio.file}}\label{simpler-approach-using-java.nio.file}

\begin{Shaded}
\begin{Highlighting}[]
\KeywordTok{import} \ImportTok{java}\OperatorTok{.}\ImportTok{nio}\OperatorTok{.}\ImportTok{file}\OperatorTok{.*;}
\KeywordTok{import} \ImportTok{java}\OperatorTok{.}\ImportTok{io}\OperatorTok{.}\ImportTok{IOException}\OperatorTok{;}
\KeywordTok{import} \ImportTok{java}\OperatorTok{.}\ImportTok{util}\OperatorTok{.}\ImportTok{List}\OperatorTok{;}

\KeywordTok{public} \KeywordTok{class}\NormalTok{ ModernFileReader }\OperatorTok{\{}
    \KeywordTok{public} \DataTypeTok{static} \DataTypeTok{void} \FunctionTok{main}\OperatorTok{(}\BuiltInTok{String}\OperatorTok{[]}\NormalTok{ args}\OperatorTok{)} \OperatorTok{\{}
\NormalTok{        Path filePath }\OperatorTok{=}\NormalTok{ Paths}\OperatorTok{.}\FunctionTok{get}\OperatorTok{(}\StringTok{"my\_notes.txt"}\OperatorTok{);}

        \ControlFlowTok{try} \OperatorTok{\{}
            \CommentTok{// Read all lines}
            \BuiltInTok{List}\OperatorTok{\textless{}}\BuiltInTok{String}\OperatorTok{\textgreater{}}\NormalTok{ lines }\OperatorTok{=}\NormalTok{ Files}\OperatorTok{.}\FunctionTok{readAllLines}\OperatorTok{(}\NormalTok{filePath}\OperatorTok{);}
            \ControlFlowTok{for} \OperatorTok{(}\BuiltInTok{String}\NormalTok{ line }\OperatorTok{:}\NormalTok{ lines}\OperatorTok{)} \OperatorTok{\{}
                \BuiltInTok{System}\OperatorTok{.}\FunctionTok{out}\OperatorTok{.}\FunctionTok{println}\OperatorTok{(}\NormalTok{line}\OperatorTok{);}
            \OperatorTok{\}}

            \CommentTok{// Read entire file as one string}
            \BuiltInTok{String}\NormalTok{ content }\OperatorTok{=}\NormalTok{ Files}\OperatorTok{.}\FunctionTok{readString}\OperatorTok{(}\NormalTok{filePath}\OperatorTok{);}
            \BuiltInTok{System}\OperatorTok{.}\FunctionTok{out}\OperatorTok{.}\FunctionTok{println}\OperatorTok{(}\NormalTok{content}\OperatorTok{);}

        \OperatorTok{\}} \ControlFlowTok{catch} \OperatorTok{(}\BuiltInTok{IOException}\NormalTok{ e}\OperatorTok{)} \OperatorTok{\{}
            \BuiltInTok{System}\OperatorTok{.}\FunctionTok{out}\OperatorTok{.}\FunctionTok{println}\OperatorTok{(}\StringTok{"Error reading file: "} \OperatorTok{+}\NormalTok{ e}\OperatorTok{.}\FunctionTok{getMessage}\OperatorTok{());}
        \OperatorTok{\}}
    \OperatorTok{\}}
\OperatorTok{\}}
\end{Highlighting}
\end{Shaded}

\begin{center}\rule{0.5\linewidth}{0.5pt}\end{center}

\subsection{Lesson 8.2: Writing to
Files}\label{lesson-8.2-writing-to-files}

\begin{Shaded}
\begin{Highlighting}[]
\KeywordTok{import} \ImportTok{java}\OperatorTok{.}\ImportTok{io}\OperatorTok{.*;}
\KeywordTok{import} \ImportTok{java}\OperatorTok{.}\ImportTok{nio}\OperatorTok{.}\ImportTok{file}\OperatorTok{.*;}

\KeywordTok{public} \KeywordTok{class}\NormalTok{ FileWriterDemo }\OperatorTok{\{}
    \KeywordTok{public} \DataTypeTok{static} \DataTypeTok{void} \FunctionTok{main}\OperatorTok{(}\BuiltInTok{String}\OperatorTok{[]}\NormalTok{ args}\OperatorTok{)} \OperatorTok{\{}
        \CommentTok{// Writing with PrintWriter (creates new or overwrites)}
        \ControlFlowTok{try} \OperatorTok{(}\BuiltInTok{PrintWriter}\NormalTok{ writer }\OperatorTok{=} \KeywordTok{new} \BuiltInTok{PrintWriter}\OperatorTok{(}\StringTok{"output.txt"}\OperatorTok{))} \OperatorTok{\{}
\NormalTok{            writer}\OperatorTok{.}\FunctionTok{println}\OperatorTok{(}\StringTok{"Hello, Java file I/O!"}\OperatorTok{);}
\NormalTok{            writer}\OperatorTok{.}\FunctionTok{println}\OperatorTok{(}\StringTok{"Writing is easy."}\OperatorTok{);}
\NormalTok{            writer}\OperatorTok{.}\FunctionTok{printf}\OperatorTok{(}\StringTok{"Value: }\SpecialCharTok{\%d\%n}\StringTok{"}\OperatorTok{,} \DecValTok{42}\OperatorTok{);}
        \OperatorTok{\}} \ControlFlowTok{catch} \OperatorTok{(}\BuiltInTok{FileNotFoundException}\NormalTok{ e}\OperatorTok{)} \OperatorTok{\{}
            \BuiltInTok{System}\OperatorTok{.}\FunctionTok{out}\OperatorTok{.}\FunctionTok{println}\OperatorTok{(}\StringTok{"Error: Could not create file."}\OperatorTok{);}
        \OperatorTok{\}}

        \CommentTok{// Appending with FileWriter(true = append mode)}
        \ControlFlowTok{try} \OperatorTok{(}\BuiltInTok{PrintWriter}\NormalTok{ appender }\OperatorTok{=} \KeywordTok{new} \BuiltInTok{PrintWriter}\OperatorTok{(}\KeywordTok{new} \BuiltInTok{FileWriter}\OperatorTok{(}\StringTok{"output.txt"}\OperatorTok{,} \KeywordTok{true}\OperatorTok{)))} \OperatorTok{\{}
\NormalTok{            appender}\OperatorTok{.}\FunctionTok{println}\OperatorTok{(}\StringTok{"This line is appended."}\OperatorTok{);}
        \OperatorTok{\}} \ControlFlowTok{catch} \OperatorTok{(}\BuiltInTok{IOException}\NormalTok{ e}\OperatorTok{)} \OperatorTok{\{}
            \BuiltInTok{System}\OperatorTok{.}\FunctionTok{out}\OperatorTok{.}\FunctionTok{println}\OperatorTok{(}\StringTok{"Error appending to file."}\OperatorTok{);}
        \OperatorTok{\}}

        \CommentTok{// Modern NIO approach}
        \ControlFlowTok{try} \OperatorTok{\{}
            \BuiltInTok{String}\NormalTok{ content }\OperatorTok{=} \StringTok{"Hello from NIO!}\SpecialCharTok{\textbackslash{}n}\StringTok{Second line."}\OperatorTok{;}
\NormalTok{            Files}\OperatorTok{.}\FunctionTok{writeString}\OperatorTok{(}\NormalTok{Paths}\OperatorTok{.}\FunctionTok{get}\OperatorTok{(}\StringTok{"modern\_output.txt"}\OperatorTok{),}\NormalTok{ content}\OperatorTok{);}

            \CommentTok{// Append}
\NormalTok{            Files}\OperatorTok{.}\FunctionTok{writeString}\OperatorTok{(}\NormalTok{Paths}\OperatorTok{.}\FunctionTok{get}\OperatorTok{(}\StringTok{"modern\_output.txt"}\OperatorTok{),}
                \StringTok{"}\SpecialCharTok{\textbackslash{}n}\StringTok{Appended line."}\OperatorTok{,}
\NormalTok{                StandardOpenOption}\OperatorTok{.}\FunctionTok{APPEND}\OperatorTok{);}
        \OperatorTok{\}} \ControlFlowTok{catch} \OperatorTok{(}\BuiltInTok{IOException}\NormalTok{ e}\OperatorTok{)} \OperatorTok{\{}
            \BuiltInTok{System}\OperatorTok{.}\FunctionTok{out}\OperatorTok{.}\FunctionTok{println}\OperatorTok{(}\StringTok{"Error: "} \OperatorTok{+}\NormalTok{ e}\OperatorTok{.}\FunctionTok{getMessage}\OperatorTok{());}
        \OperatorTok{\}}
    \OperatorTok{\}}
\OperatorTok{\}}
\end{Highlighting}
\end{Shaded}

\begin{center}\rule{0.5\linewidth}{0.5pt}\end{center}

\subsection{Lesson 8.3: Exception
Handling}\label{lesson-8.3-exception-handling}

\subsubsection{\texorpdfstring{\texttt{try-catch-finally}}{try-catch-finally}}\label{try-catch-finally}

\begin{Shaded}
\begin{Highlighting}[]
\KeywordTok{public} \KeywordTok{class}\NormalTok{ ExceptionHandling }\OperatorTok{\{}
    \KeywordTok{public} \DataTypeTok{static} \DataTypeTok{void} \FunctionTok{main}\OperatorTok{(}\BuiltInTok{String}\OperatorTok{[]}\NormalTok{ args}\OperatorTok{)} \OperatorTok{\{}
        \CommentTok{// Basic try{-}catch}
        \ControlFlowTok{try} \OperatorTok{\{}
            \DataTypeTok{int}\NormalTok{ result }\OperatorTok{=} \DecValTok{10} \OperatorTok{/} \DecValTok{0}\OperatorTok{;}
        \OperatorTok{\}} \ControlFlowTok{catch} \OperatorTok{(}\BuiltInTok{ArithmeticException}\NormalTok{ e}\OperatorTok{)} \OperatorTok{\{}
            \BuiltInTok{System}\OperatorTok{.}\FunctionTok{out}\OperatorTok{.}\FunctionTok{println}\OperatorTok{(}\StringTok{"Error: "} \OperatorTok{+}\NormalTok{ e}\OperatorTok{.}\FunctionTok{getMessage}\OperatorTok{());}
        \OperatorTok{\}}

        \CommentTok{// Multiple catch blocks}
        \ControlFlowTok{try} \OperatorTok{\{}
            \BuiltInTok{String}\NormalTok{ input }\OperatorTok{=} \StringTok{"abc"}\OperatorTok{;}
            \DataTypeTok{int}\NormalTok{ number }\OperatorTok{=} \BuiltInTok{Integer}\OperatorTok{.}\FunctionTok{parseInt}\OperatorTok{(}\NormalTok{input}\OperatorTok{);}
        \OperatorTok{\}} \ControlFlowTok{catch} \OperatorTok{(}\BuiltInTok{NumberFormatException}\NormalTok{ e}\OperatorTok{)} \OperatorTok{\{}
            \BuiltInTok{System}\OperatorTok{.}\FunctionTok{out}\OperatorTok{.}\FunctionTok{println}\OperatorTok{(}\StringTok{"Invalid number format: "} \OperatorTok{+}\NormalTok{ e}\OperatorTok{.}\FunctionTok{getMessage}\OperatorTok{());}
        \OperatorTok{\}} \ControlFlowTok{catch} \OperatorTok{(}\BuiltInTok{Exception}\NormalTok{ e}\OperatorTok{)} \OperatorTok{\{}
            \BuiltInTok{System}\OperatorTok{.}\FunctionTok{out}\OperatorTok{.}\FunctionTok{println}\OperatorTok{(}\StringTok{"Unexpected error: "} \OperatorTok{+}\NormalTok{ e}\OperatorTok{.}\FunctionTok{getMessage}\OperatorTok{());}
        \OperatorTok{\}}

        \CommentTok{// try{-}catch{-}finally}
        \BuiltInTok{PrintWriter}\NormalTok{ writer }\OperatorTok{=} \KeywordTok{null}\OperatorTok{;}
        \ControlFlowTok{try} \OperatorTok{\{}
\NormalTok{            writer }\OperatorTok{=} \KeywordTok{new} \BuiltInTok{PrintWriter}\OperatorTok{(}\StringTok{"data.txt"}\OperatorTok{);}
\NormalTok{            writer}\OperatorTok{.}\FunctionTok{println}\OperatorTok{(}\StringTok{"Some data"}\OperatorTok{);}
        \OperatorTok{\}} \ControlFlowTok{catch} \OperatorTok{(}\BuiltInTok{FileNotFoundException}\NormalTok{ e}\OperatorTok{)} \OperatorTok{\{}
            \BuiltInTok{System}\OperatorTok{.}\FunctionTok{out}\OperatorTok{.}\FunctionTok{println}\OperatorTok{(}\StringTok{"Could not create file."}\OperatorTok{);}
        \OperatorTok{\}} \ControlFlowTok{finally} \OperatorTok{\{}
            \ControlFlowTok{if} \OperatorTok{(}\NormalTok{writer }\OperatorTok{!=} \KeywordTok{null}\OperatorTok{)} \OperatorTok{\{}
\NormalTok{                writer}\OperatorTok{.}\FunctionTok{close}\OperatorTok{();}
                \BuiltInTok{System}\OperatorTok{.}\FunctionTok{out}\OperatorTok{.}\FunctionTok{println}\OperatorTok{(}\StringTok{"File closed."}\OperatorTok{);}
            \OperatorTok{\}}
        \OperatorTok{\}}

        \CommentTok{// try{-}with{-}resources (automatically closes resources)}
        \ControlFlowTok{try} \OperatorTok{(}\BuiltInTok{PrintWriter}\NormalTok{ pw }\OperatorTok{=} \KeywordTok{new} \BuiltInTok{PrintWriter}\OperatorTok{(}\StringTok{"data2.txt"}\OperatorTok{))} \OperatorTok{\{}
\NormalTok{            pw}\OperatorTok{.}\FunctionTok{println}\OperatorTok{(}\StringTok{"Better approach!"}\OperatorTok{);}
        \OperatorTok{\}} \ControlFlowTok{catch} \OperatorTok{(}\BuiltInTok{FileNotFoundException}\NormalTok{ e}\OperatorTok{)} \OperatorTok{\{}
            \BuiltInTok{System}\OperatorTok{.}\FunctionTok{out}\OperatorTok{.}\FunctionTok{println}\OperatorTok{(}\StringTok{"Could not create file."}\OperatorTok{);}
        \OperatorTok{\}}
    \OperatorTok{\}}
\OperatorTok{\}}
\end{Highlighting}
\end{Shaded}

\subsubsection{Common Java Exceptions}\label{common-java-exceptions}

{\def\LTcaptype{none} % do not increment counter
\begin{longtable}[]{@{}
  >{\raggedright\arraybackslash}p{(\linewidth - 2\tabcolsep) * \real{0.3516}}
  >{\raggedright\arraybackslash}p{(\linewidth - 2\tabcolsep) * \real{0.6484}}@{}}
\toprule\noalign{}
\begin{minipage}[b]{\linewidth}\raggedright
Exception
\end{minipage} & \begin{minipage}[b]{\linewidth}\raggedright
When it occurs
\end{minipage} \\
\midrule\noalign{}
\endhead
\bottomrule\noalign{}
\endlastfoot
\texttt{ArithmeticException} & Division by zero \\
\texttt{NumberFormatException} & Invalid number conversion (e.g.,
\texttt{Integer.parseInt("abc")}) \\
\texttt{ArrayIndexOutOfBoundsException} & Array index out of range \\
\texttt{NullPointerException} & Calling method on null reference \\
\texttt{FileNotFoundException} & File doesn't exist \\
\texttt{IOException} & General I/O error \\
\texttt{ClassCastException} & Invalid type cast \\
\end{longtable}
}

\subsubsection{Custom Exceptions}\label{custom-exceptions}

\begin{Shaded}
\begin{Highlighting}[]
\KeywordTok{public} \KeywordTok{class}\NormalTok{ InsufficientFundsException }\KeywordTok{extends} \BuiltInTok{Exception} \OperatorTok{\{}
    \KeywordTok{private} \DataTypeTok{double}\NormalTok{ amount}\OperatorTok{;}

    \KeywordTok{public} \FunctionTok{InsufficientFundsException}\OperatorTok{(}\DataTypeTok{double}\NormalTok{ amount}\OperatorTok{)} \OperatorTok{\{}
        \KeywordTok{super}\OperatorTok{(}\StringTok{"Insufficient funds: need "} \OperatorTok{+}\NormalTok{ amount }\OperatorTok{+} \StringTok{" more."}\OperatorTok{);}
        \KeywordTok{this}\OperatorTok{.}\FunctionTok{amount} \OperatorTok{=}\NormalTok{ amount}\OperatorTok{;}
    \OperatorTok{\}}

    \KeywordTok{public} \DataTypeTok{double} \FunctionTok{getAmount}\OperatorTok{()} \OperatorTok{\{} \ControlFlowTok{return}\NormalTok{ amount}\OperatorTok{;} \OperatorTok{\}}
\OperatorTok{\}}

\CommentTok{// Usage:}
\KeywordTok{public} \DataTypeTok{void} \FunctionTok{withdraw}\OperatorTok{(}\DataTypeTok{double}\NormalTok{ amount}\OperatorTok{)} \KeywordTok{throws}\NormalTok{ InsufficientFundsException }\OperatorTok{\{}
    \ControlFlowTok{if} \OperatorTok{(}\NormalTok{amount }\OperatorTok{\textgreater{}}\NormalTok{ balance}\OperatorTok{)} \OperatorTok{\{}
        \ControlFlowTok{throw} \KeywordTok{new} \FunctionTok{InsufficientFundsException}\OperatorTok{(}\NormalTok{amount }\OperatorTok{{-}}\NormalTok{ balance}\OperatorTok{);}
    \OperatorTok{\}}
\NormalTok{    balance }\OperatorTok{{-}=}\NormalTok{ amount}\OperatorTok{;}
\OperatorTok{\}}
\end{Highlighting}
\end{Shaded}

\subsubsection{Skills Checklist}\label{skills-checklist-17}

\begin{itemize}
\tightlist
\item[$\square$]
  I can read from a file using \texttt{Scanner} or
  \texttt{BufferedReader}.
\item[$\square$]
  I can write to a file using \texttt{PrintWriter}.
\item[$\square$]
  I can use \texttt{try-catch-finally} to handle exceptions.
\item[$\square$]
  I can use try-with-resources for automatic resource cleanup.
\item[$\square$]
  I know common Java exception types.
\end{itemize}

\newpage

\section{Unit 9: Collections, Data Analysis, and
Testing}\label{unit-9-collections-data-analysis-and-testing}

\subsection{Overview}\label{overview-9}

This unit introduces Java's Collections Framework, basic data analysis,
and unit testing with JUnit.

\begin{center}\rule{0.5\linewidth}{0.5pt}\end{center}

\subsection{Lesson 9.1: Java Collections
Framework}\label{lesson-9.1-java-collections-framework}

\subsubsection{Key Interfaces and
Classes}\label{key-interfaces-and-classes}

\begin{verbatim}
Collection
|--- List
|   |--- ArrayList
|   +--- LinkedList
+--- Set
    |--- HashSet
    +--- TreeSet

Map
|--- HashMap
+--- TreeMap
\end{verbatim}

\subsubsection{\texorpdfstring{\texttt{HashMap} --- Key-Value
Storage}{HashMap --- Key-Value Storage}}\label{hashmap-key-value-storage}

\begin{Shaded}
\begin{Highlighting}[]
\KeywordTok{import} \ImportTok{java}\OperatorTok{.}\ImportTok{util}\OperatorTok{.}\ImportTok{HashMap}\OperatorTok{;}
\KeywordTok{import} \ImportTok{java}\OperatorTok{.}\ImportTok{util}\OperatorTok{.}\ImportTok{Map}\OperatorTok{;}

\BuiltInTok{HashMap}\OperatorTok{\textless{}}\BuiltInTok{String}\OperatorTok{,} \BuiltInTok{Integer}\OperatorTok{\textgreater{}}\NormalTok{ inventory }\OperatorTok{=} \KeywordTok{new} \BuiltInTok{HashMap}\OperatorTok{\textless{}\textgreater{}();}

\CommentTok{// Add entries}
\NormalTok{inventory}\OperatorTok{.}\FunctionTok{put}\OperatorTok{(}\StringTok{"apples"}\OperatorTok{,} \DecValTok{50}\OperatorTok{);}
\NormalTok{inventory}\OperatorTok{.}\FunctionTok{put}\OperatorTok{(}\StringTok{"bananas"}\OperatorTok{,} \DecValTok{30}\OperatorTok{);}
\NormalTok{inventory}\OperatorTok{.}\FunctionTok{put}\OperatorTok{(}\StringTok{"oranges"}\OperatorTok{,} \DecValTok{40}\OperatorTok{);}

\CommentTok{// Access}
\BuiltInTok{System}\OperatorTok{.}\FunctionTok{out}\OperatorTok{.}\FunctionTok{println}\OperatorTok{(}\NormalTok{inventory}\OperatorTok{.}\FunctionTok{get}\OperatorTok{(}\StringTok{"apples"}\OperatorTok{));}           \CommentTok{// 50}
\BuiltInTok{System}\OperatorTok{.}\FunctionTok{out}\OperatorTok{.}\FunctionTok{println}\OperatorTok{(}\NormalTok{inventory}\OperatorTok{.}\FunctionTok{getOrDefault}\OperatorTok{(}\StringTok{"grapes"}\OperatorTok{,} \DecValTok{0}\OperatorTok{));} \CommentTok{// 0}

\CommentTok{// Check}
\BuiltInTok{System}\OperatorTok{.}\FunctionTok{out}\OperatorTok{.}\FunctionTok{println}\OperatorTok{(}\NormalTok{inventory}\OperatorTok{.}\FunctionTok{containsKey}\OperatorTok{(}\StringTok{"bananas"}\OperatorTok{));}   \CommentTok{// true}
\BuiltInTok{System}\OperatorTok{.}\FunctionTok{out}\OperatorTok{.}\FunctionTok{println}\OperatorTok{(}\NormalTok{inventory}\OperatorTok{.}\FunctionTok{containsValue}\OperatorTok{(}\DecValTok{30}\OperatorTok{));}         \CommentTok{// true}

\CommentTok{// Update}
\NormalTok{inventory}\OperatorTok{.}\FunctionTok{put}\OperatorTok{(}\StringTok{"apples"}\OperatorTok{,} \DecValTok{55}\OperatorTok{);}

\CommentTok{// Remove}
\NormalTok{inventory}\OperatorTok{.}\FunctionTok{remove}\OperatorTok{(}\StringTok{"bananas"}\OperatorTok{);}

\CommentTok{// Iterate}
\ControlFlowTok{for} \OperatorTok{(}\BuiltInTok{Map}\OperatorTok{.}\FunctionTok{Entry}\OperatorTok{\textless{}}\BuiltInTok{String}\OperatorTok{,} \BuiltInTok{Integer}\OperatorTok{\textgreater{}}\NormalTok{ entry }\OperatorTok{:}\NormalTok{ inventory}\OperatorTok{.}\FunctionTok{entrySet}\OperatorTok{())} \OperatorTok{\{}
    \BuiltInTok{System}\OperatorTok{.}\FunctionTok{out}\OperatorTok{.}\FunctionTok{println}\OperatorTok{(}\NormalTok{entry}\OperatorTok{.}\FunctionTok{getKey}\OperatorTok{()} \OperatorTok{+} \StringTok{": "} \OperatorTok{+}\NormalTok{ entry}\OperatorTok{.}\FunctionTok{getValue}\OperatorTok{());}
\OperatorTok{\}}

\CommentTok{// Keys only}
\ControlFlowTok{for} \OperatorTok{(}\BuiltInTok{String}\NormalTok{ key }\OperatorTok{:}\NormalTok{ inventory}\OperatorTok{.}\FunctionTok{keySet}\OperatorTok{())} \OperatorTok{\{}
    \BuiltInTok{System}\OperatorTok{.}\FunctionTok{out}\OperatorTok{.}\FunctionTok{println}\OperatorTok{(}\NormalTok{key}\OperatorTok{);}
\OperatorTok{\}}

\CommentTok{// Values only}
\ControlFlowTok{for} \OperatorTok{(}\DataTypeTok{int}\NormalTok{ value }\OperatorTok{:}\NormalTok{ inventory}\OperatorTok{.}\FunctionTok{values}\OperatorTok{())} \OperatorTok{\{}
    \BuiltInTok{System}\OperatorTok{.}\FunctionTok{out}\OperatorTok{.}\FunctionTok{println}\OperatorTok{(}\NormalTok{value}\OperatorTok{);}
\OperatorTok{\}}

\BuiltInTok{System}\OperatorTok{.}\FunctionTok{out}\OperatorTok{.}\FunctionTok{println}\OperatorTok{(}\StringTok{"Size: "} \OperatorTok{+}\NormalTok{ inventory}\OperatorTok{.}\FunctionTok{size}\OperatorTok{());}
\end{Highlighting}
\end{Shaded}

\subsubsection{\texorpdfstring{\texttt{HashSet} --- Unique
Elements}{HashSet --- Unique Elements}}\label{hashset-unique-elements}

\begin{Shaded}
\begin{Highlighting}[]
\KeywordTok{import} \ImportTok{java}\OperatorTok{.}\ImportTok{util}\OperatorTok{.}\ImportTok{HashSet}\OperatorTok{;}

\BuiltInTok{HashSet}\OperatorTok{\textless{}}\BuiltInTok{String}\OperatorTok{\textgreater{}}\NormalTok{ uniqueWords }\OperatorTok{=} \KeywordTok{new} \BuiltInTok{HashSet}\OperatorTok{\textless{}\textgreater{}();}
\NormalTok{uniqueWords}\OperatorTok{.}\FunctionTok{add}\OperatorTok{(}\StringTok{"apple"}\OperatorTok{);}
\NormalTok{uniqueWords}\OperatorTok{.}\FunctionTok{add}\OperatorTok{(}\StringTok{"banana"}\OperatorTok{);}
\NormalTok{uniqueWords}\OperatorTok{.}\FunctionTok{add}\OperatorTok{(}\StringTok{"apple"}\OperatorTok{);}   \CommentTok{// Duplicate {-}{-}{-} ignored}

\BuiltInTok{System}\OperatorTok{.}\FunctionTok{out}\OperatorTok{.}\FunctionTok{println}\OperatorTok{(}\NormalTok{uniqueWords}\OperatorTok{.}\FunctionTok{size}\OperatorTok{());}          \CommentTok{// 2}
\BuiltInTok{System}\OperatorTok{.}\FunctionTok{out}\OperatorTok{.}\FunctionTok{println}\OperatorTok{(}\NormalTok{uniqueWords}\OperatorTok{.}\FunctionTok{contains}\OperatorTok{(}\StringTok{"banana"}\OperatorTok{));} \CommentTok{// true}
\end{Highlighting}
\end{Shaded}

\subsubsection{\texorpdfstring{\texttt{TreeMap} --- Sorted
Keys}{TreeMap --- Sorted Keys}}\label{treemap-sorted-keys}

\begin{Shaded}
\begin{Highlighting}[]
\KeywordTok{import} \ImportTok{java}\OperatorTok{.}\ImportTok{util}\OperatorTok{.}\ImportTok{TreeMap}\OperatorTok{;}

\BuiltInTok{TreeMap}\OperatorTok{\textless{}}\BuiltInTok{String}\OperatorTok{,} \BuiltInTok{Integer}\OperatorTok{\textgreater{}}\NormalTok{ sorted }\OperatorTok{=} \KeywordTok{new} \BuiltInTok{TreeMap}\OperatorTok{\textless{}\textgreater{}();}
\NormalTok{sorted}\OperatorTok{.}\FunctionTok{put}\OperatorTok{(}\StringTok{"banana"}\OperatorTok{,} \DecValTok{2}\OperatorTok{);}
\NormalTok{sorted}\OperatorTok{.}\FunctionTok{put}\OperatorTok{(}\StringTok{"apple"}\OperatorTok{,} \DecValTok{1}\OperatorTok{);}
\NormalTok{sorted}\OperatorTok{.}\FunctionTok{put}\OperatorTok{(}\StringTok{"cherry"}\OperatorTok{,} \DecValTok{3}\OperatorTok{);}

\CommentTok{// Iterates in alphabetical key order}
\ControlFlowTok{for} \OperatorTok{(}\BuiltInTok{Map}\OperatorTok{.}\FunctionTok{Entry}\OperatorTok{\textless{}}\BuiltInTok{String}\OperatorTok{,} \BuiltInTok{Integer}\OperatorTok{\textgreater{}}\NormalTok{ e }\OperatorTok{:}\NormalTok{ sorted}\OperatorTok{.}\FunctionTok{entrySet}\OperatorTok{())} \OperatorTok{\{}
    \BuiltInTok{System}\OperatorTok{.}\FunctionTok{out}\OperatorTok{.}\FunctionTok{println}\OperatorTok{(}\NormalTok{e}\OperatorTok{.}\FunctionTok{getKey}\OperatorTok{()} \OperatorTok{+} \StringTok{": "} \OperatorTok{+}\NormalTok{ e}\OperatorTok{.}\FunctionTok{getValue}\OperatorTok{());}
\OperatorTok{\}}
\end{Highlighting}
\end{Shaded}

\subsubsection{Sorting with Collections}\label{sorting-with-collections}

\begin{Shaded}
\begin{Highlighting}[]
\KeywordTok{import} \ImportTok{java}\OperatorTok{.}\ImportTok{util}\OperatorTok{.*;}

\BuiltInTok{List}\OperatorTok{\textless{}}\BuiltInTok{Integer}\OperatorTok{\textgreater{}}\NormalTok{ numbers }\OperatorTok{=} \KeywordTok{new} \BuiltInTok{ArrayList}\OperatorTok{\textless{}\textgreater{}(}\BuiltInTok{Arrays}\OperatorTok{.}\FunctionTok{asList}\OperatorTok{(}\DecValTok{5}\OperatorTok{,} \DecValTok{2}\OperatorTok{,} \DecValTok{9}\OperatorTok{,} \DecValTok{1}\OperatorTok{,} \DecValTok{7}\OperatorTok{));}
\BuiltInTok{Collections}\OperatorTok{.}\FunctionTok{sort}\OperatorTok{(}\NormalTok{numbers}\OperatorTok{);}
\BuiltInTok{System}\OperatorTok{.}\FunctionTok{out}\OperatorTok{.}\FunctionTok{println}\OperatorTok{(}\NormalTok{numbers}\OperatorTok{);}   \CommentTok{// [1, 2, 5, 7, 9]}

\BuiltInTok{Collections}\OperatorTok{.}\FunctionTok{sort}\OperatorTok{(}\NormalTok{numbers}\OperatorTok{,} \BuiltInTok{Collections}\OperatorTok{.}\FunctionTok{reverseOrder}\OperatorTok{());}
\BuiltInTok{System}\OperatorTok{.}\FunctionTok{out}\OperatorTok{.}\FunctionTok{println}\OperatorTok{(}\NormalTok{numbers}\OperatorTok{);}   \CommentTok{// [9, 7, 5, 2, 1]}

\DataTypeTok{int}\NormalTok{ max }\OperatorTok{=} \BuiltInTok{Collections}\OperatorTok{.}\FunctionTok{max}\OperatorTok{(}\NormalTok{numbers}\OperatorTok{);}
\DataTypeTok{int}\NormalTok{ min }\OperatorTok{=} \BuiltInTok{Collections}\OperatorTok{.}\FunctionTok{min}\OperatorTok{(}\NormalTok{numbers}\OperatorTok{);}
\end{Highlighting}
\end{Shaded}

\begin{center}\rule{0.5\linewidth}{0.5pt}\end{center}

\subsection{Lesson 9.2: Word Frequency
Analysis}\label{lesson-9.2-word-frequency-analysis}

\begin{Shaded}
\begin{Highlighting}[]
\KeywordTok{import} \ImportTok{java}\OperatorTok{.}\ImportTok{util}\OperatorTok{.*;}

\KeywordTok{public} \KeywordTok{class}\NormalTok{ WordFrequency }\OperatorTok{\{}
    \KeywordTok{public} \DataTypeTok{static} \BuiltInTok{Map}\OperatorTok{\textless{}}\BuiltInTok{String}\OperatorTok{,} \BuiltInTok{Integer}\OperatorTok{\textgreater{}} \FunctionTok{countWords}\OperatorTok{(}\BuiltInTok{String}\NormalTok{ text}\OperatorTok{)} \OperatorTok{\{}
        \BuiltInTok{Map}\OperatorTok{\textless{}}\BuiltInTok{String}\OperatorTok{,} \BuiltInTok{Integer}\OperatorTok{\textgreater{}}\NormalTok{ frequency }\OperatorTok{=} \KeywordTok{new} \BuiltInTok{HashMap}\OperatorTok{\textless{}\textgreater{}();}
        \BuiltInTok{String}\OperatorTok{[]}\NormalTok{ words }\OperatorTok{=}\NormalTok{ text}\OperatorTok{.}\FunctionTok{toLowerCase}\OperatorTok{()}
                             \OperatorTok{.}\FunctionTok{replaceAll}\OperatorTok{(}\StringTok{"[\^{}a{-}z}\SpecialCharTok{\textbackslash{}\textbackslash{}}\StringTok{s]"}\OperatorTok{,} \StringTok{""}\OperatorTok{)}
                             \OperatorTok{.}\FunctionTok{trim}\OperatorTok{()}
                             \OperatorTok{.}\FunctionTok{split}\OperatorTok{(}\StringTok{"}\SpecialCharTok{\textbackslash{}\textbackslash{}}\StringTok{s+"}\OperatorTok{);}

        \ControlFlowTok{for} \OperatorTok{(}\BuiltInTok{String}\NormalTok{ word }\OperatorTok{:}\NormalTok{ words}\OperatorTok{)} \OperatorTok{\{}
\NormalTok{            frequency}\OperatorTok{.}\FunctionTok{put}\OperatorTok{(}\NormalTok{word}\OperatorTok{,}\NormalTok{ frequency}\OperatorTok{.}\FunctionTok{getOrDefault}\OperatorTok{(}\NormalTok{word}\OperatorTok{,} \DecValTok{0}\OperatorTok{)} \OperatorTok{+} \DecValTok{1}\OperatorTok{);}
        \OperatorTok{\}}
        \ControlFlowTok{return}\NormalTok{ frequency}\OperatorTok{;}
    \OperatorTok{\}}

    \KeywordTok{public} \DataTypeTok{static} \DataTypeTok{void} \FunctionTok{main}\OperatorTok{(}\BuiltInTok{String}\OperatorTok{[]}\NormalTok{ args}\OperatorTok{)} \OperatorTok{\{}
        \BuiltInTok{String}\NormalTok{ text }\OperatorTok{=} \StringTok{"the quick brown fox jumps over the lazy dog the quick brown fox"}\OperatorTok{;}
        \BuiltInTok{Map}\OperatorTok{\textless{}}\BuiltInTok{String}\OperatorTok{,} \BuiltInTok{Integer}\OperatorTok{\textgreater{}}\NormalTok{ freq }\OperatorTok{=} \FunctionTok{countWords}\OperatorTok{(}\NormalTok{text}\OperatorTok{);}

        \CommentTok{// Sort by frequency (most common first)}
        \BuiltInTok{List}\OperatorTok{\textless{}}\BuiltInTok{Map}\OperatorTok{.}\FunctionTok{Entry}\OperatorTok{\textless{}}\BuiltInTok{String}\OperatorTok{,} \BuiltInTok{Integer}\OperatorTok{\textgreater{}\textgreater{}}\NormalTok{ entries }\OperatorTok{=} \KeywordTok{new} \BuiltInTok{ArrayList}\OperatorTok{\textless{}\textgreater{}(}\NormalTok{freq}\OperatorTok{.}\FunctionTok{entrySet}\OperatorTok{());}
\NormalTok{        entries}\OperatorTok{.}\FunctionTok{sort}\OperatorTok{((}\NormalTok{a}\OperatorTok{,}\NormalTok{ b}\OperatorTok{)} \OperatorTok{{-}\textgreater{}}\NormalTok{ b}\OperatorTok{.}\FunctionTok{getValue}\OperatorTok{()} \OperatorTok{{-}}\NormalTok{ a}\OperatorTok{.}\FunctionTok{getValue}\OperatorTok{());}

        \BuiltInTok{System}\OperatorTok{.}\FunctionTok{out}\OperatorTok{.}\FunctionTok{println}\OperatorTok{(}\StringTok{"Word Frequencies:"}\OperatorTok{);}
        \ControlFlowTok{for} \OperatorTok{(}\BuiltInTok{Map}\OperatorTok{.}\FunctionTok{Entry}\OperatorTok{\textless{}}\BuiltInTok{String}\OperatorTok{,} \BuiltInTok{Integer}\OperatorTok{\textgreater{}}\NormalTok{ entry }\OperatorTok{:}\NormalTok{ entries}\OperatorTok{)} \OperatorTok{\{}
            \BuiltInTok{System}\OperatorTok{.}\FunctionTok{out}\OperatorTok{.}\FunctionTok{printf}\OperatorTok{(}\StringTok{"\textquotesingle{}}\SpecialCharTok{\%s}\StringTok{\textquotesingle{}: }\SpecialCharTok{\%d\%n}\StringTok{"}\OperatorTok{,}\NormalTok{ entry}\OperatorTok{.}\FunctionTok{getKey}\OperatorTok{(),}\NormalTok{ entry}\OperatorTok{.}\FunctionTok{getValue}\OperatorTok{());}
        \OperatorTok{\}}
    \OperatorTok{\}}
\OperatorTok{\}}
\end{Highlighting}
\end{Shaded}

\begin{center}\rule{0.5\linewidth}{0.5pt}\end{center}

\subsection{Lesson 9.3: Working with CSV
Files}\label{lesson-9.3-working-with-csv-files}

\begin{Shaded}
\begin{Highlighting}[]
\KeywordTok{import} \ImportTok{java}\OperatorTok{.}\ImportTok{io}\OperatorTok{.*;}
\KeywordTok{import} \ImportTok{java}\OperatorTok{.}\ImportTok{util}\OperatorTok{.*;}

\KeywordTok{public} \KeywordTok{class}\NormalTok{ CSVReader }\OperatorTok{\{}
    \KeywordTok{public} \DataTypeTok{static} \DataTypeTok{void} \FunctionTok{main}\OperatorTok{(}\BuiltInTok{String}\OperatorTok{[]}\NormalTok{ args}\OperatorTok{)} \OperatorTok{\{}
        \BuiltInTok{String}\NormalTok{ filename }\OperatorTok{=} \StringTok{"students.csv"}\OperatorTok{;}

        \ControlFlowTok{try} \OperatorTok{(}\BuiltInTok{BufferedReader}\NormalTok{ reader }\OperatorTok{=} \KeywordTok{new} \BuiltInTok{BufferedReader}\OperatorTok{(}\KeywordTok{new} \BuiltInTok{FileReader}\OperatorTok{(}\NormalTok{filename}\OperatorTok{)))} \OperatorTok{\{}
            \BuiltInTok{String}\NormalTok{ headerLine }\OperatorTok{=}\NormalTok{ reader}\OperatorTok{.}\FunctionTok{readLine}\OperatorTok{();}
            \BuiltInTok{String}\OperatorTok{[]}\NormalTok{ headers }\OperatorTok{=}\NormalTok{ headerLine}\OperatorTok{.}\FunctionTok{split}\OperatorTok{(}\StringTok{","}\OperatorTok{);}
            \BuiltInTok{System}\OperatorTok{.}\FunctionTok{out}\OperatorTok{.}\FunctionTok{println}\OperatorTok{(}\StringTok{"Columns: "} \OperatorTok{+} \BuiltInTok{Arrays}\OperatorTok{.}\FunctionTok{toString}\OperatorTok{(}\NormalTok{headers}\OperatorTok{));}
            \BuiltInTok{System}\OperatorTok{.}\FunctionTok{out}\OperatorTok{.}\FunctionTok{println}\OperatorTok{(}\StringTok{"{-}{-}{-}"}\OperatorTok{);}

            \BuiltInTok{String}\NormalTok{ line}\OperatorTok{;}
            \ControlFlowTok{while} \OperatorTok{((}\NormalTok{line }\OperatorTok{=}\NormalTok{ reader}\OperatorTok{.}\FunctionTok{readLine}\OperatorTok{())} \OperatorTok{!=} \KeywordTok{null}\OperatorTok{)} \OperatorTok{\{}
                \BuiltInTok{String}\OperatorTok{[]}\NormalTok{ data }\OperatorTok{=}\NormalTok{ line}\OperatorTok{.}\FunctionTok{split}\OperatorTok{(}\StringTok{","}\OperatorTok{);}
                \ControlFlowTok{if} \OperatorTok{(}\NormalTok{data}\OperatorTok{.}\FunctionTok{length} \OperatorTok{==}\NormalTok{ headers}\OperatorTok{.}\FunctionTok{length}\OperatorTok{)} \OperatorTok{\{}
                    \BuiltInTok{String}\NormalTok{ name }\OperatorTok{=}\NormalTok{ data}\OperatorTok{[}\DecValTok{0}\OperatorTok{].}\FunctionTok{trim}\OperatorTok{();}
                    \DataTypeTok{int}\NormalTok{ age }\OperatorTok{=} \BuiltInTok{Integer}\OperatorTok{.}\FunctionTok{parseInt}\OperatorTok{(}\NormalTok{data}\OperatorTok{[}\DecValTok{1}\OperatorTok{].}\FunctionTok{trim}\OperatorTok{());}
                    \BuiltInTok{String}\NormalTok{ grade }\OperatorTok{=}\NormalTok{ data}\OperatorTok{[}\DecValTok{2}\OperatorTok{].}\FunctionTok{trim}\OperatorTok{();}
                    \BuiltInTok{System}\OperatorTok{.}\FunctionTok{out}\OperatorTok{.}\FunctionTok{printf}\OperatorTok{(}\StringTok{"Name: }\SpecialCharTok{\%{-}15s}\StringTok{ Age: }\SpecialCharTok{\%3d}\StringTok{  Grade: }\SpecialCharTok{\%s\%n}\StringTok{"}\OperatorTok{,}
\NormalTok{                                      name}\OperatorTok{,}\NormalTok{ age}\OperatorTok{,}\NormalTok{ grade}\OperatorTok{);}
                \OperatorTok{\}}
            \OperatorTok{\}}
        \OperatorTok{\}} \ControlFlowTok{catch} \OperatorTok{(}\BuiltInTok{FileNotFoundException}\NormalTok{ e}\OperatorTok{)} \OperatorTok{\{}
            \BuiltInTok{System}\OperatorTok{.}\FunctionTok{out}\OperatorTok{.}\FunctionTok{println}\OperatorTok{(}\StringTok{"Error: File not found {-} "} \OperatorTok{+}\NormalTok{ filename}\OperatorTok{);}
        \OperatorTok{\}} \ControlFlowTok{catch} \OperatorTok{(}\BuiltInTok{IOException}\NormalTok{ e}\OperatorTok{)} \OperatorTok{\{}
            \BuiltInTok{System}\OperatorTok{.}\FunctionTok{out}\OperatorTok{.}\FunctionTok{println}\OperatorTok{(}\StringTok{"Error reading file: "} \OperatorTok{+}\NormalTok{ e}\OperatorTok{.}\FunctionTok{getMessage}\OperatorTok{());}
        \OperatorTok{\}} \ControlFlowTok{catch} \OperatorTok{(}\BuiltInTok{NumberFormatException}\NormalTok{ e}\OperatorTok{)} \OperatorTok{\{}
            \BuiltInTok{System}\OperatorTok{.}\FunctionTok{out}\OperatorTok{.}\FunctionTok{println}\OperatorTok{(}\StringTok{"Error: Invalid number format in CSV."}\OperatorTok{);}
        \OperatorTok{\}}
    \OperatorTok{\}}
\OperatorTok{\}}
\end{Highlighting}
\end{Shaded}

\begin{center}\rule{0.5\linewidth}{0.5pt}\end{center}

\subsection{Lesson 9.4: Unit Testing with
JUnit}\label{lesson-9.4-unit-testing-with-junit}

\subsubsection{Adding JUnit to Your
Project}\label{adding-junit-to-your-project}

If using VS Code with Maven, add to \texttt{pom.xml}:

\begin{Shaded}
\begin{Highlighting}[]
\NormalTok{\textless{}}\KeywordTok{dependency}\NormalTok{\textgreater{}}
\NormalTok{    \textless{}}\KeywordTok{groupId}\NormalTok{\textgreater{}org.junit.jupiter\textless{}/}\KeywordTok{groupId}\NormalTok{\textgreater{}}
\NormalTok{    \textless{}}\KeywordTok{artifactId}\NormalTok{\textgreater{}junit{-}jupiter\textless{}/}\KeywordTok{artifactId}\NormalTok{\textgreater{}}
\NormalTok{    \textless{}}\KeywordTok{version}\NormalTok{\textgreater{}5.10.0\textless{}/}\KeywordTok{version}\NormalTok{\textgreater{}}
\NormalTok{    \textless{}}\KeywordTok{scope}\NormalTok{\textgreater{}test\textless{}/}\KeywordTok{scope}\NormalTok{\textgreater{}}
\NormalTok{\textless{}/}\KeywordTok{dependency}\NormalTok{\textgreater{}}
\end{Highlighting}
\end{Shaded}

\subsubsection{Writing Tests}\label{writing-tests}

\begin{Shaded}
\begin{Highlighting}[]
\CommentTok{// MyMath.java}
\KeywordTok{public} \KeywordTok{class}\NormalTok{ MyMath }\OperatorTok{\{}
    \KeywordTok{public} \DataTypeTok{static} \DataTypeTok{int} \FunctionTok{add}\OperatorTok{(}\DataTypeTok{int}\NormalTok{ a}\OperatorTok{,} \DataTypeTok{int}\NormalTok{ b}\OperatorTok{)} \OperatorTok{\{} \ControlFlowTok{return}\NormalTok{ a }\OperatorTok{+}\NormalTok{ b}\OperatorTok{;} \OperatorTok{\}}
    \KeywordTok{public} \DataTypeTok{static} \DataTypeTok{double} \FunctionTok{circleArea}\OperatorTok{(}\DataTypeTok{double}\NormalTok{ radius}\OperatorTok{)} \OperatorTok{\{} \ControlFlowTok{return} \BuiltInTok{Math}\OperatorTok{.}\FunctionTok{PI} \OperatorTok{*}\NormalTok{ radius }\OperatorTok{*}\NormalTok{ radius}\OperatorTok{;} \OperatorTok{\}}
    \KeywordTok{public} \DataTypeTok{static} \BuiltInTok{String} \FunctionTok{getGrade}\OperatorTok{(}\DataTypeTok{int}\NormalTok{ score}\OperatorTok{)} \OperatorTok{\{}
        \ControlFlowTok{if} \OperatorTok{(}\NormalTok{score }\OperatorTok{\textgreater{}=} \DecValTok{90}\OperatorTok{)} \ControlFlowTok{return} \StringTok{"A"}\OperatorTok{;}
        \ControlFlowTok{if} \OperatorTok{(}\NormalTok{score }\OperatorTok{\textgreater{}=} \DecValTok{80}\OperatorTok{)} \ControlFlowTok{return} \StringTok{"B"}\OperatorTok{;}
        \ControlFlowTok{if} \OperatorTok{(}\NormalTok{score }\OperatorTok{\textgreater{}=} \DecValTok{70}\OperatorTok{)} \ControlFlowTok{return} \StringTok{"C"}\OperatorTok{;}
        \ControlFlowTok{return} \StringTok{"F"}\OperatorTok{;}
    \OperatorTok{\}}
\OperatorTok{\}}
\end{Highlighting}
\end{Shaded}

\begin{Shaded}
\begin{Highlighting}[]
\CommentTok{// TestMyMath.java}
\KeywordTok{import} \ImportTok{org}\OperatorTok{.}\ImportTok{junit}\OperatorTok{.}\ImportTok{jupiter}\OperatorTok{.}\ImportTok{api}\OperatorTok{.*;}
\KeywordTok{import} \KeywordTok{static}\ImportTok{ org}\OperatorTok{.}\ImportTok{junit}\OperatorTok{.}\ImportTok{jupiter}\OperatorTok{.}\ImportTok{api}\OperatorTok{.}\ImportTok{Assertions}\OperatorTok{.*;}

\KeywordTok{public} \KeywordTok{class}\NormalTok{ TestMyMath }\OperatorTok{\{}

    \AttributeTok{@Test}
    \KeywordTok{public} \DataTypeTok{void} \FunctionTok{testAddPositiveNumbers}\OperatorTok{()} \OperatorTok{\{}
        \FunctionTok{assertEquals}\OperatorTok{(}\DecValTok{8}\OperatorTok{,}\NormalTok{ MyMath}\OperatorTok{.}\FunctionTok{add}\OperatorTok{(}\DecValTok{5}\OperatorTok{,} \DecValTok{3}\OperatorTok{));}
    \OperatorTok{\}}

    \AttributeTok{@Test}
    \KeywordTok{public} \DataTypeTok{void} \FunctionTok{testAddNegativeNumbers}\OperatorTok{()} \OperatorTok{\{}
        \FunctionTok{assertEquals}\OperatorTok{({-}}\DecValTok{8}\OperatorTok{,}\NormalTok{ MyMath}\OperatorTok{.}\FunctionTok{add}\OperatorTok{({-}}\DecValTok{5}\OperatorTok{,} \OperatorTok{{-}}\DecValTok{3}\OperatorTok{));}
    \OperatorTok{\}}

    \AttributeTok{@Test}
    \KeywordTok{public} \DataTypeTok{void} \FunctionTok{testCircleArea}\OperatorTok{()} \OperatorTok{\{}
        \DataTypeTok{double}\NormalTok{ expected }\OperatorTok{=} \BuiltInTok{Math}\OperatorTok{.}\FunctionTok{PI} \OperatorTok{*} \DecValTok{25}\OperatorTok{;}
        \FunctionTok{assertEquals}\OperatorTok{(}\NormalTok{expected}\OperatorTok{,}\NormalTok{ MyMath}\OperatorTok{.}\FunctionTok{circleArea}\OperatorTok{(}\FloatTok{5.0}\OperatorTok{),} \FloatTok{0.0001}\OperatorTok{);}
    \OperatorTok{\}}

    \AttributeTok{@Test}
    \KeywordTok{public} \DataTypeTok{void} \FunctionTok{testGradeA}\OperatorTok{()} \OperatorTok{\{}
        \FunctionTok{assertEquals}\OperatorTok{(}\StringTok{"A"}\OperatorTok{,}\NormalTok{ MyMath}\OperatorTok{.}\FunctionTok{getGrade}\OperatorTok{(}\DecValTok{95}\OperatorTok{));}
    \OperatorTok{\}}

    \AttributeTok{@Test}
    \KeywordTok{public} \DataTypeTok{void} \FunctionTok{testGradeF}\OperatorTok{()} \OperatorTok{\{}
        \FunctionTok{assertEquals}\OperatorTok{(}\StringTok{"F"}\OperatorTok{,}\NormalTok{ MyMath}\OperatorTok{.}\FunctionTok{getGrade}\OperatorTok{(}\DecValTok{55}\OperatorTok{));}
    \OperatorTok{\}}
\OperatorTok{\}}
\end{Highlighting}
\end{Shaded}

\subsubsection{Common JUnit Assertions}\label{common-junit-assertions}

{\def\LTcaptype{none} % do not increment counter
\begin{longtable}[]{@{}
  >{\raggedright\arraybackslash}p{(\linewidth - 2\tabcolsep) * \real{0.6269}}
  >{\raggedright\arraybackslash}p{(\linewidth - 2\tabcolsep) * \real{0.3731}}@{}}
\toprule\noalign{}
\begin{minipage}[b]{\linewidth}\raggedright
Method
\end{minipage} & \begin{minipage}[b]{\linewidth}\raggedright
Description
\end{minipage} \\
\midrule\noalign{}
\endhead
\bottomrule\noalign{}
\endlastfoot
\texttt{assertEquals(expected,\ actual)} & Check equality \\
\texttt{assertTrue(condition)} & Check condition is true \\
\texttt{assertFalse(condition)} & Check condition is false \\
\texttt{assertNull(object)} & Check object is null \\
\texttt{assertNotNull(object)} & Check object is not null \\
\texttt{assertThrows(ExceptionClass,\ executable)} & Check exception is
thrown \\
\end{longtable}
}

\subsubsection{Skills Checklist}\label{skills-checklist-18}

\begin{itemize}
\tightlist
\item[$\square$]
  I can use \texttt{HashMap} for key-value storage in Java.
\item[$\square$]
  I can iterate through a map's entries, keys, and values.
\item[$\square$]
  I can write unit tests using JUnit 5.
\item[$\square$]
  I can use common assertion methods.
\end{itemize}

\begin{center}\rule{0.5\linewidth}{0.5pt}\end{center}

\section{Unit 10: Project-Based
Learning}\label{unit-10-project-based-learning}

\subsection{Overview}\label{overview-10}

This unit guides students through building text-based Java applications.

\begin{center}\rule{0.5\linewidth}{0.5pt}\end{center}

\subsection{Lesson 10.1: Designing a Text-Based Adventure
Game}\label{lesson-10.1-designing-a-text-based-adventure-game}

\subsubsection{Room and Player
Structure}\label{room-and-player-structure}

\begin{Shaded}
\begin{Highlighting}[]
\KeywordTok{import} \ImportTok{java}\OperatorTok{.}\ImportTok{util}\OperatorTok{.*;}

\KeywordTok{public} \KeywordTok{class}\NormalTok{ Room }\OperatorTok{\{}
    \KeywordTok{private} \BuiltInTok{String}\NormalTok{ name}\OperatorTok{;}
    \KeywordTok{private} \BuiltInTok{String}\NormalTok{ description}\OperatorTok{;}
    \KeywordTok{private} \BuiltInTok{Map}\OperatorTok{\textless{}}\BuiltInTok{String}\OperatorTok{,}\NormalTok{ Room}\OperatorTok{\textgreater{}}\NormalTok{ exits}\OperatorTok{;}
    \KeywordTok{private} \BuiltInTok{List}\OperatorTok{\textless{}}\BuiltInTok{String}\OperatorTok{\textgreater{}}\NormalTok{ items}\OperatorTok{;}

    \KeywordTok{public} \FunctionTok{Room}\OperatorTok{(}\BuiltInTok{String}\NormalTok{ name}\OperatorTok{,} \BuiltInTok{String}\NormalTok{ description}\OperatorTok{)} \OperatorTok{\{}
        \KeywordTok{this}\OperatorTok{.}\FunctionTok{name} \OperatorTok{=}\NormalTok{ name}\OperatorTok{;}
        \KeywordTok{this}\OperatorTok{.}\FunctionTok{description} \OperatorTok{=}\NormalTok{ description}\OperatorTok{;}
        \KeywordTok{this}\OperatorTok{.}\FunctionTok{exits} \OperatorTok{=} \KeywordTok{new} \BuiltInTok{HashMap}\OperatorTok{\textless{}\textgreater{}();}
        \KeywordTok{this}\OperatorTok{.}\FunctionTok{items} \OperatorTok{=} \KeywordTok{new} \BuiltInTok{ArrayList}\OperatorTok{\textless{}\textgreater{}();}
    \OperatorTok{\}}

    \KeywordTok{public} \DataTypeTok{void} \FunctionTok{addExit}\OperatorTok{(}\BuiltInTok{String}\NormalTok{ direction}\OperatorTok{,}\NormalTok{ Room room}\OperatorTok{)} \OperatorTok{\{}
\NormalTok{        exits}\OperatorTok{.}\FunctionTok{put}\OperatorTok{(}\NormalTok{direction}\OperatorTok{.}\FunctionTok{toLowerCase}\OperatorTok{(),}\NormalTok{ room}\OperatorTok{);}
    \OperatorTok{\}}

    \KeywordTok{public} \DataTypeTok{void} \FunctionTok{addItem}\OperatorTok{(}\BuiltInTok{String}\NormalTok{ item}\OperatorTok{)} \OperatorTok{\{}
\NormalTok{        items}\OperatorTok{.}\FunctionTok{add}\OperatorTok{(}\NormalTok{item}\OperatorTok{);}
    \OperatorTok{\}}

    \KeywordTok{public}\NormalTok{ Room }\FunctionTok{getExit}\OperatorTok{(}\BuiltInTok{String}\NormalTok{ direction}\OperatorTok{)} \OperatorTok{\{}
        \ControlFlowTok{return}\NormalTok{ exits}\OperatorTok{.}\FunctionTok{get}\OperatorTok{(}\NormalTok{direction}\OperatorTok{.}\FunctionTok{toLowerCase}\OperatorTok{());}
    \OperatorTok{\}}

    \KeywordTok{public} \BuiltInTok{String} \FunctionTok{getDescription}\OperatorTok{()} \OperatorTok{\{}
        \BuiltInTok{StringBuilder}\NormalTok{ sb }\OperatorTok{=} \KeywordTok{new} \BuiltInTok{StringBuilder}\OperatorTok{();}
\NormalTok{        sb}\OperatorTok{.}\FunctionTok{append}\OperatorTok{(}\StringTok{"{-}{-}{-} "}\OperatorTok{).}\FunctionTok{append}\OperatorTok{(}\NormalTok{name}\OperatorTok{).}\FunctionTok{append}\OperatorTok{(}\StringTok{" {-}{-}{-}}\SpecialCharTok{\textbackslash{}n}\StringTok{"}\OperatorTok{);}
\NormalTok{        sb}\OperatorTok{.}\FunctionTok{append}\OperatorTok{(}\NormalTok{description}\OperatorTok{).}\FunctionTok{append}\OperatorTok{(}\StringTok{"}\SpecialCharTok{\textbackslash{}n}\StringTok{"}\OperatorTok{);}
        \ControlFlowTok{if} \OperatorTok{(!}\NormalTok{items}\OperatorTok{.}\FunctionTok{isEmpty}\OperatorTok{())} \OperatorTok{\{}
\NormalTok{            sb}\OperatorTok{.}\FunctionTok{append}\OperatorTok{(}\StringTok{"You see: "}\OperatorTok{).}\FunctionTok{append}\OperatorTok{(}\BuiltInTok{String}\OperatorTok{.}\FunctionTok{join}\OperatorTok{(}\StringTok{", "}\OperatorTok{,}\NormalTok{ items}\OperatorTok{)).}\FunctionTok{append}\OperatorTok{(}\StringTok{"}\SpecialCharTok{\textbackslash{}n}\StringTok{"}\OperatorTok{);}
        \OperatorTok{\}}
\NormalTok{        sb}\OperatorTok{.}\FunctionTok{append}\OperatorTok{(}\StringTok{"Exits: "}\OperatorTok{).}\FunctionTok{append}\OperatorTok{(}\BuiltInTok{String}\OperatorTok{.}\FunctionTok{join}\OperatorTok{(}\StringTok{", "}\OperatorTok{,}\NormalTok{ exits}\OperatorTok{.}\FunctionTok{keySet}\OperatorTok{()));}
        \ControlFlowTok{return}\NormalTok{ sb}\OperatorTok{.}\FunctionTok{toString}\OperatorTok{();}
    \OperatorTok{\}}

    \KeywordTok{public} \BuiltInTok{String} \FunctionTok{getName}\OperatorTok{()} \OperatorTok{\{} \ControlFlowTok{return}\NormalTok{ name}\OperatorTok{;} \OperatorTok{\}}
    \KeywordTok{public} \BuiltInTok{List}\OperatorTok{\textless{}}\BuiltInTok{String}\OperatorTok{\textgreater{}} \FunctionTok{getItems}\OperatorTok{()} \OperatorTok{\{} \ControlFlowTok{return}\NormalTok{ items}\OperatorTok{;} \OperatorTok{\}}
\OperatorTok{\}}
\end{Highlighting}
\end{Shaded}

\begin{Shaded}
\begin{Highlighting}[]
\KeywordTok{import} \ImportTok{java}\OperatorTok{.}\ImportTok{util}\OperatorTok{.*;}

\KeywordTok{public} \KeywordTok{class}\NormalTok{ Player }\OperatorTok{\{}
    \KeywordTok{private} \BuiltInTok{String}\NormalTok{ name}\OperatorTok{;}
    \KeywordTok{private}\NormalTok{ Room currentRoom}\OperatorTok{;}
    \KeywordTok{private} \BuiltInTok{List}\OperatorTok{\textless{}}\BuiltInTok{String}\OperatorTok{\textgreater{}}\NormalTok{ inventory}\OperatorTok{;}
    \KeywordTok{private} \DataTypeTok{int}\NormalTok{ health}\OperatorTok{;}

    \KeywordTok{public} \FunctionTok{Player}\OperatorTok{(}\BuiltInTok{String}\NormalTok{ name}\OperatorTok{,}\NormalTok{ Room startRoom}\OperatorTok{)} \OperatorTok{\{}
        \KeywordTok{this}\OperatorTok{.}\FunctionTok{name} \OperatorTok{=}\NormalTok{ name}\OperatorTok{;}
        \KeywordTok{this}\OperatorTok{.}\FunctionTok{currentRoom} \OperatorTok{=}\NormalTok{ startRoom}\OperatorTok{;}
        \KeywordTok{this}\OperatorTok{.}\FunctionTok{inventory} \OperatorTok{=} \KeywordTok{new} \BuiltInTok{ArrayList}\OperatorTok{\textless{}\textgreater{}();}
        \KeywordTok{this}\OperatorTok{.}\FunctionTok{health} \OperatorTok{=} \DecValTok{100}\OperatorTok{;}
    \OperatorTok{\}}

    \KeywordTok{public} \DataTypeTok{boolean} \FunctionTok{move}\OperatorTok{(}\BuiltInTok{String}\NormalTok{ direction}\OperatorTok{)} \OperatorTok{\{}
\NormalTok{        Room nextRoom }\OperatorTok{=}\NormalTok{ currentRoom}\OperatorTok{.}\FunctionTok{getExit}\OperatorTok{(}\NormalTok{direction}\OperatorTok{);}
        \ControlFlowTok{if} \OperatorTok{(}\NormalTok{nextRoom }\OperatorTok{!=} \KeywordTok{null}\OperatorTok{)} \OperatorTok{\{}
\NormalTok{            currentRoom }\OperatorTok{=}\NormalTok{ nextRoom}\OperatorTok{;}
            \ControlFlowTok{return} \KeywordTok{true}\OperatorTok{;}
        \OperatorTok{\}}
        \ControlFlowTok{return} \KeywordTok{false}\OperatorTok{;}
    \OperatorTok{\}}

    \KeywordTok{public} \DataTypeTok{boolean} \FunctionTok{takeItem}\OperatorTok{(}\BuiltInTok{String}\NormalTok{ itemName}\OperatorTok{)} \OperatorTok{\{}
        \BuiltInTok{List}\OperatorTok{\textless{}}\BuiltInTok{String}\OperatorTok{\textgreater{}}\NormalTok{ roomItems }\OperatorTok{=}\NormalTok{ currentRoom}\OperatorTok{.}\FunctionTok{getItems}\OperatorTok{();}
        \ControlFlowTok{for} \OperatorTok{(}\BuiltInTok{String}\NormalTok{ item }\OperatorTok{:}\NormalTok{ roomItems}\OperatorTok{)} \OperatorTok{\{}
            \ControlFlowTok{if} \OperatorTok{(}\NormalTok{item}\OperatorTok{.}\FunctionTok{equalsIgnoreCase}\OperatorTok{(}\NormalTok{itemName}\OperatorTok{))} \OperatorTok{\{}
\NormalTok{                inventory}\OperatorTok{.}\FunctionTok{add}\OperatorTok{(}\NormalTok{item}\OperatorTok{);}
\NormalTok{                roomItems}\OperatorTok{.}\FunctionTok{remove}\OperatorTok{(}\NormalTok{item}\OperatorTok{);}
                \ControlFlowTok{return} \KeywordTok{true}\OperatorTok{;}
            \OperatorTok{\}}
        \OperatorTok{\}}
        \ControlFlowTok{return} \KeywordTok{false}\OperatorTok{;}
    \OperatorTok{\}}

    \KeywordTok{public} \BuiltInTok{String} \FunctionTok{getStatus}\OperatorTok{()} \OperatorTok{\{}
        \ControlFlowTok{return} \StringTok{"Player: "} \OperatorTok{+}\NormalTok{ name }\OperatorTok{+} \StringTok{" | Health: "} \OperatorTok{+}\NormalTok{ health }\OperatorTok{+}
               \StringTok{" | Inventory: "} \OperatorTok{+} \OperatorTok{(}\NormalTok{inventory}\OperatorTok{.}\FunctionTok{isEmpty}\OperatorTok{()} \OperatorTok{?} \StringTok{"empty"} \OperatorTok{:} \BuiltInTok{String}\OperatorTok{.}\FunctionTok{join}\OperatorTok{(}\StringTok{", "}\OperatorTok{,}\NormalTok{ inventory}\OperatorTok{));}
    \OperatorTok{\}}

    \KeywordTok{public}\NormalTok{ Room }\FunctionTok{getCurrentRoom}\OperatorTok{()} \OperatorTok{\{} \ControlFlowTok{return}\NormalTok{ currentRoom}\OperatorTok{;} \OperatorTok{\}}
    \KeywordTok{public} \BuiltInTok{List}\OperatorTok{\textless{}}\BuiltInTok{String}\OperatorTok{\textgreater{}} \FunctionTok{getInventory}\OperatorTok{()} \OperatorTok{\{} \ControlFlowTok{return}\NormalTok{ inventory}\OperatorTok{;} \OperatorTok{\}}
\OperatorTok{\}}
\end{Highlighting}
\end{Shaded}

\begin{Shaded}
\begin{Highlighting}[]
\KeywordTok{import} \ImportTok{java}\OperatorTok{.}\ImportTok{util}\OperatorTok{.}\ImportTok{Scanner}\OperatorTok{;}

\KeywordTok{public} \KeywordTok{class}\NormalTok{ AdventureGame }\OperatorTok{\{}
    \KeywordTok{public} \DataTypeTok{static} \DataTypeTok{void} \FunctionTok{main}\OperatorTok{(}\BuiltInTok{String}\OperatorTok{[]}\NormalTok{ args}\OperatorTok{)} \OperatorTok{\{}
        \CommentTok{// Build the world}
\NormalTok{        Room cave }\OperatorTok{=} \KeywordTok{new} \FunctionTok{Room}\OperatorTok{(}\StringTok{"Dark Cave"}\OperatorTok{,} \StringTok{"You are in a dimly lit cave."}\OperatorTok{);}
\NormalTok{        Room tunnel }\OperatorTok{=} \KeywordTok{new} \FunctionTok{Room}\OperatorTok{(}\StringTok{"Stone Tunnel"}\OperatorTok{,} \StringTok{"A narrow tunnel leads north and south."}\OperatorTok{);}
\NormalTok{        Room clearing }\OperatorTok{=} \KeywordTok{new} \FunctionTok{Room}\OperatorTok{(}\StringTok{"Sunny Clearing"}\OperatorTok{,} \StringTok{"A bright clearing with tall trees."}\OperatorTok{);}

\NormalTok{        cave}\OperatorTok{.}\FunctionTok{addExit}\OperatorTok{(}\StringTok{"north"}\OperatorTok{,}\NormalTok{ tunnel}\OperatorTok{);}
\NormalTok{        cave}\OperatorTok{.}\FunctionTok{addItem}\OperatorTok{(}\StringTok{"torch"}\OperatorTok{);}
\NormalTok{        tunnel}\OperatorTok{.}\FunctionTok{addExit}\OperatorTok{(}\StringTok{"south"}\OperatorTok{,}\NormalTok{ cave}\OperatorTok{);}
\NormalTok{        tunnel}\OperatorTok{.}\FunctionTok{addExit}\OperatorTok{(}\StringTok{"north"}\OperatorTok{,}\NormalTok{ clearing}\OperatorTok{);}
\NormalTok{        clearing}\OperatorTok{.}\FunctionTok{addExit}\OperatorTok{(}\StringTok{"south"}\OperatorTok{,}\NormalTok{ tunnel}\OperatorTok{);}

\NormalTok{        Player player }\OperatorTok{=} \KeywordTok{new} \FunctionTok{Player}\OperatorTok{(}\StringTok{"Hero"}\OperatorTok{,}\NormalTok{ cave}\OperatorTok{);}
        \BuiltInTok{Scanner}\NormalTok{ scanner }\OperatorTok{=} \KeywordTok{new} \BuiltInTok{Scanner}\OperatorTok{(}\BuiltInTok{System}\OperatorTok{.}\FunctionTok{in}\OperatorTok{);}

        \BuiltInTok{System}\OperatorTok{.}\FunctionTok{out}\OperatorTok{.}\FunctionTok{println}\OperatorTok{(}\StringTok{"=== Text Adventure ==="}\OperatorTok{);}
        \BuiltInTok{System}\OperatorTok{.}\FunctionTok{out}\OperatorTok{.}\FunctionTok{println}\OperatorTok{(}\StringTok{"Commands: go [direction], take [item], inventory, look, quit"}\OperatorTok{);}
        \BuiltInTok{System}\OperatorTok{.}\FunctionTok{out}\OperatorTok{.}\FunctionTok{println}\OperatorTok{();}

        \ControlFlowTok{while} \OperatorTok{(}\KeywordTok{true}\OperatorTok{)} \OperatorTok{\{}
            \BuiltInTok{System}\OperatorTok{.}\FunctionTok{out}\OperatorTok{.}\FunctionTok{println}\OperatorTok{(}\NormalTok{player}\OperatorTok{.}\FunctionTok{getCurrentRoom}\OperatorTok{().}\FunctionTok{getDescription}\OperatorTok{());}
            \BuiltInTok{System}\OperatorTok{.}\FunctionTok{out}\OperatorTok{.}\FunctionTok{println}\OperatorTok{(}\NormalTok{player}\OperatorTok{.}\FunctionTok{getStatus}\OperatorTok{());}
            \BuiltInTok{System}\OperatorTok{.}\FunctionTok{out}\OperatorTok{.}\FunctionTok{print}\OperatorTok{(}\StringTok{"}\SpecialCharTok{\textbackslash{}n}\StringTok{\textgreater{} "}\OperatorTok{);}

            \BuiltInTok{String}\NormalTok{ input }\OperatorTok{=}\NormalTok{ scanner}\OperatorTok{.}\FunctionTok{nextLine}\OperatorTok{().}\FunctionTok{trim}\OperatorTok{().}\FunctionTok{toLowerCase}\OperatorTok{();}
            \BuiltInTok{String}\OperatorTok{[]}\NormalTok{ parts }\OperatorTok{=}\NormalTok{ input}\OperatorTok{.}\FunctionTok{split}\OperatorTok{(}\StringTok{"}\SpecialCharTok{\textbackslash{}\textbackslash{}}\StringTok{s+"}\OperatorTok{,} \DecValTok{2}\OperatorTok{);}
            \BuiltInTok{String}\NormalTok{ command }\OperatorTok{=}\NormalTok{ parts}\OperatorTok{[}\DecValTok{0}\OperatorTok{];}
            \BuiltInTok{String}\NormalTok{ argument }\OperatorTok{=}\NormalTok{ parts}\OperatorTok{.}\FunctionTok{length} \OperatorTok{\textgreater{}} \DecValTok{1} \OperatorTok{?}\NormalTok{ parts}\OperatorTok{[}\DecValTok{1}\OperatorTok{]} \OperatorTok{:} \StringTok{""}\OperatorTok{;}

            \ControlFlowTok{switch} \OperatorTok{(}\NormalTok{command}\OperatorTok{)} \OperatorTok{\{}
                \ControlFlowTok{case} \StringTok{"go"}\OperatorTok{:}
                    \ControlFlowTok{if} \OperatorTok{(}\NormalTok{argument}\OperatorTok{.}\FunctionTok{isEmpty}\OperatorTok{())} \OperatorTok{\{}
                        \BuiltInTok{System}\OperatorTok{.}\FunctionTok{out}\OperatorTok{.}\FunctionTok{println}\OperatorTok{(}\StringTok{"Go where?"}\OperatorTok{);}
                    \OperatorTok{\}} \ControlFlowTok{else} \ControlFlowTok{if} \OperatorTok{(!}\NormalTok{player}\OperatorTok{.}\FunctionTok{move}\OperatorTok{(}\NormalTok{argument}\OperatorTok{))} \OperatorTok{\{}
                        \BuiltInTok{System}\OperatorTok{.}\FunctionTok{out}\OperatorTok{.}\FunctionTok{println}\OperatorTok{(}\StringTok{"You can\textquotesingle{}t go that way."}\OperatorTok{);}
                    \OperatorTok{\}}
                    \ControlFlowTok{break}\OperatorTok{;}
                \ControlFlowTok{case} \StringTok{"take"}\OperatorTok{:}
                    \ControlFlowTok{if} \OperatorTok{(}\NormalTok{player}\OperatorTok{.}\FunctionTok{takeItem}\OperatorTok{(}\NormalTok{argument}\OperatorTok{))} \OperatorTok{\{}
                        \BuiltInTok{System}\OperatorTok{.}\FunctionTok{out}\OperatorTok{.}\FunctionTok{println}\OperatorTok{(}\StringTok{"You picked up the "} \OperatorTok{+}\NormalTok{ argument }\OperatorTok{+} \StringTok{"."}\OperatorTok{);}
                    \OperatorTok{\}} \ControlFlowTok{else} \OperatorTok{\{}
                        \BuiltInTok{System}\OperatorTok{.}\FunctionTok{out}\OperatorTok{.}\FunctionTok{println}\OperatorTok{(}\StringTok{"There\textquotesingle{}s no "} \OperatorTok{+}\NormalTok{ argument }\OperatorTok{+} \StringTok{" here."}\OperatorTok{);}
                    \OperatorTok{\}}
                    \ControlFlowTok{break}\OperatorTok{;}
                \ControlFlowTok{case} \StringTok{"inventory"}\OperatorTok{:}
                    \BuiltInTok{System}\OperatorTok{.}\FunctionTok{out}\OperatorTok{.}\FunctionTok{println}\OperatorTok{(}\StringTok{"Inventory: "} \OperatorTok{+}\NormalTok{ player}\OperatorTok{.}\FunctionTok{getInventory}\OperatorTok{());}
                    \ControlFlowTok{break}\OperatorTok{;}
                \ControlFlowTok{case} \StringTok{"look"}\OperatorTok{:}
                    \CommentTok{// Room description already shown at top of loop}
                    \ControlFlowTok{break}\OperatorTok{;}
                \ControlFlowTok{case} \StringTok{"quit"}\OperatorTok{:}
                    \BuiltInTok{System}\OperatorTok{.}\FunctionTok{out}\OperatorTok{.}\FunctionTok{println}\OperatorTok{(}\StringTok{"Thanks for playing!"}\OperatorTok{);}
\NormalTok{                    scanner}\OperatorTok{.}\FunctionTok{close}\OperatorTok{();}
                    \ControlFlowTok{return}\OperatorTok{;}
                \KeywordTok{default}\OperatorTok{:}
                    \BuiltInTok{System}\OperatorTok{.}\FunctionTok{out}\OperatorTok{.}\FunctionTok{println}\OperatorTok{(}\StringTok{"Unknown command: "} \OperatorTok{+}\NormalTok{ command}\OperatorTok{);}
            \OperatorTok{\}}
            \BuiltInTok{System}\OperatorTok{.}\FunctionTok{out}\OperatorTok{.}\FunctionTok{println}\OperatorTok{();}
        \OperatorTok{\}}
    \OperatorTok{\}}
\OperatorTok{\}}
\end{Highlighting}
\end{Shaded}

\begin{center}\rule{0.5\linewidth}{0.5pt}\end{center}

\subsection{Lesson 10.2: Command-Line To-Do
Manager}\label{lesson-10.2-command-line-to-do-manager}

\begin{Shaded}
\begin{Highlighting}[]
\KeywordTok{import} \ImportTok{java}\OperatorTok{.}\ImportTok{util}\OperatorTok{.*;}
\KeywordTok{import} \ImportTok{java}\OperatorTok{.}\ImportTok{io}\OperatorTok{.*;}
\KeywordTok{import} \ImportTok{com}\OperatorTok{.}\ImportTok{google}\OperatorTok{.}\ImportTok{gson}\OperatorTok{.*;}   \CommentTok{// Add Gson library for JSON support}

\KeywordTok{public} \KeywordTok{class}\NormalTok{ TodoManager }\OperatorTok{\{}
    \KeywordTok{private} \BuiltInTok{List}\OperatorTok{\textless{}}\BuiltInTok{Map}\OperatorTok{\textless{}}\BuiltInTok{String}\OperatorTok{,} \BuiltInTok{Object}\OperatorTok{\textgreater{}\textgreater{}}\NormalTok{ todos }\OperatorTok{=} \KeywordTok{new} \BuiltInTok{ArrayList}\OperatorTok{\textless{}\textgreater{}();}
    \KeywordTok{private} \DataTypeTok{final} \BuiltInTok{String}\NormalTok{ FILE\_PATH }\OperatorTok{=} \StringTok{"todos.json"}\OperatorTok{;}

    \KeywordTok{public} \DataTypeTok{void} \FunctionTok{addTask}\OperatorTok{(}\BuiltInTok{String}\NormalTok{ description}\OperatorTok{)} \OperatorTok{\{}
        \BuiltInTok{Map}\OperatorTok{\textless{}}\BuiltInTok{String}\OperatorTok{,} \BuiltInTok{Object}\OperatorTok{\textgreater{}}\NormalTok{ task }\OperatorTok{=} \KeywordTok{new} \BuiltInTok{HashMap}\OperatorTok{\textless{}\textgreater{}();}
\NormalTok{        task}\OperatorTok{.}\FunctionTok{put}\OperatorTok{(}\StringTok{"description"}\OperatorTok{,}\NormalTok{ description}\OperatorTok{);}
\NormalTok{        task}\OperatorTok{.}\FunctionTok{put}\OperatorTok{(}\StringTok{"completed"}\OperatorTok{,} \KeywordTok{false}\OperatorTok{);}
\NormalTok{        todos}\OperatorTok{.}\FunctionTok{add}\OperatorTok{(}\NormalTok{task}\OperatorTok{);}
        \BuiltInTok{System}\OperatorTok{.}\FunctionTok{out}\OperatorTok{.}\FunctionTok{println}\OperatorTok{(}\StringTok{"Task added: "} \OperatorTok{+}\NormalTok{ description}\OperatorTok{);}
    \OperatorTok{\}}

    \KeywordTok{public} \DataTypeTok{void} \FunctionTok{viewTasks}\OperatorTok{()} \OperatorTok{\{}
        \ControlFlowTok{if} \OperatorTok{(}\NormalTok{todos}\OperatorTok{.}\FunctionTok{isEmpty}\OperatorTok{())} \OperatorTok{\{}
            \BuiltInTok{System}\OperatorTok{.}\FunctionTok{out}\OperatorTok{.}\FunctionTok{println}\OperatorTok{(}\StringTok{"No tasks in your list."}\OperatorTok{);}
            \ControlFlowTok{return}\OperatorTok{;}
        \OperatorTok{\}}
        \BuiltInTok{System}\OperatorTok{.}\FunctionTok{out}\OperatorTok{.}\FunctionTok{println}\OperatorTok{(}\StringTok{"{-}{-}{-} Your To{-}Do List {-}{-}{-}"}\OperatorTok{);}
        \ControlFlowTok{for} \OperatorTok{(}\DataTypeTok{int}\NormalTok{ i }\OperatorTok{=} \DecValTok{0}\OperatorTok{;}\NormalTok{ i }\OperatorTok{\textless{}}\NormalTok{ todos}\OperatorTok{.}\FunctionTok{size}\OperatorTok{();}\NormalTok{ i}\OperatorTok{++)} \OperatorTok{\{}
            \BuiltInTok{Map}\OperatorTok{\textless{}}\BuiltInTok{String}\OperatorTok{,} \BuiltInTok{Object}\OperatorTok{\textgreater{}}\NormalTok{ task }\OperatorTok{=}\NormalTok{ todos}\OperatorTok{.}\FunctionTok{get}\OperatorTok{(}\NormalTok{i}\OperatorTok{);}
            \BuiltInTok{String}\NormalTok{ status }\OperatorTok{=} \OperatorTok{(}\BuiltInTok{Boolean}\OperatorTok{)}\NormalTok{ task}\OperatorTok{.}\FunctionTok{get}\OperatorTok{(}\StringTok{"completed"}\OperatorTok{)} \OperatorTok{?} \StringTok{"[X]"} \OperatorTok{:} \StringTok{"[ ]"}\OperatorTok{;}
            \BuiltInTok{System}\OperatorTok{.}\FunctionTok{out}\OperatorTok{.}\FunctionTok{printf}\OperatorTok{(}\StringTok{"}\SpecialCharTok{\%d}\StringTok{. }\SpecialCharTok{\%s}\StringTok{ }\SpecialCharTok{\%s\%n}\StringTok{"}\OperatorTok{,}\NormalTok{ i }\OperatorTok{+} \DecValTok{1}\OperatorTok{,}\NormalTok{ status}\OperatorTok{,}\NormalTok{ task}\OperatorTok{.}\FunctionTok{get}\OperatorTok{(}\StringTok{"description"}\OperatorTok{));}
        \OperatorTok{\}}
        \BuiltInTok{System}\OperatorTok{.}\FunctionTok{out}\OperatorTok{.}\FunctionTok{println}\OperatorTok{(}\StringTok{"{-}{-}{-}{-}{-}{-}{-}{-}{-}{-}{-}{-}{-}{-}{-}{-}{-}{-}{-}{-}{-}{-}{-}"}\OperatorTok{);}
    \OperatorTok{\}}

    \KeywordTok{public} \DataTypeTok{void} \FunctionTok{markComplete}\OperatorTok{(}\DataTypeTok{int}\NormalTok{ taskNumber}\OperatorTok{)} \OperatorTok{\{}
        \DataTypeTok{int}\NormalTok{ index }\OperatorTok{=}\NormalTok{ taskNumber }\OperatorTok{{-}} \DecValTok{1}\OperatorTok{;}
        \ControlFlowTok{if} \OperatorTok{(}\NormalTok{index }\OperatorTok{\textgreater{}=} \DecValTok{0} \OperatorTok{\&\&}\NormalTok{ index }\OperatorTok{\textless{}}\NormalTok{ todos}\OperatorTok{.}\FunctionTok{size}\OperatorTok{())} \OperatorTok{\{}
\NormalTok{            todos}\OperatorTok{.}\FunctionTok{get}\OperatorTok{(}\NormalTok{index}\OperatorTok{).}\FunctionTok{put}\OperatorTok{(}\StringTok{"completed"}\OperatorTok{,} \KeywordTok{true}\OperatorTok{);}
            \BuiltInTok{System}\OperatorTok{.}\FunctionTok{out}\OperatorTok{.}\FunctionTok{println}\OperatorTok{(}\StringTok{"Task marked complete: "} \OperatorTok{+}\NormalTok{ todos}\OperatorTok{.}\FunctionTok{get}\OperatorTok{(}\NormalTok{index}\OperatorTok{).}\FunctionTok{get}\OperatorTok{(}\StringTok{"description"}\OperatorTok{));}
        \OperatorTok{\}} \ControlFlowTok{else} \OperatorTok{\{}
            \BuiltInTok{System}\OperatorTok{.}\FunctionTok{out}\OperatorTok{.}\FunctionTok{println}\OperatorTok{(}\StringTok{"Invalid task number."}\OperatorTok{);}
        \OperatorTok{\}}
    \OperatorTok{\}}

    \KeywordTok{public} \DataTypeTok{static} \DataTypeTok{void} \FunctionTok{main}\OperatorTok{(}\BuiltInTok{String}\OperatorTok{[]}\NormalTok{ args}\OperatorTok{)} \OperatorTok{\{}
\NormalTok{        TodoManager manager }\OperatorTok{=} \KeywordTok{new} \FunctionTok{TodoManager}\OperatorTok{();}
        \BuiltInTok{Scanner}\NormalTok{ scanner }\OperatorTok{=} \KeywordTok{new} \BuiltInTok{Scanner}\OperatorTok{(}\BuiltInTok{System}\OperatorTok{.}\FunctionTok{in}\OperatorTok{);}

        \BuiltInTok{System}\OperatorTok{.}\FunctionTok{out}\OperatorTok{.}\FunctionTok{println}\OperatorTok{(}\StringTok{"=== To{-}Do Manager ==="}\OperatorTok{);}

        \ControlFlowTok{while} \OperatorTok{(}\KeywordTok{true}\OperatorTok{)} \OperatorTok{\{}
            \BuiltInTok{System}\OperatorTok{.}\FunctionTok{out}\OperatorTok{.}\FunctionTok{println}\OperatorTok{(}\StringTok{"}\SpecialCharTok{\textbackslash{}n}\StringTok{1. Add Task  2. View Tasks  3. Mark Complete  4. Quit"}\OperatorTok{);}
            \BuiltInTok{System}\OperatorTok{.}\FunctionTok{out}\OperatorTok{.}\FunctionTok{print}\OperatorTok{(}\StringTok{"Choice: "}\OperatorTok{);}
            \BuiltInTok{String}\NormalTok{ choice }\OperatorTok{=}\NormalTok{ scanner}\OperatorTok{.}\FunctionTok{nextLine}\OperatorTok{().}\FunctionTok{trim}\OperatorTok{();}

            \ControlFlowTok{switch} \OperatorTok{(}\NormalTok{choice}\OperatorTok{)} \OperatorTok{\{}
                \ControlFlowTok{case} \StringTok{"1"}\OperatorTok{:}
                    \BuiltInTok{System}\OperatorTok{.}\FunctionTok{out}\OperatorTok{.}\FunctionTok{print}\OperatorTok{(}\StringTok{"Task description: "}\OperatorTok{);}
\NormalTok{                    manager}\OperatorTok{.}\FunctionTok{addTask}\OperatorTok{(}\NormalTok{scanner}\OperatorTok{.}\FunctionTok{nextLine}\OperatorTok{());}
                    \ControlFlowTok{break}\OperatorTok{;}
                \ControlFlowTok{case} \StringTok{"2"}\OperatorTok{:}
\NormalTok{                    manager}\OperatorTok{.}\FunctionTok{viewTasks}\OperatorTok{();}
                    \ControlFlowTok{break}\OperatorTok{;}
                \ControlFlowTok{case} \StringTok{"3"}\OperatorTok{:}
\NormalTok{                    manager}\OperatorTok{.}\FunctionTok{viewTasks}\OperatorTok{();}
                    \BuiltInTok{System}\OperatorTok{.}\FunctionTok{out}\OperatorTok{.}\FunctionTok{print}\OperatorTok{(}\StringTok{"Task number: "}\OperatorTok{);}
                    \ControlFlowTok{try} \OperatorTok{\{}
\NormalTok{                        manager}\OperatorTok{.}\FunctionTok{markComplete}\OperatorTok{(}\BuiltInTok{Integer}\OperatorTok{.}\FunctionTok{parseInt}\OperatorTok{(}\NormalTok{scanner}\OperatorTok{.}\FunctionTok{nextLine}\OperatorTok{()));}
                    \OperatorTok{\}} \ControlFlowTok{catch} \OperatorTok{(}\BuiltInTok{NumberFormatException}\NormalTok{ e}\OperatorTok{)} \OperatorTok{\{}
                        \BuiltInTok{System}\OperatorTok{.}\FunctionTok{out}\OperatorTok{.}\FunctionTok{println}\OperatorTok{(}\StringTok{"Invalid input."}\OperatorTok{);}
                    \OperatorTok{\}}
                    \ControlFlowTok{break}\OperatorTok{;}
                \ControlFlowTok{case} \StringTok{"4"}\OperatorTok{:}
                    \BuiltInTok{System}\OperatorTok{.}\FunctionTok{out}\OperatorTok{.}\FunctionTok{println}\OperatorTok{(}\StringTok{"Goodbye!"}\OperatorTok{);}
\NormalTok{                    scanner}\OperatorTok{.}\FunctionTok{close}\OperatorTok{();}
                    \ControlFlowTok{return}\OperatorTok{;}
                \KeywordTok{default}\OperatorTok{:}
                    \BuiltInTok{System}\OperatorTok{.}\FunctionTok{out}\OperatorTok{.}\FunctionTok{println}\OperatorTok{(}\StringTok{"Invalid choice."}\OperatorTok{);}
            \OperatorTok{\}}
        \OperatorTok{\}}
    \OperatorTok{\}}
\OperatorTok{\}}
\end{Highlighting}
\end{Shaded}

\begin{center}\rule{0.5\linewidth}{0.5pt}\end{center}

\section{Unit 11: Capstone Project --- Personal Data Manager in
Java}\label{unit-11-capstone-project-personal-data-manager-in-java}

\subsection{Overview}\label{overview-11}

Students build a fully functional command-line personal data manager in
Java.

\subsection{Project Options}\label{project-options}

\begin{itemize}
\tightlist
\item
  \textbf{Contact Manager:} name, phone, email, notes
\item
  \textbf{Book Tracker:} title, author, genre, read status
\item
  \textbf{Inventory System:} item name, quantity, location, notes
\item
  \textbf{To-Do App:} description, due date, priority, status
\end{itemize}

\subsection{Required Features}\label{required-features}

\begin{enumerate}
\def\labelenumi{\arabic{enumi}.}
\tightlist
\item
  \textbf{Add item} --- prompt for all attributes and save to list.
\item
  \textbf{View all items} --- display all stored items clearly.
\item
  \textbf{Edit item} --- select by number/name and modify.
\item
  \textbf{Delete item} --- remove selected item.
\item
  \textbf{Search/filter} --- find items by keyword.
\item
  \textbf{Persistence} --- save/load to a JSON or CSV file.
\end{enumerate}

\subsection{Architecture
Recommendations}\label{architecture-recommendations}

\begin{verbatim}
src/
|--- model/
|   +--- Contact.java        (data model class)
|--- storage/
|   +--- DataManager.java    (file I/O)
|--- ui/
|   +--- ConsoleUI.java      (user interface)
+--- Main.java               (entry point)
\end{verbatim}

\subsection{Assessment Criteria}\label{assessment-criteria}

\subsubsection{Functionality (60\%)}\label{functionality-60}

\begin{itemize}
\tightlist
\item[$\square$]
  Add, view, edit, and delete items work correctly.
\item[$\square$]
  Data persists between program runs.
\item[$\square$]
  Program handles invalid input without crashing.
\end{itemize}

\subsubsection{Code Quality (20\%)}\label{code-quality-20}

\begin{itemize}
\tightlist
\item[$\square$]
  Methods are small and have a single responsibility.
\item[$\square$]
  Classes follow OOP principles (encapsulation).
\item[$\square$]
  Meaningful names for all variables, methods, and classes.
\item[$\square$]
  Appropriate comments and Javadoc.
\end{itemize}

\subsubsection{Accessibility \& UX (15\%)}\label{accessibility-ux-15}

\begin{itemize}
\tightlist
\item[$\square$]
  Clear, descriptive prompts for all input.
\item[$\square$]
  Explicit confirmation messages after every action.
\item[$\square$]
  Consistent output format.
\item[$\square$]
  Informative error messages.
\end{itemize}

\subsubsection{Documentation (5\%)}\label{documentation-5}

\begin{itemize}
\tightlist
\item[$\square$]
  \texttt{README.md} with setup and usage instructions.
\item[$\square$]
  Reflection on challenges and solutions.
\end{itemize}

\end{document}